\chapter{Estrutura Curricular do Curso}

O currículo do curso de \nomecursonovo da UFU contempla as indicações e sugestões realizadas pela ACM (\textit{Association for Computing Machinery}), pelo IEEE (\textit{Institute of Electrical and Electronics Engineer}) nos currículos de referência criados em conjunto por ambas~\cite{ACM:CC2020}, pela SBC (Sociedade Brasileira de Computação), e fundamentalmente, pela Resolução CNE/CES nº 5 de 16 de novembro de 2016 do MEC que institui as Diretrizes Curriculares Nacionais específicas para os cursos da área de Computação~\cite{MEC:CNE:CES:5:2016}, em particular, para cursos de Engenharia da Computação.

O currículo do Curso está organizado em 8 (oito) períodos ou semestres, sendo que os componentes curriculares do curso estão divididos em: Disciplinas Obrigatórias, Disciplinas Optativas, Atividade de Conclusão de curso (Estágio Supervisionado ou Trabalho de Conclusão de Curso), Atividades Complementares e Atividades de Extensão. As disciplinas obrigatórias e optativas, por sua vez, possuem atividades classificadas nas modalidades: Prática, Teórica ou Remota, ou seja, algumas disciplinas obrigatórias e/ou optativas terão característica semipresencial ou remota.

A estrutura curricular apresenta um total de \textbf{\input{include/auto/par_carga_horaria_total_do_curso}horas} distribuídas em núcleos de formações: \begin{enumerate*}[label=\roman*.] \item Básica (Quadro~\ref{tab:nucleo_formacao_basica}), \item Humanística e de Extensão (Quadro~\ref{tab:nucleo_formacao_humanistica_extensao}), \item Tecnológica e Profissional (Quadro~\ref{tab:nucleo_formacao_tecnologica_profissional}) e \item Optativa e Complementar (Quadro~\ref{tab:nucleo_formacao_optativa_complementar}). \end{enumerate*}

A seguir, são apresentados cada um dos núcleos supracitados e suas respectivas composições. 

Note que o quadro do núcleo de formação optativa (Quadro~\ref{tab:nucleo_formacao_optativa_complementar}) elenca disciplinas que podem ser escolhidas pelo discente, sendo necessária a integralização de pelo menos \input{include/auto/par_carga_horaria_opt.tex}entre elas, além das \textbf{90 horas}
de ``Atividades Acadêmicas Complementares''. Já o quadro de formação humanística e de extensão (Quadro~\ref{tab:nucleo_formacao_humanistica_extensao}) elenca as disciplinas de ``Atividades Curriculares de Extensão'' que devem ser integralizadas a partir de participações em projetos de extensão universitária e da integralização da disciplina de ``Memorial'', totalizando \textbf{345 horas}
\unskip.


% Note que o quadro do núcleo de formação optativa (Quadro~\ref{tab:nucleo_formacao_optativa_complementar}) elenca disciplinas que podem ser escolhidas pelo discente, sendo necessária a integralização de pelo menos \textbf{90 horas} entre elas, além das \textbf{90 horas} de ``Atividades Acadêmicas Complementares''.


% Note que o quadro do núcleo de formação optativa (Quadro~\ref{tab:nucleo_formacao_optativa_complementar}) elenca disciplinas que podem ser escolhidas pelo discente, sendo necessária a integralização de pelo menos \textbf{495 horas} entre elas, além das 75 horas de ``Atividades Acadêmicas Complementares''.

Ao final da seção, apresentam-se a síntese de distribuição de carga horária por núcleo de formação (Quadro~\ref{tab:carga_horaria_por_nucleo_formacao}) e a síntese de distribuição de carga horária por componente curricular (Quadro~\ref{tab:carga_horaria_por_componente_curricular}).

% \newpage
%\section{Estrutura Curricular do Curso}

% \todo{Completar tabelas}

% \renewcommand{\familydefault}{\sfdefault} \normalfont
% \begin{longtblr}[
%     theme = ecp,
%     caption = {Núcleo de Formação Básica},
%     label = {tab:formacao_basica},
% ]{
%   colspec = {Q[l,m,wd=65mm]Q[c,m,wd=24mm]Q[c,m,wd=8mm]Q[c,m,wd=8mm]Q[c,m,wd=8mm]Q[c,m,wd=10mm]},
%   rowhead = 1,
%   row{odd} = {bg=CinzaClaro},
%   row{1} = {bg=AzulEscuro, fg=white},
%   row{19} = {bg=AzulClaro, fg=white},
%   cells  = {font=\fontsize{10pt}{12pt}\selectfont},
% } 
%     % Componente & Código & CH Teór. & CH Prát.  & CH Rem. & CH Total \\
%     \textbf{Componente} & \textbf{Código}  & \textbf{CH Teór.} & \textbf{CH Prát.}  & \textbf{CH Rem.} & \textbf{CH Total} \\
%     Cálculo Diferencial e Integral I & FAMAT31011 & 90 & 0 & 0 & 90 \\
%     Geometria Analítica & FAMAT31021 & 60 & 0 & 0 & 60 \\
%     Cálculo Diferencial e Integral II & FAMAT31012 & 90 & 0 & 0 & 90 \\
%     Álgebra Linear & FAMAT31022 & 45 & 0 & 0 & 45 \\
%     Física Básica: Mecânica & INFIS39031 & 60 & 0 & 0 & 60 \\
%     Estrutura e Interpretação de Programas de Computadores I & FEELT? & 30 & 0 & 0 & 30 \\
%     Estrutura e Interpretação de Programas de Computadores II & FEELT? & 30 & 0 & 0 & 30 \\
%     Lógica e Matemática Discreta & FEELT31203 & 45 & 0 & 0 & 45 \\
%     Metrologia & FEELT31204 & 30 & 30 & 0 & 60 \\
%     Cálculo Diferencial e Integral III & FAMAT31013 & 90 & 0 & 0 & 90 \\
%     Física Básica: Eletricidade e Magnetismo & INFIS39033 & 60 & 0 & 0 & 60 \\
%     Estatística & FAMAT31033 & 60 & 0 & 0 & 60 \\
%     Métodos Matemáticos & FAMAT31031 & 75 & 0 & 0 & 75 \\
%     Cálculo Numérico & FAMAT31032 & 60 & 0 & 0 & 60 \\
%     Gerenciamento de Projetos & FAGEN? & 30 & 0 & 0 & 30 \\
%     Engenharia Econômica & IERI? & 30 & 0 & 0 & 30 \\
%     Direito e Legislação & FADIR49111 & 30 & 0 & 0 & 30 \\
%     \textbf{TOTAL} &  & \textbf{915} & \textbf{30} & \textbf{0} & \textbf{945} \\
% \end{longtblr}
% \renewcommand{\familydefault}{\rmdefault} \normalfont

\begin{longtblr}[
        theme = ecp,
        caption = {Núcleo de Formação Básica},
        label = {tab:nucleo_formacao_basica},]{
    colspec = {Q[l,m,wd=63mm]Q[l,m,wd=23mm]Q[c,m,wd=7mm]Q[c,m,wd=7mm]Q[c,m,wd=7mm]Q[c,m,wd=7mm]Q[c,m,wd=9mm]},
    rowhead = 1,
    row{odd} = {bg=CinzaClaro},
    row{1} = {bg=AzulEscuro, fg=white},
    row{16} = {bg=AzulClaro, fg=white},
    cells  = {font=\fontsize{10pt}{12pt}\selectfont},
    }
        \textbf{Componente} & \textbf{Código}  & \textbf{CH Teór.} & \textbf{CH Prát.}  & \textbf{CH EaD} & \textbf{CH Ext.} & \textbf{CH Total} \\
        Álgebra Linear & FAMAT31022 & 45 & 0 & 0 & 0 & 45 \\
    Cálculo Diferencial e Integral I & FAMAT31011 & 90 & 0 & 0 & 0 & 90 \\
    Cálculo Diferencial e Integral II & FAMAT31012 & 90 & 0 & 0 & 0 & 90 \\
    Cálculo Diferencial e Integral III & FAMAT31013 & 90 & 0 & 0 & 0 & 90 \\
    Cálculo Numérico & FAMAT31032 & 60 & 0 & 0 & 0 & 60 \\
    Estatística & FAMAT31033 & 60 & 0 & 0 & 0 & 60 \\
    Física Básica: Eletricidade e Magnetismo & INFIS39033 & 60 & 0 & 0 & 0 & 60 \\
    Experimental de Física Básica: Eletricidade e Magnetismo & INFIS39034 & 0 & 30 & 0 & 0 & 30 \\
    Física Básica: Mecânica & INFIS39031 & 60 & 0 & 0 & 0 & 60 \\
    Experimental de Física Básica: Mecânica & INFIS39032 & 0 & 30 & 0 & 0 & 30 \\
    Geometria Analítica & FAMAT31021 & 60 & 0 & 0 & 0 & 60 \\
    Métodos Matemáticos & FAMAT31031 & 75 & 0 & 0 & 0 & 75 \\
    Lógica e Matemática Discreta para Computação & FEELT!LMD & 45 & 0 & 0 & 0 & 45 \\
    Metrologia & FEELT31204 & 30 & 30 & 0 & 0 & 60 \\*
    \textbf{TOTAL} &  & \textbf{765} & \textbf{90} & \textbf{0} & \textbf{0} & \textbf{855} \\
    \end{longtblr}
    

\begin{longtblr}[
        theme = ecp,
        caption = {Núcleo de Formação Tecnológica e Profissional},
        label = {tab:nucleo_formacao_tecnologica_profissional},
        note{$\ast$} = {\scriptsize Contando com as 300 horas da Atividade de Conclusão de Curso.}]{
    colspec = {Q[l,m,wd=63mm]Q[l,m,wd=23mm]Q[c,m,wd=7mm]Q[c,m,wd=7mm]Q[c,m,wd=7mm]Q[c,m,wd=7mm]Q[c,m,wd=9mm]},
    rowhead = 1,
    row{odd} = {bg=CinzaClaro},
    row{1} = {bg=AzulEscuro, fg=white},
    row{38} = {bg=AzulClaro, fg=white},
    cells  = {font=\fontsize{10pt}{12pt}\selectfont},
    }
        \textbf{Componente} & \textbf{Código}  & \textbf{CH Teór.} & \textbf{CH Prát.}  & \textbf{CH EaD} & \textbf{CH Ext.} & \textbf{CH Total} \\
        Atividade de Conclusão de Curso & ACC & -- & -- & -- & -- & 300\TblrNote{$\ast$} \\
    Empreendedorismo e Inovação & FAGEN39401 & 30 & 0 & 0 & 0 & 30 \\
    Engenharia Econômica & IERI39601 & 30 & 0 & 0 & 0 & 30 \\
    Análise de Algoritmos & FEELT!AA & 60 & 0 & 0 & 0 & 60 \\
    Aprendizagem Profunda e Modelos Generativos & FEELT!APMG & 30 & 15 & 0 & 0 & 45 \\
    Aprendizagem de Máquina & FEELT!AMAQ & 30 & 15 & 0 & 0 & 45 \\
    Arquitetura de Software Aplicada & FEELT!ASA & 15 & 45 & 0 & 0 & 60 \\
    Arquitetura e Organização de Computadores & FEELT!AOC & 60 & 0 & 0 & 0 & 60 \\
    Banco de Dados & FEELT!BD & 15 & 45 & 0 & 0 & 60 \\
    Circuitos Elétricos I & FEELT31301 & 75 & 0 & 0 & 0 & 75 \\
    Experimental de Circuitos Elétricos I & FEELT31302 & 0 & 15 & 0 & 0 & 15 \\
    Elementos de Sistemas Computacionais & FEELT!ESC & 30 & 0 & 30 & 0 & 60 \\
    Eletrônica Analógica I & FEELT31401 & 60 & 0 & 0 & 0 & 60 \\
    Experimental de Eletrônica Analógica I & FEELT31402 & 0 & 30 & 0 & 0 & 30 \\
    Engenharia de Contexto em IA Generativa & FEELT!ECIA & 15 & 0 & 30 & 0 & 45 \\
    Engenharia de Software & FEELT!ESOF & 45 & 0 & 0 & 0 & 45 \\
    Estrutura de Dados & FEELT!ED & 60 & 0 & 0 & 0 & 60 \\
    Fundamentos da Inteligência Artificial & FEELT!FIA & 30 & 15 & 0 & 0 & 45 \\
    IA Aplicada à Cibersegurança & FEELT!IACS & 15 & 0 & 30 & 0 & 45 \\
    Introdução à Prototipagem Eletrônica & FEELT!IPE & 0 & 30 & 0 & 0 & 30 \\
    Processamento Digital de Sinais & FEELT31628 & 45 & 15 & 0 & 0 & 60 \\
    Programação Funcional & FEELT!PF & 30 & 15 & 0 & 0 & 45 \\
    Programação Orientada a Objetos & FEELT!POO & 30 & 30 & 0 & 0 & 60 \\
    Programação Procedimental & FEELT!PP & 30 & 30 & 0 & 0 & 60 \\
    Programação Script & FEELT!PS & 30 & 30 & 0 & 0 & 60 \\
    Redes de Comunicações I & FEELT31724 & 45 & 15 & 0 & 0 & 60 \\
    Redes de Comunicações II & FEELT31831 & 45 & 15 & 0 & 0 & 60 \\
    Segurança de Sistemas Computacionais & FEELT!SEG & 15 & 0 & 30 & 0 & 45 \\
    Sistemas Computacionais em Tempo Real & FEELT!SCTR & 15 & 30 & 0 & 0 & 45 \\
    Sistemas Digitais & FEELT31409 & 30 & 0 & 0 & 0 & 30 \\
    Experimental de Sistemas Digitais & FEELT31410 & 0 & 30 & 0 & 0 & 30 \\
    Sistemas Distribuídos & FEELT!SD & 45 & 0 & 0 & 0 & 45 \\
    Sistemas Embarcados I & FEELT31523 & 45 & 30 & 0 & 0 & 75 \\
    Sistemas Operacionais & FEELT!SO & 45 & 0 & 0 & 0 & 45 \\
    Sistemas e Controle & FEELT!SCON & 30 & 15 & 0 & 0 & 45 \\
    Teoria da Computação & FEELT!TCOMP & 45 & 0 & 0 & 0 & 45 \\*
    \textbf{TOTAL} &  & \textbf{1125} & \textbf{465} & \textbf{120} & \textbf{0} & \textbf{2010$^\ast$} \\
    \end{longtblr}
    

\begin{longtblr}[
        theme = ecp,
        caption = {Núcleo de Formação Humanística e de Extensão},
        label = {tab:nucleo_formacao_humanistica_extensao},]{
    colspec = {Q[l,m,wd=63mm]Q[l,m,wd=23mm]Q[c,m,wd=7mm]Q[c,m,wd=7mm]Q[c,m,wd=7mm]Q[c,m,wd=7mm]Q[c,m,wd=9mm]},
    rowhead = 1,
    row{odd} = {bg=CinzaClaro},
    row{1} = {bg=AzulEscuro, fg=white},
    row{9} = {bg=AzulClaro, fg=white},
    cells  = {font=\fontsize{10pt}{12pt}\selectfont},
    }
        \textbf{Componente} & \textbf{Código}  & \textbf{CH Teór.} & \textbf{CH Prát.}  & \textbf{CH EaD} & \textbf{CH Ext.} & \textbf{CH Total} \\
        Atividades Curriculares de Extensão I & ACE1 & 0 & 0 & 0 & 90 & 90 \\
    Atividades Curriculares de Extensão II & ACE2 & 0 & 0 & 0 & 90 & 90 \\
    Atividades Curriculares de Extensão III & ACE3 & 0 & 0 & 0 & 75 & 75 \\
    Atividades Curriculares de Extensão IV & ACE4 & 0 & 0 & 0 & 60 & 60 \\
    Inteligência Artificial, Ética e Sociedade & FEELT!IAES & 30 & 0 & 0 & 0 & 30 \\
    Introdução à Engenharia de Computação & FEELT31103 & 30 & 0 & 0 & 0 & 30 \\
    Memorial de Atividades Curriculares de Extensão & FEELT!MACE & 0 & 0 & 0 & 30 & 30 \\*
    \textbf{TOTAL} &  & \textbf{60} & \textbf{0} & \textbf{0} & \textbf{345} & \textbf{405} \\
    \end{longtblr}
    

\begin{longtblr}[
        theme = ecp,
        caption = {Núcleo de Formação Optativa e Complementar},
        label = {tab:nucleo_formacao_optativa_complementar},
        remark{\scriptsize \textbf{Nota}} = {\scriptsize As disciplinas optativas pré-aprovadas são apresentadas na \autoref{sec:optativas}.}]{
    colspec = {Q[l,m,wd=63mm]Q[l,m,wd=23mm]Q[c,m,wd=7mm]Q[c,m,wd=7mm]Q[c,m,wd=7mm]Q[c,m,wd=7mm]Q[c,m,wd=9mm]},
    rowhead = 1,
    row{odd} = {bg=CinzaClaro},
    row{1} = {bg=AzulEscuro, fg=white},
    row{4} = {bg=AzulClaro, fg=white},
    cells  = {font=\fontsize{10pt}{12pt}\selectfont},
    }
        \textbf{Componente} & \textbf{Código}  & \textbf{CH Teór.} & \textbf{CH Prát.}  & \textbf{CH EaD} & \textbf{CH Ext.} & \textbf{CH Total} \\
        Atividades Acadêmicas Complementares & AAC & -- & -- & -- & -- & 90 \\
    Disciplinas Optativas & OPT & -- & -- & -- & -- & 90 \\*
    \textbf{TOTAL} &  & \textbf{--} & \textbf{--} & \textbf{--} & \textbf{--} & \textbf{180} \\
    \end{longtblr}
    

\begin{longtblr}[
        theme = ecp,
        caption = {Carga Horária por Núcleo de Formação},
        label = {tab:carga_horaria_por_nucleo_formacao}
    ]{
    colspec = {Q[l,m,wd=100mm]Q[c,m,wd=20mm]Q[c,m,wd=20mm]},
    rowhead = 1,
    row{odd} = {bg=CinzaClaro},
    row{1} = {bg=AzulEscuro, fg=white},
    row{6} = {bg=AzulClaro, fg=white},
    cells  = {font = \fontsize{10pt}{12pt}\selectfont}
    }
        \textbf{Componentes Curriculares} & \textbf{CH Total} & \textbf{Percentual} \\
        Núcleo de Formação Básica & 855 & 24,8\% \\
    Núcleo de Formação Tecnológica e Profissional & 2010 & 58,3\% \\
    Núcleo de Formação Humanística e de Extensão & 405 & 11,7\% \\
    Núcleo de Formação Optativa e Complementar & 180 & 5,2\% \\*
   \textbf{TOTAL} &  \textbf{3450} & \textbf{100,0\%} \\
    \end{longtblr}
    

\begin{longtblr}[
        theme = ecp,
        caption = {Carga Horária por Componente Curricular},   remark{\scriptsize \textbf{Sobre Estágio e Atividades Complementares}} = {\scriptsize A somatória de 390 horas dos componentes Atividade de Conclusão de Curso (possível Estágio Supervisionado, ver Seção~\ref{sec:acc_conclusao}) e Atividades Acadêmicas Complementares totalizam \textbf{11,3\%} da carga horária total do curso, respeitando o limite legal de 20\% para estágios e atividades complementares estabelecido pela legislação vigente~\cite{MEC:CNE:CES:2:2007}.}, 
        label = {tab:carga_horaria_por_componente_curricular}
    ]{
    colspec = {Q[l,m,wd=100mm]Q[c,m,wd=20mm]Q[c,m,wd=20mm]},
    rowhead = 1,
    row{odd} = {bg=CinzaClaro},
    row{1} = {bg=AzulEscuro, fg=white},
    row{7} = {bg=AzulClaro, fg=white},
    cells  = {font = \fontsize{10pt}{12pt}\selectfont}
    }
        \textbf{Componentes Curriculares} & \textbf{CH Total} & \textbf{Percentual} \\
        Disciplinas Obrigatórias & 2625 & 76,1\% \\
    Disciplinas Optativas & 90 & 2,6\% \\
    Atividades Curriculares de Extensão & 345 & 10,0\% \\
    Atividade de Conclusão de Curso & 300 & 8,7\% \\
    Atividades Acadêmicas Complementares & 90 & 2,6\% \\*
   \textbf{TOTAL} &  \textbf{3450} & \textbf{100,0\%} \\
    \end{longtblr}
    

\begin{longtblr}[
        theme = ecp,
        caption = {EaD e Extensão por Componente Curricular},
        label = {tab:ead_extensao_por_componente_curricular},
        remark{\scriptsize \textbf{Sobre EaD}} = {\scriptsize A carga horária em formato EaD mínima (CHD\textsubscript{min}) e máxima (CHD\textsubscript{max}) foi calculada considerando-se um percentual flexível de até 20,0\% sobre a carga horária presencial (CHT+CHP) das disciplinas obrigatórias, conforme Seção~\ref{sec:planejamento_ead}. Com isso, comprova-se que a flexibilidade de \textbf{3,5\%} a \textbf{20,6\%} de sua carga horária em formato EaD mantém-se, no melhor e no pior caso, em conformidade com o limite legal de 30\% estabelecido pela legislação vigente~\cite{Decreto:12456:2025}.},
        remark{\scriptsize \textbf{Sobre Extensão}} = {\scriptsize O curso possui \textbf{10,0\%} de sua carga horária total dedicada à Extensão (CHE), em conformidade com as diretrizes e normativas vigentes~\cite{MEC:CNE:CES:7:2018,UFU:CONSUN:25:2019}.}
    ]{
    colspec = {Q[l,m,wd=72mm]Q[c,m,wd=15mm]Q[c,m,wd=15mm]Q[c,m,wd=15mm]Q[c,m,wd=15mm]},
    rowhead = 1,
    row{odd} = {bg=CinzaClaro},
    row{1} = {bg=AzulEscuro, fg=white},
    row{7} = {bg=AzulClaro, fg=white},
    row{8} = {bg=AzulClaro, fg=white},
    cells  = {font = \fontsize{10pt}{12pt}\selectfont},
    }
        \textbf{Componentes Curriculares} & \textbf{CHD$_{\min}$} & \textbf{CHD$_{\max}$} & \textbf{CHE} & \textbf{CH Total} \\
        Disciplinas Obrigatórias  & 120  & 621  & 0 & 2625 \\
    Disciplinas Optativas  & 0  & 90  & 0 & 90 \\
    Atividades Curriculares de Extensão  & 0  & 0  & 345 & 345 \\
    Atividade de Conclusão de Curso  & 0  & 0  & 0 & 300 \\
    Atividades Acadêmicas Complementares  & 0  & 0  & 0 & 90 \\*
   \textbf{TOTAL} &  \textbf{120} &  \textbf{711} &  \textbf{345} &  \textbf{3450} \\*
        \textbf{Percentual} & \textbf{3,5\%} & \textbf{20,6\%} & \textbf{10,0\%} & \textbf{100,0\%} \\
    \end{longtblr}
    

\section{Conhecimentos e Temas Transversais}

Com base nas disposições da Resolução CNE/CP nº 1/2004 (Educação das Relações Étnico-Raciais e para o Ensino de História e Cultura Afro-Brasileira e Africana), do Decreto nº 5.626/2005 (Ensino de Libras e Formação de Professores), da Resolução CNE nº 1/2012 (Educação em Direitos Humanos), da Resolução CNE nº 2/2012 (Educação Ambiental) e da Lei nº 13.425/2017 (Educação para Prevenção e Redução de Desastres, especialmente aplicável às Engenharias), apresentam-se, a seguir, os quadros que identificam os componentes curriculares do curso que contemplam esses conteúdos obrigatórios de forma transversal e integrada ao longo da formação.

\begin{longtblr}[
        theme = ecp,
        caption = {Distribuição dos Conteúdos Curriculares Referentes a Educação em Relações Étnico-Raciais},
        label = {tab:disciplinas_ref_etnico_raciais},
        remark{\small \textbf{Referência}} = {\small Resolução CNE/CP nº 1/2004 e Parecer CNE/CP nº 3/2004.}]{
    colspec = {Q[c,m,wd=20mm]Q[l,m,wd=90mm]Q[l,m,wd=25mm]},
    rowhead = 1,
    row{odd} = {bg=CinzaClaro},
    row{1} = {bg=AzulEscuro, fg=white},
    row{5} = {bg=AzulClaro, fg=white},
    cells  = {font=\fontsize{10pt}{12pt}\selectfont},
    }
        \textbf{Período} & \textbf{Componente} & \textbf{Código} \\
        1\textordmasculine & Introdução à Engenharia de Computação & FEELT31103 \\
    3\textordmasculine & Inteligência Artificial, Ética e Sociedade & FEELT!IAES \\
    Optativa & Direito e Legislação & FADIR39801 \\*

    \end{longtblr}
    

\begin{longtblr}[
        theme = ecp,
        caption = {Distribuição dos Conteúdos Curriculares Referentes a Ensino de Libras},
        label = {tab:disciplinas_ref_libras},
        remark{\small \textbf{Referência}} = {\small Decreto nº 5.626/2005.}]{
    colspec = {Q[c,m,wd=20mm]Q[l,m,wd=90mm]Q[l,m,wd=25mm]},
    rowhead = 1,
    row{odd} = {bg=CinzaClaro},
    row{1} = {bg=AzulEscuro, fg=white},
    row{4} = {bg=AzulClaro, fg=white},
    cells  = {font=\fontsize{10pt}{12pt}\selectfont},
    }
        \textbf{Período} & \textbf{Componente} & \textbf{Código} \\
        Optativa & Língua Brasileira de Sinais - Libras I & LIBRAS01 \\
    Optativa & Língua Brasileira de Sinais - Libras II & LIBRAS02 \\*

    \end{longtblr}
    

\begin{longtblr}[
        theme = ecp,
        caption = {Distribuição dos Conteúdos Curriculares Referentes a Educação em Direitos Humanos},
        label = {tab:disciplinas_ref_humano},
        remark{\small \textbf{Referência}} = {\small Resolução CNE nº 1/2012.}]{
    colspec = {Q[c,m,wd=20mm]Q[l,m,wd=90mm]Q[l,m,wd=25mm]},
    rowhead = 1,
    row{odd} = {bg=CinzaClaro},
    row{1} = {bg=AzulEscuro, fg=white},
    row{5} = {bg=AzulClaro, fg=white},
    cells  = {font=\fontsize{10pt}{12pt}\selectfont},
    }
        \textbf{Período} & \textbf{Componente} & \textbf{Código} \\
        1\textordmasculine & Introdução à Engenharia de Computação & FEELT31103 \\
    3\textordmasculine & Inteligência Artificial, Ética e Sociedade & FEELT!IAES \\
    Optativa & Direito e Legislação & FADIR39801 \\*

    \end{longtblr}
    

\begin{longtblr}[
        theme = ecp,
        caption = {Distribuição dos Conteúdos Curriculares Referentes a Educação Ambiental},
        label = {tab:disciplinas_ref_ambiental},
        remark{\small \textbf{Referência}} = {\small Resolução CNE nº 2/2012.}]{
    colspec = {Q[c,m,wd=20mm]Q[l,m,wd=90mm]Q[l,m,wd=25mm]},
    rowhead = 1,
    row{odd} = {bg=CinzaClaro},
    row{1} = {bg=AzulEscuro, fg=white},
    row{7} = {bg=AzulClaro, fg=white},
    cells  = {font=\fontsize{10pt}{12pt}\selectfont},
    }
        \textbf{Período} & \textbf{Componente} & \textbf{Código} \\
        1\textordmasculine & Introdução à Engenharia de Computação & FEELT31103 \\
    2\textordmasculine & Metrologia & FEELT31204 \\
    3\textordmasculine & Inteligência Artificial, Ética e Sociedade & FEELT!IAES \\
    4\textordmasculine & Experimental de Circuitos Elétricos I & FEELT31302 \\
    Optativa & Direito e Legislação & FADIR39801 \\*

    \end{longtblr}
    

\begin{longtblr}[
        theme = ecp,
        caption = {Distribuição dos Conteúdos Curriculares Referentes a Educação em Prevenção a Desastres para Engenharias},
        label = {tab:disciplinas_ref_desastre},
        remark{\small \textbf{Referência}} = {\small Lei nº 13.425/2017.}]{
    colspec = {Q[c,m,wd=20mm]Q[l,m,wd=90mm]Q[l,m,wd=25mm]},
    rowhead = 1,
    row{odd} = {bg=CinzaClaro},
    row{1} = {bg=AzulEscuro, fg=white},
    row{5} = {bg=AzulClaro, fg=white},
    cells  = {font=\fontsize{10pt}{12pt}\selectfont},
    }
        \textbf{Período} & \textbf{Componente} & \textbf{Código} \\
        1\textordmasculine & Introdução à Engenharia de Computação & FEELT31103 \\
    2\textordmasculine & Metrologia & FEELT31204 \\
    4\textordmasculine & Experimental de Circuitos Elétricos I & FEELT31302 \\*

    \end{longtblr}
    


\section{Conhecimento em Computação}

Esta seção descreve a distribuição dos conhecimentos estruturantes da Computação ao longo do curso, em conformidade com a Resolução CNE/CES nº 5/2016 (DCN Computação) e com as recomendações da ACM/IEEE (CC2020). São apresentados, a seguir, os quadros que relacionam os componentes curriculares aos eixos de competência da área, incluindo fundamentos teóricos, sistemas, desenvolvimento de software, dados e inteligência artificial.

\subsection{Refente à Inteligência Artificial}
\label{sec:sobre_IA}

Reconhecendo a \textit{Inteligência Artificial} como tecnologia transformadora que redefine as demandas sobre a infraestrutura computacional, o currículo foi estrategicamente expandido com um robusto núcleo de formação nessa área de competência. Esta alocação reflete a perspectiva singular da Engenharia de Computação, que trata \textit{Sistemas Inteligentes (IA)} não apenas como aplicações de software, mas como cargas de trabalho de alta complexidade cuja viabilidade depende fundamentalmente do projeto, da otimização e da gestão da arquitetura de hardware especializado (GPUs, TPUs), de plataformas de treinamento distribuído e de infraestruturas de produção escaláveis e seguras. O quadro a seguir detalha a distribuição dos conteúdos curriculares referentes à \textit{Inteligência Artificial} no curso.

\begin{longtblr}[
        theme = ecp,
        caption = {Distribuição dos Conteúdos Curriculares Referentes à Área de Inteligência Artificial},
        label = {tab:disciplinas_ref_ia},
        remark{\small \textbf{Referência}} = {\small Disciplinas obrigatórias, existem opções entre as optativas.}]{
    colspec = {Q[c,m,wd=20mm]Q[l,m,wd=90mm]Q[l,m,wd=25mm]},
    rowhead = 1,
    row{odd} = {bg=CinzaClaro},
    row{1} = {bg=AzulEscuro, fg=white},
    row{8} = {bg=AzulClaro, fg=white},
    cells  = {font=\fontsize{10pt}{12pt}\selectfont},
    }
        \textbf{Período} & \textbf{Componente} & \textbf{Código} \\
        3\textordmasculine & Inteligência Artificial, Ética e Sociedade & FEELT!IAES \\
    5\textordmasculine & Fundamentos da Inteligência Artificial & FEELT!FIA \\
    6\textordmasculine & Aprendizagem de Máquina & FEELT!AMAQ \\
    7\textordmasculine & Aprendizagem Profunda e Modelos Generativos & FEELT!APMG \\
    8\textordmasculine & Engenharia de Contexto em IA Generativa & FEELT!ECIA \\
    8\textordmasculine & IA Aplicada à Cibersegurança & FEELT!IACS \\*

    \end{longtblr}
    

\subsection{Refente à DCN Computação}

Com base na Resolução CNE/CES nº 5/2016, fundamentada no Parecer CNE/ CES nº 136/2012, que estabelece os conteúdos curriculares para os cursos da área de Computação (DCN Computação), apresentam-se a seguir os componentes curriculares do curso que contemplam esses conteúdos de forma transversal e integrada.

No referido parecer, estão definidos os \textbf{conteúdos tecnológicos e básicos comuns aos cursos de Bacharelado e Licenciatura} da área, abrangendo fundamentos teóricos da computação, desenvolvimento de software, infraestrutura e sistemas computacionais, além de aspectos éticos, sociais e econômicos relacionados à profissão.

O documento também especifica os \textbf{conteúdos próprios da Engenharia de Computação}, voltados ao domínio de hardware, sistemas embarcados, eletrônica, controle e automação, consolidando a interface entre a Computação e as Engenharias.

A seleção dos conteúdos apresentados nos quadros seguintes reflete o entendimento do NDE e do Colegiado do Curso quanto à aderência entre os componentes curriculares e os eixos formativos estabelecidos pelas Diretrizes Nacionais, garantindo a formação abrangente e atualizada prevista na norma.

\begin{longtblr}[
        theme = ecp,
        caption = {Distribuição dos Conteúdos Curriculares Referentes à Formação Tecnológica e Básica para todos os Cursos de Bacharelado e de Licenciatura (DCN Computação)},
        label = {tab:disciplinas_ref_dcn_base},
        remark{\small \textbf{Referência}} = {\small Parecer CNE/CES Nº 136/2012, item 3.1.}]{
    colspec = {Q[c,m,wd=20mm]Q[l,m,wd=50mm]Q[l,m,wd=25mm]Q[l,m,wd=40mm]},
    rowhead = 1,
    row{odd} = {bg=CinzaClaro},
    row{1} = {bg=AzulEscuro, fg=white},
    row{60} = {bg=AzulClaro, fg=white},
    cells  = {font=\fontsize{10pt}{12pt}\selectfont},
    }
        \textbf{Período} & \textbf{Componente} & \textbf{Código} & \textbf{Conteúdos}\\
        1\textordmasculine & Cálculo Diferencial e Integral I & FAMAT31011 & {\scriptsize matemática do contínuo} \\
    1\textordmasculine & Geometria Analítica & FAMAT31021 & {\scriptsize matemática do contínuo} \\
    1\textordmasculine & Introdução à Engenharia de Computação & FEELT31103 & {\scriptsize computação e sociedade; empreendedorismo; filosofia; meio ambiente; metodologia cientifica; ética e legislação} \\
    1\textordmasculine & Introdução à Prototipagem Eletrônica & FEELT!IPE & {\scriptsize automação; sistemas embarcados} \\
    1\textordmasculine & Lógica e Matemática Discreta para Computação & FEELT!LMD & {\scriptsize análise combinatória; lógica; matemática discreta; métodos formais} \\
    1\textordmasculine & Programação Script & FEELT!PS & {\scriptsize programação} \\
    2\textordmasculine & Álgebra Linear & FAMAT31022 & {\scriptsize estruturas algébricas; matemática do contínuo} \\
    2\textordmasculine & Cálculo Diferencial e Integral II & FAMAT31012 & {\scriptsize matemática do contínuo} \\
    2\textordmasculine & Física Básica: Mecânica & INFIS39031 & {\scriptsize matemática do contínuo} \\
    2\textordmasculine & Experimental de Física Básica: Mecânica & INFIS39032 & {\scriptsize matemática do contínuo} \\
    2\textordmasculine & Metrologia & FEELT31204 & {\scriptsize automação; matemática do contínuo; meio ambiente} \\
    2\textordmasculine & Programação Procedimental & FEELT!PP & {\scriptsize programação} \\
    2\textordmasculine & Sistemas Digitais & FEELT31409 & {\scriptsize circuitos digitais} \\
    2\textordmasculine & Experimental de Sistemas Digitais & FEELT31410 & {\scriptsize circuitos digitais} \\
    3\textordmasculine & Cálculo Diferencial e Integral III & FAMAT31013 & {\scriptsize matemática do contínuo} \\
    3\textordmasculine & Elementos de Sistemas Computacionais & FEELT!ESC & {\scriptsize arquitetura e organização de computadores; circuitos digitais; compiladores} \\
    3\textordmasculine & Estrutura de Dados & FEELT!ED & {\scriptsize abstração e estruturas de dados; programação} \\
    3\textordmasculine & Física Básica: Eletricidade e Magnetismo & INFIS39033 & {\scriptsize matemática do contínuo} \\
    3\textordmasculine & Experimental de Física Básica: Eletricidade e Magnetismo & INFIS39034 & {\scriptsize matemática do contínuo} \\
    3\textordmasculine & Inteligência Artificial, Ética e Sociedade & FEELT!IAES & {\scriptsize computação e sociedade; filosofia; inteligência artificial e computacional; meio ambiente; ética e legislação} \\
    3\textordmasculine & Programação Funcional & FEELT!PF & {\scriptsize estruturas algébricas; fundamentos de linguagens; novos paradigmas de computação; programação} \\
    3\textordmasculine & Programação Orientada a Objetos & FEELT!POO & {\scriptsize engenharia de software; programação} \\
    4\textordmasculine & Análise de Algoritmos & FEELT!AA & {\scriptsize abstração e estruturas de dados; algoritmos e complexidade; teoria dos grafos} \\
    4\textordmasculine & Arquitetura e Organização de Computadores & FEELT!AOC & {\scriptsize arquitetura e organização de computadores} \\
    4\textordmasculine & Atividades Curriculares de Extensão I & ACE1 & {\scriptsize computação e sociedade} \\
    4\textordmasculine & Banco de Dados & FEELT!BD & {\scriptsize banco de dados; programação} \\
    4\textordmasculine & Estatística & FAMAT31033 & {\scriptsize probabilidade e estatística} \\
    4\textordmasculine & Métodos Matemáticos & FAMAT31031 & {\scriptsize matemática do contínuo} \\
    5\textordmasculine & Atividades Curriculares de Extensão II & ACE2 & {\scriptsize computação e sociedade} \\
    5\textordmasculine & Cálculo Numérico & FAMAT31032 & {\scriptsize matemática do contínuo; modelagem computacional} \\
    5\textordmasculine & Empreendedorismo e Inovação & FAGEN39401 & {\scriptsize empreendedorismo; fundamentos de administração} \\
    5\textordmasculine & Engenharia de Software & FEELT!ESOF & {\scriptsize análise, especificação, verificação e testes de sistemas; engenharia de software} \\
    5\textordmasculine & Fundamentos da Inteligência Artificial & FEELT!FIA & {\scriptsize inteligência artificial e computacional; lógica; programação} \\
    5\textordmasculine & Sistemas Operacionais & FEELT!SO & {\scriptsize sistemas operacionais} \\
    6\textordmasculine & Aprendizagem de Máquina & FEELT!AMAQ & {\scriptsize inteligência artificial e computacional; modelagem computacional; probabilidade e estatística} \\
    6\textordmasculine & Atividades Curriculares de Extensão III & ACE3 & {\scriptsize computação e sociedade} \\
    6\textordmasculine & Processamento Digital de Sinais & FEELT31628 & {\scriptsize matemática do contínuo; modelagem computacional} \\
    6\textordmasculine & Redes de Comunicações I & FEELT31724 & {\scriptsize avaliação de desempenho; redes de computadores} \\
    6\textordmasculine & Sistemas Embarcados I & FEELT31523 & {\scriptsize automação; sistemas de tempo real; sistemas embarcados; sistemas operacionais} \\
    6\textordmasculine & Teoria da Computação & FEELT!TCOMP & {\scriptsize algoritmos e complexidade; linguagens formais e autômatos; métodos formais; teoria da computação} \\
    7\textordmasculine & Aprendizagem Profunda e Modelos Generativos & FEELT!APMG & {\scriptsize inteligência artificial e computacional; modelagem computacional; processamento de imagens} \\
    7\textordmasculine & Atividades Curriculares de Extensão IV & ACE4 & {\scriptsize computação e sociedade} \\
    7\textordmasculine & Redes de Comunicações II & FEELT31831 & {\scriptsize avaliação de desempenho; redes de computadores} \\
    7\textordmasculine & Segurança de Sistemas Computacionais & FEELT!SEG & {\scriptsize redes de computadores; segurança; sistemas operacionais} \\
    7\textordmasculine & Sistemas Computacionais em Tempo Real & FEELT!SCTR & {\scriptsize avaliação de desempenho; dependabilidade; sistemas de tempo real; sistemas operacionais} \\
    7\textordmasculine & Sistemas e Controle & FEELT!SCON & {\scriptsize automação; modelagem computacional} \\
    8\textordmasculine & Arquitetura de Software Aplicada & FEELT!ASA & {\scriptsize engenharia de software; processamento distribuído; processamento paralelo} \\
    8\textordmasculine & Engenharia Econômica & IERI39601 & {\scriptsize fundamentos de administração; fundamentos de economia} \\
    8\textordmasculine & Engenharia de Contexto em IA Generativa & FEELT!ECIA & {\scriptsize computação e sociedade; inteligência artificial e computacional; interação humano-computador; ética e legislação} \\
    8\textordmasculine & IA Aplicada à Cibersegurança & FEELT!IACS & {\scriptsize inteligência artificial e computacional; redes de computadores; segurança} \\
    8\textordmasculine & Memorial de Atividades Curriculares de Extensão & FEELT!MACE & {\scriptsize computação e sociedade; metodologia cientifica} \\
    8\textordmasculine & Sistemas Distribuídos & FEELT!SD & {\scriptsize avaliação de desempenho; banco de dados; dependabilidade; processamento distribuído} \\
    Optativa & Compiladores & FACOM39043 & {\scriptsize compiladores} \\
    Optativa & Computação Gráfica RV-RA & FEELT31721 & {\scriptsize computação gráfica; realidade virtual} \\
    Optativa & Multimídia & GBC213 & {\scriptsize multimídia} \\
    Optativa & Otimização & FACOM31601 & {\scriptsize pesquisa operacional e otimização} \\
    Optativa & Pesquisa Operacional & FAGEN31703 & {\scriptsize pesquisa operacional e otimização} \\
    Optativa & Robótica & FEELT31722 & {\scriptsize robótica} \\*

    \end{longtblr}
    

\begin{longtblr}[
        theme = ecp,
        caption = {Distribuição dos Conteúdos Curriculares Referentes à Formação Tecnológica e Básica dos Cursos de Bacharelado em Engenharia de Computação (DCN Computação)},
        label = {tab:disciplinas_ref_dcn_ecp},
        remark{\small \textbf{Referência}} = {\small Parecer CNE/CES Nº 136/2012, item 3.3.}]{
    colspec = {Q[c,m,wd=20mm]Q[l,m,wd=50mm]Q[l,m,wd=25mm]Q[l,m,wd=40mm]},
    rowhead = 1,
    row{odd} = {bg=CinzaClaro},
    row{1} = {bg=AzulEscuro, fg=white},
    row{23} = {bg=AzulClaro, fg=white},
    cells  = {font=\fontsize{10pt}{12pt}\selectfont},
    }
        \textbf{Período} & \textbf{Componente} & \textbf{Código} & \textbf{Conteúdos}\\
        1\textordmasculine & Introdução à Prototipagem Eletrônica & FEELT!IPE & {\scriptsize circuitos elétricos; eletrônica digital; projeto de sistemas digitais; transdutores} \\
    2\textordmasculine & Física Básica: Mecânica & INFIS39031 & {\scriptsize física} \\
    2\textordmasculine & Experimental de Física Básica: Mecânica & INFIS39032 & {\scriptsize física} \\
    2\textordmasculine & Metrologia & FEELT31204 & {\scriptsize transdutores} \\
    2\textordmasculine & Sistemas Digitais & FEELT31409 & {\scriptsize eletrônica digital; projeto de sistemas digitais} \\
    2\textordmasculine & Experimental de Sistemas Digitais & FEELT31410 & {\scriptsize eletrônica digital; projeto de sistemas digitais} \\
    3\textordmasculine & Elementos de Sistemas Computacionais & FEELT!ESC & {\scriptsize eletrônica digital; microeletrônica e nanoeletrônica; projeto de circuitos integrados; projeto de sistemas digitais} \\
    3\textordmasculine & Física Básica: Eletricidade e Magnetismo & INFIS39033 & {\scriptsize eletricidade; física; teoria eletromagnética} \\
    3\textordmasculine & Experimental de Física Básica: Eletricidade e Magnetismo & INFIS39034 & {\scriptsize eletricidade; física; teoria eletromagnética} \\
    4\textordmasculine & Circuitos Elétricos I & FEELT31301 & {\scriptsize circuitos elétricos; eletricidade} \\
    4\textordmasculine & Experimental de Circuitos Elétricos I & FEELT31302 & {\scriptsize circuitos elétricos; eletricidade} \\
    5\textordmasculine & Eletrônica Analógica I & FEELT31401 & {\scriptsize circuitos elétricos; eletrônica analógica; microeletrônica e nanoeletrônica; teoria dos semicondutores} \\
    5\textordmasculine & Experimental de Eletrônica Analógica I & FEELT31402 & {\scriptsize circuitos elétricos; eletrônica analógica; microeletrônica e nanoeletrônica; teoria dos semicondutores} \\
    6\textordmasculine & Processamento Digital de Sinais & FEELT31628 & {\scriptsize processamento digital de sinais} \\
    6\textordmasculine & Redes de Comunicações I & FEELT31724 & {\scriptsize comunicação de dados} \\
    6\textordmasculine & Sistemas Embarcados I & FEELT31523 & {\scriptsize automação de projeto; eletrônica digital; microeletrônica e nanoeletrônica; projeto de sistemas digitais; sistemas de controle; transdutores} \\
    7\textordmasculine & Redes de Comunicações II & FEELT31831 & {\scriptsize comunicação de dados} \\
    7\textordmasculine & Sistemas Computacionais em Tempo Real & FEELT!SCTR & {\scriptsize automação de projeto; projeto de sistemas digitais; sistemas de controle} \\
    7\textordmasculine & Sistemas e Controle & FEELT!SCON & {\scriptsize automação de projeto; sistemas de controle} \\
    8\textordmasculine & Sistemas Distribuídos & FEELT!SD & {\scriptsize comunicação de dados} \\
    Optativa & Projetos Digitais com FPGA & FEELT!FPGA & {\scriptsize projeto de circuitos integrados} \\*

    \end{longtblr}
    


\subsection{Refente à CC2020 da ACM/IEEE}

% \section{Sobre a CC2020 da ACM}

% O \textit{Computing Curricula 2020} (CC2020) da ACM (\textit{Association for Computing Machinery}), em conjunto com outras entidades como a IEEE (\textit{Institute of Electrical and Electronics Engineers}), estabelece diretrizes para a educação em computação, diferenciando claramente os diversos campos dentro da área. Ele fornece um modelo unificado que identifica conhecimentos essenciais, habilidades e competências, garantindo que a formação acadêmica esteja alinhada às demandas do mercado e aos avanços tecnológicos.

% Uma das principais contribuições do CC2020 é a diferenciação entre os diversos cursos na área da computação, ajudando a estruturar currículos de maneira específica para cada especialidade. As principais áreas abordadas pelo CC2020 são:

%     \begin{itemize}
%         \item \textbf{\textit{Computer Engineering} (Engenharia de Computação)}: Integra hardware e software, abordando sistemas embarcados, redes, eletrônica digital, computação paralela e arquiteturas computacionais.
%         \item \textbf{\textit{Computer Science} (Ciência da Computação)}: Foca nos fundamentos teóricos e práticos da computação, incluindo algoritmos, complexidade computacional, linguagens de programação e paradigmas computacionais.
%         \item \textbf{\textit{Cybersecurity} (Cibersegurança)}: Explora a proteção de sistemas, redes e dados contra ameaças e ataques, abordando criptografia, segurança da informação, defesa cibernética e gestão de riscos.
%         \item \textbf{\textit{Information Systems} (Sistemas de Informação)}: Analisa a relação entre computação e organizações, focando na gestão de dados, processos de negócios e suporte à tomada de decisão.
%         \item \textbf{\textit{Information Technology} (Tecnologia da Informação)}: Concentra-se na infraestrutura de TI, redes, suporte técnico, administração de sistemas e segurança operacional.
%         \item \textbf{\textit{Software Engineering} (Engenharia de Software)}: Foca no desenvolvimento estruturado de software, abrangendo metodologias de desenvolvimento, qualidade, testes e escalabilidade de sistemas.
%         \item \textbf{\textit{Data Science} (Ciência de Dados)}: Ainda em evolução dentro do CC2020; explora o uso de estatística, aprendizado de máquina e engenharia de dados para extrair insights de grandes volumes de dados.
%     \end{itemize}

% Para um curso de Engenharia de Computação que segue as Diretrizes Curriculares Nacionais de 2016 (Resolução CNE/CES nº 5/2016), a observação do CC2020 é fundamental, pois ajuda a estruturar um currículo que integra hardware e software, diferenciando-o de cursos como Ciência da Computação e Engenharia de Software. Além disso, o alinhamento com as diretrizes do CC2020 assegura que os egressos tenham uma formação sólida, moderna e interdisciplinar, contemplando áreas emergentes como cibersegurança, computação paralela e distribuída, inteligência artificial e sistemas embarcados. Dessa forma, a adoção dessas diretrizes reforça a coerência curricular, moderniza o conteúdo didático, promove a internacionalização do ensino e amplia as oportunidades de empregabilidade dos alunos no mercado global da computação.

% \subsection{Competências em Engenharia de Computação}

% A área da computação abrange uma ampla gama de conhecimentos, que vão desde os impactos sociais da tecnologia até a implementação de sistemas complexos. O Quadro~\ref{tab:elementos_conhecimento} organiza esses elementos em categorias fundamentais, cobrindo aspectos que incluem segurança, modelagem de sistemas, infraestrutura, desenvolvimento de software e fundamentos de hardware. Esses tópicos são essenciais para a formação de profissionais que atuarão em diversos setores, garantindo que possuam habilidades técnicas sólidas e uma compreensão abrangente da área. Cada elemento descrito reflete tanto a teoria quanto as aplicações práticas necessárias para o desenvolvimento de soluções computacionais eficientes e inovadoras.

% % \newenvironment{ecptable}[3]{
% %     \renewcommand{\familydefault}{\sfdefault} \normalfont
% %     \begin{table}[ht]
% %     \centering
% %     \footnotesize
% %     \renewcommand{\arraystretch}{1.2} % Ajusta o espaçamento entre linhas
% %     \caption{#1}
% %     \label{#2}
% %     \begin{tabular}{|#3|}
% %     \hline
% % }{
% %     \hline
% %     \end{tabular}
% %     \renewcommand{\familydefault}{\rmdefault} \normalfont
% %     \end{table}
% % }

% \renewcommand{\familydefault}{\sfdefault} \normalfont{}
% \begin{table}[h!]
%     \centering
%     \caption{Elementos de Conhecimento em Computação \textsuperscript{\scriptsize\textcolor{AzulClaro}{CC2020, Tab.4.1}}}
%     \label{tab:elementos_conhecimento}
%     \tiny 
%     \renewcommand{\arraystretch}{1.2}
%     \begin{tabular}{|m{5.5cm}|m{8.5cm}|}
% % \begin{ecptable}{Elementos de Conhecimento em Computação (CC2020, tab.4.1)}{tab:elementos_conhecimento}{l|p{10cm}}
%     \hline
%     \rowcolor{AzulEscuro} \textcolor{white}{\textbf{Elemento de Conhecimento}} & \textcolor{white}{\textbf{Elaboração}} \\
%     \hline
%     \rowcolor{AzulClaro} \multicolumn{2}{|l|}{\textcolor{white}{\textbf{Usuários e Organizações}}} \\
%     \hline
%     $\cdot$ Questões Sociais e Prática Profissional & Impacto da tecnologia na sociedade, responsabilidades éticas dos profissionais de computação e conformidade com padrões de conduta profissional (ex: privacidade, sustentabilidade, inclusão). \\
%     \hline
%     $\cdot$ Política e Gestão de Segurança & Definição de políticas para proteção de dados, gestão de riscos e conformidade com regulamentações (ex: LGPD, ISO 27001). \\
%     \hline
%     $\cdot$ Gestão e Liderança de SI & Coordenação de equipes e recursos tecnológicos para alinhar sistemas de informação aos objetivos estratégicos da organização. \\
%     \hline
%     $\cdot$ Arquitetura Empresarial & Desenho e integração de sistemas, processos e infraestrutura para otimizar operações organizacionais. \\
%     \hline
%     $\cdot$ Gestão de Projetos & Aplicação de metodologias (ex: Agile, PMBOK) para planejar, executar e monitorar projetos de TI. \\
%     \hline
%     $\cdot$ Design de Experiência do Usuário (UX) & Criação de interfaces intuitivas e acessíveis, focadas nas necessidades e expectativas dos usuários finais. \\
%     \hline
%     \rowcolor{AzulClaro} \multicolumn{2}{|l|}{\textcolor{white}{\textbf{Modelagem de Sistemas}}} \\
%     \hline
%     $\cdot$ Questões e Princípios de Segurança & Identificação de vulnerabilidades e implementação de princípios de segurança (ex: confidencialidade, integridade) no ciclo de vida do sistema. \\
%     \hline
%     $\cdot$ Análise e Design de Sistemas & Tradução de requisitos em modelos funcionais e técnicos, utilizando técnicas como UML ou diagramas de fluxo. \\
%     \hline
%     $\cdot$ Análise de Requisitos e Especificações & Coleta, documentação e validação de necessidades de stakeholders para guiar o desenvolvimento. \\
%     \hline
%     $\cdot$ Gestão de Dados e Informação & Organização, armazenamento e recuperação de dados, incluindo uso de bancos de dados e governança de dados. \\
%     \hline
%     \rowcolor{AzulClaro} \multicolumn{2}{|l|}{\textcolor{white}{\textbf{Arquitetura e Infraestrutura de Sistemas}}} \\
%     \hline
%     $\cdot$ Sistemas e Serviços Virtuais & Implementação de ambientes virtualizados (ex: cloud computing, containers) para otimizar recursos físicos. \\
%     \hline
%     $\cdot$ Sistemas Inteligentes (IA) & Desenvolvimento de algoritmos e modelos de aprendizado de máquina para automação e tomada de decisão. \\
%     \hline
%     $\cdot$ Internet das Coisas (IoT) & Integração de dispositivos conectados à internet para coleta e processamento de dados em tempo real. \\
%     \hline
%     $\cdot$ Computação Paralela e Distribuída & Projeto de sistemas que executam tarefas simultaneamente em múltiplos processadores ou nós de rede. \\
%     \hline
%     $\cdot$ Redes de Computadores & Configuração e gerenciamento de protocolos, topologias e infraestrutura de comunicação (ex: TCP/IP, redes 5G). \\
%     \hline
%     $\cdot$ Sistemas Embarcados & Desenvolvimento de software e hardware integrados a dispositivos específicos (ex: sistemas automotivos, robótica). \\
%     \hline
%     $\cdot$ Tecnologia de Sistemas Integrados & Combinação de subsistemas heterogêneos (hardware/software) para funcionamento unificado. \\
%     \hline
%     $\cdot$ Tecnologias de Plataforma & Utilização de frameworks e ambientes pré-configurados para acelerar o desenvolvimento (ex: AWS, Kubernetes). \\
%     \hline
%     $\cdot$ Tecnologia e Implementação de Segurança & Aplicação de ferramentas e técnicas para proteger infraestrutura (ex: firewalls, criptografia). \\
%     \hline
%     \rowcolor{AzulClaro} \multicolumn{2}{|l|}{\textcolor{white}{\textbf{Desenvolvimento de Software}}} \\
%     \hline
%     $\cdot$ Qualidade de Software, Verificação e Validação & Garantia de que o software atende aos requisitos através de testes, revisões e métricas de qualidade (ex: testes unitários, CMMI). \\
%     \hline
%     $\cdot$ Processos de Software & Adoção de metodologias estruturadas para desenvolvimento (ex: DevOps, Scrum). \\
%     \hline
%     $\cdot$ Modelagem e Análise de Software & Criação de representações abstratas do sistema para identificar falhas e otimizar desempenho. \\
%     \hline
%     $\cdot$ Design de Software & Definição de arquitetura modular, padrões de projeto (ex: MVC) e boas práticas de codificação. \\
%     \hline
%     $\cdot$ Desenvolvimento Baseado em Plataformas & Utilização de ecossistemas pré-existentes (ex: Android Studio, Salesforce) para construção de aplicações. \\
%     \hline
%     \rowcolor{AzulClaro} \multicolumn{2}{|l|}{\textcolor{white}{\textbf{Fundamentos de Software}}} \\
%     \hline
%     $\cdot$ Gráficos e Visualização & Técnicas para representação visual de dados e criação de interfaces gráficas (ex: OpenGL, ferramentas de BI). \\
%     \hline
%     $\cdot$ Sistemas Operacionais & Estudo de gerenciamento de recursos (memória, processos) e interação entre hardware e software (ex: Linux, Windows). \\
%     \hline
%     $\cdot$ Estruturas de Dados, Algoritmos e Complexidade & Projeto e implementação de estruturas eficientes (ex: árvores, grafos, tabelas hash), desenvolvimento de algoritmos otimizados (ex: ordenação, busca, algoritmos em grafos) e/ou análise de eficiência computacional (ex: notação Big-O, escalabilidade, trade-offs entre tempo e espaço). \\
%     \hline
%     $\cdot$ Linguagens de Programação & Domínio de paradigmas (ex: orientação a objetos, funcional) e sintaxe de linguagens (ex: Python, Java). \\
%     \hline
%     $\cdot$ Fundamentos de Programação & Princípios básicos de lógica, sintaxe e depuração para construção de programas. \\
%     \hline
%     $\cdot$ Fundamentos de Sistemas Computacionais & Compreensão da interação entre hardware, software e redes em sistemas computacionais. \\
%     \hline
%     \rowcolor{AzulClaro} \multicolumn{2}{|l|}{\textcolor{white}{\textbf{Hardware}}} \\
%     \hline
%     $\cdot$ Arquitetura e Organização de Computadores & Projeto de componentes físicos (ex: CPU, memória) e sua interconexão para otimizar desempenho. \\
%     \hline
%     $\cdot$ Design Digital & Criação de circuitos lógicos e sistemas digitais usando portas lógicas e linguagens de descrição de hardware (ex: VHDL). \\
%     \hline
%     $\cdot$ Circuitos e Eletrônica & Estudo de componentes eletrônicos (ex: transistores, capacitores) e sua aplicação em sistemas computacionais. \\
%     \hline
%     $\cdot$ Processamento de Sinais & Técnicas para análise e manipulação de sinais analógicos/digitais (ex: filtros, compressão de áudio). \\
%     \hline
%     \end{tabular}
% \end{table}
% \renewcommand{\familydefault}{\rmdefault} \normalfont{}

%     Além do conhecimento técnico, a formação em computação exige competências fundamentais e profissionais que influenciam diretamente a atuação no mercado de trabalho. O Quadro~\ref{tab:elementos_fundamentais} destaca habilidades essenciais como pensamento crítico, colaboração, resolução de problemas e comunicação, bem como aspectos organizacionais como gestão de projetos e qualidade. Esses elementos são fundamentais para o sucesso de um profissional na área, permitindo que ele analise problemas complexos, trabalhe de maneira eficaz em equipe e se comunique com clareza em diferentes contextos. O desenvolvimento dessas habilidades complementa a formação técnica, tornando os profissionais mais preparados para enfrentar desafios multidisciplinares.

% \renewcommand{\familydefault}{\sfdefault} \normalfont{}
% \begin{table}[h!]
%     \centering
%     \caption{Elementos de Conhecimento Gerais \textsuperscript{\scriptsize\textcolor{AzulClaro}{CC2020, Tab.4.2}}}
%     \label{tab:elementos_fundamentais}
%     \tiny 
%     \renewcommand{\arraystretch}{1.2}
%     \begin{tabular}{|m{5cm}|m{9cm}|}
% \hline
% \rowcolor{AzulEscuro} \textcolor{white}{\textbf{Elemento de Conhecimento}} & \textcolor{white}{\textbf{Elaboração}} \\
% \hline
% $\cdot$ Pensamento Analítico e Crítico & Processo mental de simplificar informações complexas e tomar decisões \\
% \hline
% $\cdot$ Colaboração e Trabalho em Equipe & Dividir tarefas complexas e trabalhar em conjunto para completá-las \\
% \hline
% $\cdot$ Perspectivas Éticas e Interculturais & Visões éticas de diferentes pontos de vista em contextos humanos \\
% \hline
% $\cdot$ Matemática e Estatística & Uso abstrato de números e teorias, especialmente na análise de dados \\
% \hline
% $\cdot$ Priorização e Gestão de Múltiplas Tarefas & Organizar e priorizar múltiplas tarefas simultaneamente \\
% \hline
% $\cdot$ Comunicação Oral e Apresentação & Transmitir mensagens com eficácia usando recursos visuais e orais \\
% \hline
% $\cdot$ Resolução de Problemas e Solução de Erros & Buscar e resolver problemas de maneira lógica e organizada \\
% \hline
% $\cdot$ Organização e Planejamento de Projetos e Tarefas & Planejar e organizar projetos para atingir resultados \\
% \hline
% $\cdot$ Controle e Garantia de Qualidade & Usar métodos para identificar e prevenir defeitos \\
% \hline
% $\cdot$ Gestão de Relacionamentos & Manter o engajamento com clientes e parceiros \\
% \hline
% $\cdot$ Pesquisa e Aprendizagem Independente & Iniciar e conduzir projetos sem necessidade de direção \\
% \hline
% $\cdot$ Gestão de Tempo & Utilizar o tempo de maneira produtiva e eficiente \\
% \hline
% $\cdot$ Comunicação Escrita & Uso eficaz da escrita para comunicação entre pessoas e organizações \\
% \hline
% \end{tabular}
% \end{table}
% \renewcommand{\familydefault}{\rmdefault} \normalfont{}

% A aprendizagem em computação envolve diferentes níveis de complexidade cognitiva, que vão desde a memorização de conceitos até a criação de soluções inovadoras. A Taxonomia de Bloom organiza esses níveis de habilidade de forma hierárquica, permitindo estruturar o ensino de maneira progressiva. O Quadro~\ref{tab:habilidades_cognitivas} apresenta essas categorias, que incluem recordação, compreensão, aplicação, análise, avaliação e criação. Esse modelo é amplamente utilizado na educação para desenvolver materiais didáticos e avaliações que incentivem o aprendizado significativo, garantindo que os estudantes não apenas absorvam informações, mas também saibam utilizá-las para resolver problemas reais.

% \renewcommand{\familydefault}{\sfdefault} \normalfont{}
% \begin{table}[h!]
%     \centering
%     \caption{Níveis de Habilidade Cognitiva (Taxonomia de Bloom)  \textsuperscript{\scriptsize\textcolor{AzulClaro}{CC2020, Tab.4.3}}}
%     \label{tab:habilidades_cognitivas}
%     \tiny 
%     \renewcommand{\arraystretch}{1.2}
%     \begin{tabular}{|m{4cm}|m{10cm}|}
%     \hline
%     \rowcolor{AzulEscuro} \textcolor{white}{\textbf{Nível de Habilidade Cognitiva}} & \textcolor{white}{\textbf{Significado}} \\
%     \hline
%     $\cdot$ Recordação & Ato de relembrar materiais previamente aprendidos \\
%     \hline
%     $\cdot$ Compreensão & Capacidade de demonstrar entendimento (por meio de organização e comparação) \\
%     \hline
%     $\cdot$ Aplicação & Utilização do conhecimento para resolver problemas novos \\
%     \hline
%     $\cdot$ Análise & Exame e decomposição das informações em partes \\
%     \hline
%     $\cdot$ Avaliação & Julgamento sobre a validade e qualidade das informações \\
%     \hline
%     $\cdot$ Criação & Compilação das informações em padrões novos ou proposição de alternativas \\
%     \hline
%     \end{tabular}
%     \end{table}
%     \renewcommand{\familydefault}{\rmdefault} \normalfont{}

%     Além das habilidades técnicas e cognitivas, os profissionais da computação devem desenvolver disposições que orientam sua postura e atuação no ambiente profissional. O Quadro~\ref{tab:disposicoes_potenciais} apresenta características desejáveis, como adaptabilidade, colaboração, criatividade e responsabilidade, que impactam diretamente a eficiência e o sucesso no trabalho. Essas disposições são fundamentais para enfrentar os desafios do setor tecnológico, que exige constante inovação, trabalho em equipe e uma postura ética diante das transformações digitais. O cultivo dessas qualidades contribui para a formação de profissionais mais completos e preparados para lidar com os desafios dinâmicos da computação.

%     \renewcommand{\familydefault}{\sfdefault} \normalfont{}
% \begin{table}[h!]
%     \centering
%     \caption{Disposições Potenciais  \textsuperscript{\scriptsize\textcolor{AzulClaro}{CC2020, Tab.4.4}}}
%     \label{tab:disposicoes_potenciais}
%     \tiny 
%     \renewcommand{\arraystretch}{1.2}
%     \begin{tabular}{|m{4cm}|m{10cm}|}
%         \hline
%         \rowcolor{AzulEscuro} \textcolor{white}{\textbf{Disposição}} & \textcolor{white}{\textbf{Elaboração}} \\
%         \hline
%         $\cdot$ Adaptável & Flexível; ágil, ajusta-se em resposta à mudança \\
%         \hline
%         $\cdot$ Colaborativo & Colabora em equipe, demonstra disposição para trabalhar em conjunto \\
%         \hline
%         $\cdot$ Inventivo & Exploratório; busca soluções além do óbvio \\
%         \hline
%         $\cdot$ Meticuloso & Atento aos detalhes; minucioso, preciso \\
%         \hline
%         $\cdot$ Apaixonado & Convicção; forte comprometimento; envolvimento cativante \\
%         \hline
%         $\cdot$ Proativo & Atua com iniciativa; empreendedor; independente \\
%         \hline
%         $\cdot$ Profissional & Demonstra profissionalismo, discrição, comportamento ético e perspicácia \\
%         \hline
%         $\cdot$ Orientado por Propósito & Focado em metas; busca atingir objetivos; apresenta discernimento em termos de resultados e negócios \\
%         \hline
%         $\cdot$ Responsável & Exerce bom julgamento, usa discrição; age de maneira apropriada \\
%         \hline
%         $\cdot$ Solícito & Respeitoso; reage de forma rápida e positiva \\
%         \hline
%         $\cdot$ Autodirigido & Autônomo, motivado internamente, determinado, independente \\
%         \hline
%         \end{tabular}
%         \end{table}
%         \renewcommand{\familydefault}{\rmdefault} \normalfont{}


% pg. 108 pdf

% CE-CAE --- Circuits and Electronics
% CE-CAL — Computing Algorithms
% CE-CAO — Computer Architecture \& Organization
% CE-DIG — Digital Design
% CE-ESY — Embedded Systems
% CE-NWK — Computer Networks
% CE-PPP — Preparation for Professional Practice
% CE-SEC — Information Security
% CE-SGP — Signal Processing
% CE-SPE — Systems and Project Engineering
% CE-SRM — Systems Resource Management
% CE-SWD — Software Design

% \section{Conhecimentos em Engenharia de Computação}

% % pg. 65 pdf
% 1. Users and Organizations
% 1.1. Social Issues and Professional Practice 2 5
% --- 1.2. Security Policy and Management 1 3
% --- 1.3. IS Management and Leadership 0 2
% --- 1.4. Enterprise Architecture 0 1
% 1.5. Project Management 1 3
% --- 1.6. User Experience Design 1 3
% 2. Systems Modeling 
% --- 2.1. Security Issues and Principles 2 3
% --- 2.2. Systems Analysis \& Design 1 2
% 2.3. Requirements Analysis and Specification 1 2
% --- 2.4. Data and Information Management 1 2
% 3. Systems Architecture and Infrastructure
% 3.1. Virtual Systems and Services 1 3
% 3.2. Intelligent Systems (AI) 1 3
% *** 3.3. Internet of Things 2 4
% *** 3.4. Parallel and Distributed Computing 2 4
% 3.5. Computer Networks 2 4
% *** 3.6. Embedded Systems 3 5
% 3.7. Integrated Systems Technology 1 2
% --- 3.8. Platform Technologies 0 1
% 3.9. Security Technology and Implementation 2 3
% 4. Software Development
% 4.1. Software Quality, Verification and Validation 1 3
% 4.2. Software Process 1 2
% 4.3. Software Modeling and Analysis 1 3
% 4.4. Software Design 2 4
% 4.5. Platform-Based Development 0 2
% 5. Software Fundamentals
% 5.1. Graphics and Visualization 1 2
% 5.2. Operating Systems 2 4
% 5.3. Data Structures, Algorithms and Complexity 2 4
% 5.4. Programming Languages 2 3
% 5.5. Programming Fundamentals 2 4
% *** 5.6. Computing Systems Fundamentals 2 3
% 6. Hardware 
% *** 6.1. Architecture and Organization 4 5
% *** 6.2. Digital Design 4 5
% *** 6.3. Circuits and Electronics 4 5
% *** 6.4. Signal Processing 3 4 

No Capítulo~\ref{chp:apresentacao_acm}, foi apresentada a organização curricular de um curso de Engenharia de Computação alinhada com os elementos de conhecimento estabelecidos na CC2020. Para garantir uma formação abrangente, atualizada e globalizada, as disciplinas foram estruturadas de forma a cobrir os principais domínios da computação.
 
Nesta seção, são apresentados os quadros que relacionam cada categoria de conhecimento da computação aos componentes curriculares correspondentes do curso de \nomecursonovo, evidenciando como os conteúdos são distribuídos ao longo do curso.

Pela ordem de importância dentro da área de competências da Engenharia de Computação, temos:

\begin{enumerate}
\item Quadro~\ref{tab:cc2020_hardware}, Hardware (Tabela~\ref{tab:elementos_conhecimento}, item 6), com destaque para \textit{Arquitetura e Organização de Computadores}, \textit{Projetos Digitais}, \textit{Circuitos e Eletrônica} e \textit{Processamento de Sinais}.
\item Quadro~\ref{tab:cc2020_arquitetura_infraestrutura_sistemas}, Arquitetura e Infraestrutura de Sistemas (Tabela~\ref{tab:elementos_conhecimento}, item 3), com destaque \textit{Internet das Coisas (IoT)}, \textit{Computação Paralela e Distribuída} e \textit{Sistemas Embarcados}, além da inclusão de \textit{Inteligência Artificial} (ver Seção~\ref{sec:sobre_IA}).
\item Quadro~\ref{tab:cc2020_fundamentos_software}, Fundamentos de Software (Tabela~\ref{tab:elementos_conhecimento}, item 5), com destaque para \textit{Fundamentos de Sistemas Computacionais}.
\item Quadro~\ref{tab:cc2020_usuarios_organizacoes}, Usuários e Organizações (Tabela~\ref{tab:elementos_conhecimento}, item 1).
\item Quadro~\ref{tab:cc2020_modelagem_sistemas}, Modelagem de Sistemas (Tabela~\ref{tab:elementos_conhecimento}, item 2).
\item Quadro~\ref{tab:cc2020_desenvolvimento_software}, Desenvolvimento de Software (Tabela~\ref{tab:elementos_conhecimento}, item 4).
\end{enumerate}

Com relação à disciplinas obrigatórias do curso, eis os quadros com as classificações das mesmas:

\begin{longtblr}[
        theme = ecp,
        caption = {CC2020: Hardware},
        label = {tab:cc2020_hardware},]{
    colspec = {Q[l,m,wd=63mm]Q[l,m,wd=23mm]Q[c,m,wd=7mm]Q[c,m,wd=7mm]Q[c,m,wd=7mm]Q[c,m,wd=7mm]Q[c,m,wd=9mm]},
    rowhead = 1,
    row{odd} = {bg=CinzaClaro},
    row{1} = {bg=AzulEscuro, fg=white},
    row{13} = {bg=AzulClaro, fg=white},
    cells  = {font=\fontsize{10pt}{12pt}\selectfont},
    }
        \textbf{Componente} & \textbf{Código}  & \textbf{CH Teór.} & \textbf{CH Prát.}  & \textbf{CH EaD} & \textbf{CH Ext.} & \textbf{CH Total} \\
        Arquitetura e Organização de Computadores & FEELT!AOC & 60 & 0 & 0 & 0 & 60 \\
    Circuitos Elétricos I & FEELT31301 & 75 & 0 & 0 & 0 & 75 \\
    Experimental de Circuitos Elétricos I & FEELT31302 & 0 & 15 & 0 & 0 & 15 \\
    Elementos de Sistemas Computacionais & FEELT!ESC & 30 & 0 & 30 & 0 & 60 \\
    Eletrônica Analógica I & FEELT31401 & 60 & 0 & 0 & 0 & 60 \\
    Experimental de Eletrônica Analógica I & FEELT31402 & 0 & 30 & 0 & 0 & 30 \\
    Introdução à Prototipagem Eletrônica & FEELT!IPE & 0 & 30 & 0 & 0 & 30 \\
    Metrologia & FEELT31204 & 30 & 30 & 0 & 0 & 60 \\
    Processamento Digital de Sinais & FEELT31628 & 45 & 15 & 0 & 0 & 60 \\
    Sistemas Digitais & FEELT31409 & 30 & 0 & 0 & 0 & 30 \\
    Experimental de Sistemas Digitais & FEELT31410 & 0 & 30 & 0 & 0 & 30 \\*
    \textbf{TOTAL} &  & \textbf{330} & \textbf{150} & \textbf{30} & \textbf{0} & \textbf{510} \\
    \end{longtblr}
    

\begin{longtblr}[
        theme = ecp,
        caption = {CC2020: Arquitetura e Infraestrutura de Sistemas},
        label = {tab:cc2020_arquitetura_infraestrutura_sistemas},]{
    colspec = {Q[l,m,wd=63mm]Q[l,m,wd=23mm]Q[c,m,wd=7mm]Q[c,m,wd=7mm]Q[c,m,wd=7mm]Q[c,m,wd=7mm]Q[c,m,wd=9mm]},
    rowhead = 1,
    row{odd} = {bg=CinzaClaro},
    row{1} = {bg=AzulEscuro, fg=white},
    row{15} = {bg=AzulClaro, fg=white},
    cells  = {font=\fontsize{10pt}{12pt}\selectfont},
    }
        \textbf{Componente} & \textbf{Código}  & \textbf{CH Teór.} & \textbf{CH Prát.}  & \textbf{CH EaD} & \textbf{CH Ext.} & \textbf{CH Total} \\
        Aprendizagem Profunda e Modelos Generativos & FEELT!APMG & 30 & 15 & 0 & 0 & 45 \\
    Aprendizagem de Máquina & FEELT!AMAQ & 30 & 15 & 0 & 0 & 45 \\
    Arquitetura de Software Aplicada & FEELT!ASA & 15 & 45 & 0 & 0 & 60 \\
    Engenharia de Contexto em IA Generativa & FEELT!ECIA & 15 & 0 & 30 & 0 & 45 \\
    Fundamentos da Inteligência Artificial & FEELT!FIA & 30 & 15 & 0 & 0 & 45 \\
    IA Aplicada à Cibersegurança & FEELT!IACS & 15 & 0 & 30 & 0 & 45 \\
    Redes de Comunicações I & FEELT31724 & 45 & 15 & 0 & 0 & 60 \\
    Redes de Comunicações II & FEELT31831 & 45 & 15 & 0 & 0 & 60 \\
    Segurança de Sistemas Computacionais & FEELT!SEG & 15 & 0 & 30 & 0 & 45 \\
    Sistemas Computacionais em Tempo Real & FEELT!SCTR & 15 & 30 & 0 & 0 & 45 \\
    Sistemas Distribuídos & FEELT!SD & 45 & 0 & 0 & 0 & 45 \\
    Sistemas Embarcados I & FEELT31523 & 45 & 30 & 0 & 0 & 75 \\
    Sistemas e Controle & FEELT!SCON & 30 & 15 & 0 & 0 & 45 \\*
    \textbf{TOTAL} &  & \textbf{375} & \textbf{195} & \textbf{90} & \textbf{0} & \textbf{660} \\
    \end{longtblr}
    

\begin{longtblr}[
        theme = ecp,
        caption = {CC2020: Fundamentos de Software},
        label = {tab:cc2020_fundamentos_software},]{
    colspec = {Q[l,m,wd=63mm]Q[l,m,wd=23mm]Q[c,m,wd=7mm]Q[c,m,wd=7mm]Q[c,m,wd=7mm]Q[c,m,wd=7mm]Q[c,m,wd=9mm]},
    rowhead = 1,
    row{odd} = {bg=CinzaClaro},
    row{1} = {bg=AzulEscuro, fg=white},
    row{9} = {bg=AzulClaro, fg=white},
    cells  = {font=\fontsize{10pt}{12pt}\selectfont},
    }
        \textbf{Componente} & \textbf{Código}  & \textbf{CH Teór.} & \textbf{CH Prát.}  & \textbf{CH EaD} & \textbf{CH Ext.} & \textbf{CH Total} \\
        Análise de Algoritmos & FEELT!AA & 60 & 0 & 0 & 0 & 60 \\
    Estrutura de Dados & FEELT!ED & 60 & 0 & 0 & 0 & 60 \\
    Programação Funcional & FEELT!PF & 30 & 15 & 0 & 0 & 45 \\
    Programação Procedimental & FEELT!PP & 30 & 30 & 0 & 0 & 60 \\
    Programação Script & FEELT!PS & 30 & 30 & 0 & 0 & 60 \\
    Sistemas Operacionais & FEELT!SO & 45 & 0 & 0 & 0 & 45 \\
    Teoria da Computação & FEELT!TCOMP & 45 & 0 & 0 & 0 & 45 \\*
    \textbf{TOTAL} &  & \textbf{300} & \textbf{75} & \textbf{0} & \textbf{0} & \textbf{375} \\
    \end{longtblr}
    

\begin{longtblr}[
        theme = ecp,
        caption = {CC2020: Usuários e Organizações},
        label = {tab:cc2020_usuarios_organizacoes},]{
    colspec = {Q[l,m,wd=63mm]Q[l,m,wd=23mm]Q[c,m,wd=7mm]Q[c,m,wd=7mm]Q[c,m,wd=7mm]Q[c,m,wd=7mm]Q[c,m,wd=9mm]},
    rowhead = 1,
    row{odd} = {bg=CinzaClaro},
    row{1} = {bg=AzulEscuro, fg=white},
    row{6} = {bg=AzulClaro, fg=white},
    cells  = {font=\fontsize{10pt}{12pt}\selectfont},
    }
        \textbf{Componente} & \textbf{Código}  & \textbf{CH Teór.} & \textbf{CH Prát.}  & \textbf{CH EaD} & \textbf{CH Ext.} & \textbf{CH Total} \\
        Empreendedorismo e Inovação & FAGEN39401 & 30 & 0 & 0 & 0 & 30 \\
    Engenharia Econômica & IERI39601 & 30 & 0 & 0 & 0 & 30 \\
    Inteligência Artificial, Ética e Sociedade & FEELT!IAES & 30 & 0 & 0 & 0 & 30 \\
    Introdução à Engenharia de Computação & FEELT31103 & 30 & 0 & 0 & 0 & 30 \\*
    \textbf{TOTAL} &  & \textbf{120} & \textbf{0} & \textbf{0} & \textbf{0} & \textbf{120} \\
    \end{longtblr}
    

\begin{longtblr}[
        theme = ecp,
        caption = {CC2020: Modelagem de Sistemas},
        label = {tab:cc2020_modelagem_sistemas},]{
    colspec = {Q[l,m,wd=63mm]Q[l,m,wd=23mm]Q[c,m,wd=7mm]Q[c,m,wd=7mm]Q[c,m,wd=7mm]Q[c,m,wd=7mm]Q[c,m,wd=9mm]},
    rowhead = 1,
    row{odd} = {bg=CinzaClaro},
    row{1} = {bg=AzulEscuro, fg=white},
    row{4} = {bg=AzulClaro, fg=white},
    cells  = {font=\fontsize{10pt}{12pt}\selectfont},
    }
        \textbf{Componente} & \textbf{Código}  & \textbf{CH Teór.} & \textbf{CH Prát.}  & \textbf{CH EaD} & \textbf{CH Ext.} & \textbf{CH Total} \\
        Banco de Dados & FEELT!BD & 15 & 45 & 0 & 0 & 60 \\
    Programação Orientada a Objetos & FEELT!POO & 30 & 30 & 0 & 0 & 60 \\*
    \textbf{TOTAL} &  & \textbf{45} & \textbf{75} & \textbf{0} & \textbf{0} & \textbf{120} \\
    \end{longtblr}
    

\begin{longtblr}[
        theme = ecp,
        caption = {CC2020: Desenvolvimento de Software},
        label = {tab:cc2020_desenvolvimento_software},]{
    colspec = {Q[l,m,wd=63mm]Q[l,m,wd=23mm]Q[c,m,wd=7mm]Q[c,m,wd=7mm]Q[c,m,wd=7mm]Q[c,m,wd=7mm]Q[c,m,wd=9mm]},
    rowhead = 1,
    row{odd} = {bg=CinzaClaro},
    row{1} = {bg=AzulEscuro, fg=white},
    row{3} = {bg=AzulClaro, fg=white},
    cells  = {font=\fontsize{10pt}{12pt}\selectfont},
    }
        \textbf{Componente} & \textbf{Código}  & \textbf{CH Teór.} & \textbf{CH Prát.}  & \textbf{CH EaD} & \textbf{CH Ext.} & \textbf{CH Total} \\
        Engenharia de Software & FEELT!ESOF & 45 & 0 & 0 & 0 & 45 \\*
    \textbf{TOTAL} &  & \textbf{45} & \textbf{0} & \textbf{0} & \textbf{0} & \textbf{45} \\
    \end{longtblr}
    

\begin{longtblr}[
        theme = ecp,
        caption = {Carga Horária por Elemento de Conhecimento},%\protect\footnotemark[\value{footnote}]},
        label = {tab:ch_acm_semanal},
        note{$\ast$} = {\scriptsize A carga horária total de disciplinas de IA é de 255 horas.}
    ]{
    colspec = {Q[l,m,wd=80mm]Q[c,m,wd=30mm]Q[c,m,wd=30mm]},
    rowhead = 1,
    row{odd} = {bg=CinzaClaro},
    row{1} = {bg=AzulEscuro, fg=white},
    row{8} = {bg=gray, fg=white},
    row{13} = {bg=gray, fg=white},
    row{14} = {bg=AzulClaro, fg=white},
    cells  = {font = \fontsize{10pt}{12pt}\selectfont},
    } 
        \textbf{Conhecimento} & \textbf{Total (horas)} & \textbf{Percentual} \\
        Hardware & 510 & 14,8\% \\
    Arquitetura e Infraestrutura de Sistemas\TblrNote{$\ast$} & 660 & 19,1\% \\
    Fundamentos de Software & 375 & 10,9\% \\
    Usuários e Organizações & 120 & 3,5\% \\
    Modelagem de Sistemas & 120 & 3,5\% \\
    Desenvolvimento de Software & 45 & 1,3\% \\
   \textbf{Subtotal} &  \textbf{1830} & \textbf{53,0\%} \\    Formação Matemática e Física Aplicadas à Engenharia & 795 & 23,0\% \\
    Formação Optativa e Complementar & 180 & 5,2\% \\
    Atividades Curriculares de Extensão & 345 & 10,0\% \\
    Atividade de Conclusão de Curso & 300 & 8,7\% \\
   \textbf{Subtotal} &  \textbf{1620} & \textbf{47,0\%} \\   \textbf{TOTAL} &  \textbf{3450} & \textbf{100,0\%} \\
    \end{longtblr}

% \subsection*{Conhecimento em Hardware} Conhecimento em Hardware com dimensões:

% \listspace
% \begin{enumerate}[label=\Roman*]
%     \item \label{enum:hw-1} Arquitetura e Organização
%     \item \label{enum:hw-2} Design digital
%     \item \label{enum:hw-3} Circuitos e Eletrônica
%     \item \label{enum:hw-4} Processamento de Sinal
% \end{enumerate}

% \renewcommand{\familydefault}{\sfdefault} \normalfont
% \begin{longtblr}[
%     theme = ecp,
%     caption = {Conhecimento em Hardware},
%     label = {tab:conh_hardware},
% ]{
%   colspec = {Q[l,m,wd=65mm]Q[c,m,wd=24mm]Q[c,m,wd=10mm]Q[c,m,wd=20mm]},
%   rowhead = 1,
%   row{odd} = {bg=CinzaClaro},
%   row{1} = {bg=AzulEscuro, fg=white},
%   row{2} = {bg=AzulClaro, fg=white},
%   row{10} = {bg=AzulClaro, fg=white},
%   row{15} = {bg=AzulClaro, fg=white},
%   cells  = {font=\fontsize{10pt}{12pt}\selectfont},
% } 
%     % Componente & Código & CH Teór. & CH Prát.  & CH Rem. & CH Total \\
%     \textbf{Componente} & \textbf{Código}  & \textbf{CH Total} & \textbf{Dimensões}\\
%     \ \ \ \textit{Obrigatórias} &&&\\
%     Enriquecimento Instrumental & FEELT31109 & 30 & \ref{enum:hw-3} \\
%     Circuitos Eletro-Eletrônicos I & FEELT? & 60 & \ref{enum:hw-3}  \\
%     Circuitos Eletro-Eletrônicos I & FEELT? & 60 & \ref{enum:hw-3}  \\
%     Elementos de Sistemas Computacionais: Hardware & FEELT? & 45 & \ref{enum:hw-1},\ref{enum:hw-2}  \\
%     Arquitetura e Organização de Computadores & FEELT? & 60 & \ref{enum:hw-1}  \\
%     Sinais e Sistemas & FEELT32504 & 60 & \ref{enum:hw-4}  \\
%     Computação Embarcada & FEELT? & 75 & \ref{enum:hw-1},\ref{enum:hw-3}  \\
%     \ \ \ \textit{Optativas} &&&\\
%     Projetos Eletrônicos e Microprocessados & FEELT? & 45 & \ref{enum:hw-2},\ref{enum:hw-3} \\
%     Sistemas de Controle Realimentado & FEELT32603 & 60  & \ref{enum:hw-4} \\
%     Sistemas Computacionais em Tempo Real & FEELT31725 & 45  & \ref{enum:hw-2},\ref{enum:hw-4} \\
%     Sistemas Embarcados II & FEELT39015 & 60  & \ref{enum:hw-1},\ref{enum:hw-3} \\
%     \textbf{TOTAL} &  & \textbf{330} & 17,4\% (11,3\%|6,1\%) \\
% \end{longtblr}
% \renewcommand{\familydefault}{\rmdefault} \normalfont


% \renewcommand{\familydefault}{\sfdefault} \normalfont
% \begin{longtblr}[
%     theme = ecp,
%     caption = {Núcleo de Formação Humanística e de Extensão},
%     label = {tab:formacao_humanistica},
% ]{
%   colspec = {Q[l,m,wd=65mm]Q[c,m,wd=24mm]Q[c,m,wd=8mm]Q[c,m,wd=8mm]Q[c,m,wd=8mm]Q[c,m,wd=10mm]},
%   rowhead = 1,
%   row{odd} = {bg=CinzaClaro},
%   row{1} = {bg=AzulEscuro, fg=white},
%   row{9} = {bg=AzulClaro, fg=white},
%   cells  = {font = \fontsize{10pt}{12pt}\selectfont},
% } 
%     \textbf{Componente} & \textbf{Código}  & \textbf{CH Teór.} & \textbf{CH Prát.}  & \textbf{CH Rem.} & \textbf{CH Total} \\
%     Introdução a Engenharia de Computação & FEELT31103 & 30 & 0 & 0 & 30 \\
%     Computação e Sociedade &  & 30 & 0 & 0 & 30 \\
%     Atividades Curriculares de Extensão I &  & 0 & 60 & 0 & 60 \\
%     Atividades Curriculares de Extensão II &  & 0 & 75 & 0 & 75 \\
%     Atividades Curriculares de Extensão III &  & 0 & 90 & 0 & 90 \\
%     Atividades Curriculares de Extensão IV &  & 0 & 90 & 0 & 90 \\
%     Memorial de Atividades Curriculares de Extensão &  & 30 & 0 & 0 & 30 \\
%     \textbf{TOTAL} &  & \textbf{90} & \textbf{315} & \textbf{0} & \textbf{405} \\
% \end{longtblr}
% \renewcommand{\familydefault}{\rmdefault} \normalfont

% \renewcommand{\familydefault}{\sfdefault} \normalfont
% \begin{longtblr}[
%     theme = ecp,
%     caption = {Núcleo de Formação Tecnológica e Profissional},
%     label = {tab:formacao_tecnologica},
% ]{
%   colspec = {Q[l,m,wd=65mm]Q[c,m,wd=24mm]Q[c,m,wd=8mm]Q[c,m,wd=8mm]Q[c,m,wd=8mm]Q[c,m,wd=10mm]},
%   rowhead = 1,
%   row{odd} = {bg=CinzaClaro},
%   row{1} = {bg=AzulEscuro, fg=white},
%   row{26} = {bg=AzulClaro, fg=white},
%   cells  = {font = \fontsize{10pt}{12pt}\selectfont},
% } 
%     \textbf{Componente} & \textbf{Código}  & \textbf{CH Teór.} & \textbf{CH Prát.}  & \textbf{CH Rem.} & \textbf{CH Total} \\
%     Enriquecimento Instrumental & FEELT31109 & 0 & 30 & 0 & 30 \\
%     Programação Script & FEELT31107 & 30 & 30 & 0 & 60 \\
%     Programação Procedimental & FEELT31201 & 30 & 30 & 0 & 60 \\
%     Estrutura de Dados & FEELT? & 60 & 0 & 0 & 60 \\
%     Banco de Dados & FEELT31408 & 30 & 15 & 0 & 45 \\
%     Engenharia de Software & FEELT? & 45 & 0 & 0 & 45 \\
%     Análise de Algoritmos & FEELT? & 45 & 0 & 0 & 45 \\
%     Programação Orientada a Objetos & FEELT31514 & 30 & 30 & 0 & 60 \\
%     Redes de Comunicações I & FEELT31724 & 45 & 15 & 0 & 60 \\
%     Circuitos Eletro-Eletrônicos I & FEELT? & 60 & 0 & 0 & 60 \\
%     Teoria da Computação & FEELT? & 45 & 0 & 0 & 45 \\
%     Programação Funcional & FEELT? & 45 & 0 & 0 & 45 \\
%     Circuitos Eletro-Eletrônicos II & FEELT? & 60 & 0 & 0 & 60 \\
%     Elementos de Sistemas Computacionais: Hardware & FEELT? & 45 & 0 & 0 & 45 \\
%     Arquitetura e Organização de Computadores & FEELT? & 60 & 0 & 0 & 60 \\
%     Programação Lógica & FEELT? & 45 & 0 & 0 & 45 \\
%     Elementos de Sistemas Computacionais: Software & FEELT? & 45 & 0 & 0 & 45 \\
%     Sistemas Operacionais & FEELT? & 45 & 15 & 0 & 60 \\
%     Sinais e Sistemas & FEELT32504 & 60 & 0 & 0 & 60 \\
%     Computação Embarcada & FEELT? & 45 & 30 & 0 & 75 \\
%     Sistemas Paralelos e Distribuídos & FEELT? & 45 & 0 & 0 & 45 \\
%     Arquitetura de Software Aplicada & FEELT? & 60 & 0 & 0 & 60 \\
%     Redes de Sensores e Internet das Coisas & FEELT39057 & 45 & 15 & 0 & 60 \\
%     Atividade de Conclusão de Curso & FEELT31629 & 0 & 300 & 0 & 300 \\
%     \textbf{TOTAL} &  & \textbf{X} & \textbf{X} & \textbf{0} & \textbf{1530} \\
% \end{longtblr}
% \renewcommand{\familydefault}{\rmdefault} \normalfont




% \addtocounter{footnote}{1}
% \renewcommand{\familydefault}{\sfdefault} \normalfont
% \begin{longtblr}[
%     theme = ecp,
%     caption = {Núcleo de Formação Optativa e Complementar\protect\footnotemark[\value{footnote}]},
%     label = {tab:formacao_optativa},
% ]{
%   colspec = {Q[l,m,wd=65mm]Q[c,m,wd=24mm]Q[c,m,wd=8mm]Q[c,m,wd=8mm]Q[c,m,wd=8mm]Q[c,m,wd=10mm]},
%   rowhead = 1,
%   row{odd} = {bg=CinzaClaro},
%   row{1} = {bg=AzulEscuro, fg=white},
%   row{3} = {bg=AzulClaro, fg=white},
%   row{13} = {bg=AzulClaro, fg=white},
%   row{34} = {bg=AzulClaro, fg=white},
%   cells  = {font = \fontsize{10pt}{12pt}\selectfont},
% } 
%     \textbf{Componente} & \textbf{Código}  & \textbf{CH Teór.} & \textbf{CH Prát.}  & \textbf{CH Rem.} & \textbf{CH Total} \\
%     Atividades Acadêmicas Complementares &  & 0 & 75 & 0 & 75 \\
%     \ \ \ \textit{Optativas A} &&&&&\\
%      Língua Brasileira de Sinais - Libras I & LIBRAS01 & 30 & 30 & 0 & 60 \\
%      Educação Para o Meio Ambiente & IGUFU49010 & 30 & 0 & 0 & 30 \\
%      Empreendedorismo e Geração de Ideias & FAGEN32303 & 30 & 0 & 0 & 30 \\
%      Física Básica: Oscilações, Ondas e Óptica & INFIS31402 & 60 & 0 & 0 & 60 \\
%      Resolução de Problemas I &  &  &  & 0 &  \\
%      Tecnologias Web &  &  &  & 0 &  \\
%      Aprendizado de Máquina &  &  &  & 0 &  \\
%      Computação Gráfica RV-RA &  &  &  & 0 &  \\
%      Design Colaborativo &  &  &  & 0 &  \\
%     \ \ \ \textit{Optativas B} &&&&&\\
%      Arquitetura de Computação em Nuvem AWS &  & 0 & 0 & 45 & 45 \\
%      Segurança de Sistemas Computacionais &  &  &  & 0 &  \\
%      Tecnologias Mobile &  &  &  & 0 &  \\
%      Sistemas Computacionais de Tempo Real &  &  &  & 0 &  \\
%      Computação Evolucionária &  &  &  & 0 &  \\
%      Computação Quântica &  &  &  & 0 &  \\
%      Resolução de Problemas II &  &  &  & 0 &  \\
%      Estruturas de Dados Avançadas &  &  &  & 0 &  \\
%      Visão Computacional &  &  &  & 0 &  \\
%      Processamento de Linguagem Natural &  &  &  & 0 &  \\
%      Processamento Digital de Imagens &  &  &  & 0 &  \\
%      Algoritmos Avançados &  &  &  & 0 &  \\
%      Teoria dos Grafos &  &  &  & 0 &  \\
%      Otimização &  &  &  & 0 &  \\
%      Programação Paralela e Distribuída &  &  &  & 0 &  \\
%      Construção de Compiladores &  &  &  & 0 &  \\
%      Periféricos e Interfaces &  &  &  & 0 &  \\
%       &  &  &  & 0 &  \\
%       &  &  &  & 0 &  \\
%       &  &  &  & 0 &  \\
%     \textbf{TOTAL} &  & \textbf{X} & \textbf{X} & \textbf{X} & \textbf{X} \\
% \end{longtblr}
% \renewcommand{\familydefault}{\rmdefault} \normalfont
% \footnotetext{Da totalização indicada no quadro, o discente deve integralizar 495 horas em disciplinas optativas mais as 75 horas de Atividades Acadêmicas Complementares.}

% \renewcommand{\familydefault}{\sfdefault} \normalfont
% \begin{longtblr}[
%     theme = ecp,
%     caption = {Núcleo de Formação Complementar e Profissionalizante},
%     label = {tab:formacao_complementar},
% ]{
%   colspec = {Q[l,m,wd=55mm]Q[c,m,wd=30mm]Q[c,m,wd=10mm]Q[c,m,wd=10mm]Q[c,m,wd=10mm]Q[c,m,wd=10mm]},
%   rowhead = 1,
%   row{odd} = {bg=CinzaClaro},
%   row{1} = {bg=AzulEscuro, fg=white},
%   row{4} = {bg=AzulClaro, fg=white},
%   cells  = {font = \fontsize{10pt}{12pt}\selectfont},
% } 
%     \textbf{Componente} & \textbf{Código}  & \textbf{CH Teór.} & \textbf{CH Prát.}  & \textbf{CH Rem.} & \textbf{CH Total} \\
%     \textbf{TOTAL} &  & \textbf{0} & \textbf{75} & \textbf{0} & \textbf{75} \\
% \end{longtblr}
% \renewcommand{\familydefault}{\rmdefault} \normalfont

% \renewcommand{\familydefault}{\sfdefault} \normalfont
% \begin{longtblr}[
%     theme = ecp,
%     caption = {Distribuição de Carga Horária por Núcleo de Formação},
%     label = {tab:sintese_nucleo},
% ]{
%   colspec = {Q[l,m,wd=80mm]Q[c,m,wd=20mm]Q[c,m,wd=20mm]},
%   rowhead = 1,
%   row{odd} = {bg=CinzaClaro},
%   row{1} = {bg=AzulEscuro, fg=white},
%   row{6} = {bg=AzulClaro, fg=white},
%   cells  = {font = \fontsize{10pt}{12pt}\selectfont},
% } 
%     \textbf{Componentes Curriculares} & \textbf{CH Total} & \textbf{Percentual} \\
%     Núcleo de Formação Básica & 945 & 27,4\% \\
%     Núcleo de Formação Humanística e de Extensão & 405 & 11,7\% \\
%     Núcleo de Formação Tecnológica e Profissional & 1530 & 44,4\% \\
%     Núcleo de Formação Optativa e Complementar & 570 & 16,5\% \\
%     \textbf{TOTAL} &  \textbf{3450} & \textbf{100,0\%} \\
% \end{longtblr}
% \renewcommand{\familydefault}{\rmdefault} \normalfont


% \renewcommand{\familydefault}{\sfdefault} \normalfont
% \begin{longtblr}[
%     theme = ecp,
%     caption = {Distribuição de Carga Horária por Componente Curricular},
%     label = {tab:sintese_ch},
% ]{
%   colspec = {Q[l,m,wd=80mm]Q[c,m,wd=20mm]Q[c,m,wd=20mm]},
%   rowhead = 1,
%   row{odd} = {bg=CinzaClaro},
%   row{1} = {bg=AzulEscuro, fg=white},
%   row{7} = {bg=AzulClaro, fg=white},
%   cells  = {font = \fontsize{10pt}{12pt}\selectfont},
% } 
%     \textbf{Componentes Curriculares} & \textbf{CH Total} & \textbf{Percentual} \\
%     Disciplinas Obrigatórias & 2235 & 64,8\% \\
%     Disciplinas Optativas & 495 & 14,3\% \\
%     Atividades Curriculares de Extensão & 345 & 10,0\% \\
%     Atividade de Conclusão de Curso & 300 & 8,7\% \\
%     Atividades Acadêmicas Complementares & 75 & 2,2\% \\
%     \textbf{TOTAL} &  \textbf{3450} & \textbf{100,0\%} \\
% \end{longtblr}
% \renewcommand{\familydefault}{\rmdefault} \normalfont

\newpage

\section{Fluxo Curricular}

O fluxo curricular do curso é apresentado de modo a \textbf{facilitar a compreensão da trajetória acadêmica} pelos estudantes e demais leitores deste Projeto Pedagógico, evidenciando a sequência recomendada das disciplinas em cada período letivo, bem como seus pré-requisitos, natureza e carga horária correspondente.

O Quadro~\ref{tab:fluxo_curricular} foi organizado com fins didáticos e orientativos, buscando tornar mais clara a progressão natural dos estudos e o equilíbrio entre os eixos de formação ao longo do curso. Cada coluna representa:

\begin{itemize}[itemsep=0pt]
    \item \textbf{PER}: período ideal para cursar a disciplina;
    \item \textbf{Componente Curricular}: denominação de cada disciplina ou atividade;
    \item \textbf{Natureza}: especifica se a disciplina é obrigatória ou optativa;
    \item \textbf{CHT}: carga horária teórica;
    \item \textbf{CHP}: carga horária prática;
    \item \textbf{CHD}: carga horária a ser cumprida a distância (EaD);
    \item \textbf{CHE}: carga horária de extensão;
    \item \textbf{TOT}: carga horária total da disciplina;
    \item \textbf{PREQ}: pré-requisitos necessários (quando existentes, ou ``livre'' se não possui);
    \item \textbf{CREQ}: co-requisitos (quando existentes, ou ``livre'' se não possui); e
    \item \textbf{UA Oferta}: Unidade Acadêmica responsável pela oferta da disciplina.
\end{itemize}

\input{include/auto/tab_fluxo_curricular.tex}

\newpage

\section{Fluxo Curricular conforme Guia de Orientações Gerais}

% O fluxo curricular do curso é aqui \textbf{apresentado novamente}, conforme Guia de Orientações Gerais~\cite{UFU:GuiaPPC:2021}, de forma a ser reconhecível por esta Universidade, evidenciando a sequência recomendada das disciplinas em cada período letivo, bem como seus pré-requisitos e a carga horária correspondente. No Quadro~\ref{tab:fluxo_curricular_wrt_guia}, cada coluna representa:

Em atendimento ao que dispõe o Guia de Orientações Gerais para Elaboração de Projetos Pedagógicos de Cursos de Graduação da UFU (3ª edição, 2021)~\cite{UFU:GuiaPPC:2021}, apresenta-se a seguir o fluxo curricular do curso em formato adequado às normas institucionais atualmente vigentes.

Este segundo quadro tem caráter formal e normativo, assegurando a compatibilidade do documento com os padrões adotados pela Universidade enquanto não houver deliberação específica acerca da contabilização das cargas horárias de ensino a distância (CHD) e de extensão (CHE).

O Quadro~\ref{tab:fluxo_curricular_wrt_guia} segue, portanto, o formato reconhecido pelo Guia, e cada coluna representa:

\begin{itemize}[itemsep=0pt]
    \item \textbf{PER}: período ideal para cursar a disciplina;
    \item \textbf{Componente Curricular}: denominação de cada disciplina ou atividade;
    \item \textbf{Natureza}: especifica se a disciplina é obrigatória ou optativa;
    \item \textbf{CHT}: carga horária teórica;
    \item \textbf{CHP}: carga horária prática (que pode incluir atividades EaD ou de extensão);
    \item \textbf{TOT}: carga horária total da disciplina;
    \item \textbf{PREQ}: pré-requisitos necessários (quando existentes, ou ``livre'' se não possui);
    \item \textbf{CREQ}: co-requisitos (quando existentes, ou ``livre'' se não possui); e
    \item \textbf{UA Oferta}: Unidade Acadêmica responsável pela oferta da disciplina.
\end{itemize}

\begin{longtblr}[
        theme = ecp,
        caption = {Fluxo Curricular (Guia de Orientações)},%\protect\footnotemark[\value{footnote}]},
        label = {tab:fluxo_curricular_wrt_guia},
    note{$1$} = {\scriptsize Os discentes deverão integralizar \textbf{345 horas} de atividades extensionistas (ACE) ao longo do curso, incluindo defesa de memorial em disciplina específica, com distribuição sugerida neste fluxo como \textit{Atividades Curriculares de Extensão}.},
note{$2$} = {\scriptsize O Exame Nacional de Desempenho dos Estudantes (ENADE) integra o Sistema Nacional de Avaliação da Educação Superior (Sinaes) e é componente curricular obrigatório, conforme Lei n\textordmasculine{} 10.861, de 14 de abril de 2004.},
note{$3$} = {\scriptsize Para integralização curricular, o discente deverá cursar \textbf{90 horas} horas de atividades acadêmicas complementares (AAC) ao longo do curso.},
note{$4$} = {\scriptsize Para a realização do mínimo de 300 horas na Atividade de Conclusão do Curso (Estágio Supervisionado ou do Trabalho de Conclusão de Curso, opção possibilitada pela Resolução CNE/CES nº 5/2016 em seus artigos 7º e 8º), o discente deverá ter integralizado, no mínimo, \textbf{1875 horas} em disciplinas obrigatórias, o equivalente a todas do 1\textordmasculine{} ao 5\textordmasculine{} períodos.},
note{$5$} = {\scriptsize O discente deverá cursar, no mínimo, \textbf{90 horas} em disciplinas optativas (ver relação na \autoref{sec:optativas}) para integralização curricular.},
]{
    colspec = {|Q[c,m,wd=3mm]|Q[l,m,wd=20mm]|Q[c,m,wd=13mm]|Q[c,m,wd=5mm]|Q[c,m,wd=5mm]|Q[c,m,wd=5mm]|Q[c,m,wd=19mm]|Q[c,m,wd=19mm]|Q[c,m,wd=8mm]|}, %140mm
    %lastfoot = {\hline},
    rowsep = 1pt,
    rowhead = 2,
    hlines = {fg=AzulEscuro},
    vlines = {fg=AzulEscuro},
    row{odd} = {bg=CinzaClaro},
    row{1} = {bg=AzulEscuro, fg=white},
    row{2} = {bg=AzulEscuro, fg=white},
    row{Z} = {ht=1mm},
    cells  = {font = \fontsize{8pt}{8pt}\selectfont}
    } 
    
    % Cabeçalho Principal
    \SetCell[r=2]{m}\textbf{PER} & 
    \SetCell[r=2]{c,m}\textbf{Componente Curricular} & 
    \SetCell[r=2]{m}\textbf{Natureza}  & 
    \SetCell[c=3]{c,m}\textbf{Carga Horária} & & & 
    \SetCell[c=2]{m}\textbf{Requisitos} & & 
    \SetCell[r=2]{m}\textbf{UA Oferta}\\ \hline[AzulEscuro]
        
    % Segunda linha do cabeçalho (divisão das cargas horárias)
     &  &  & \textbf{CHT} & \textbf{CHP} & \textbf{TOT} & \textbf{PREQ} & \textbf{CREQ} & \\
    \hline[AzulEscuro]
                \SetCell[r=7]{h,bg=white} 1\textordmasculine & Cálculo Diferencial e Integral I & Obrigatória & 90 & 0 & 90 & Livre & Livre & FAMAT\\
 & Geometria Analítica & Obrigatória & 60 & 0 & 60 & Livre & Livre & FAMAT\\
 & Introdução à Engenharia de Computação & Obrigatória & 30 & 0 & 30 & Livre & Livre & FEELT\\
 & Introdução à Prototipagem Eletrônica & Obrigatória & 0 & 30 & 30 & Livre & Livre & FEELT\\
 & Lógica e Matemática Discreta para Computação & Obrigatória & 45 & 0 & 45 & Livre & Livre & FEELT\\
 & Programação Script & Obrigatória & 30 & 30 & 60 & Livre & Livre & FEELT\\
 & ENADE (Ingressante)\TblrNote{$2$} & Obrigatória & -- & -- & -- & Livre & Livre & --\\
\hline[AzulEscuro]
                \SetCell[r=8]{h,bg=white} 2\textordmasculine & Álgebra Linear & Obrigatória & 45 & 0 & 45 & Geometria Analítica & Livre & FAMAT\\
 & Cálculo Diferencial e Integral II & Obrigatória & 90 & 0 & 90 & Cálculo Diferencial e Integral I & Livre & FAMAT\\
 & Física Básica: Mecânica & Obrigatória & 60 & 0 & 60 & Cálculo Diferencial e Integral I & Livre & INFIS\\
 & Experimental de Física Básica: Mecânica & Obrigatória & 0 & 30 & 30 & Livre & Física Básica: Mecânica & INFIS\\
 & Metrologia & Obrigatória & 30 & 30 & 60 & Livre & Livre & FEELT\\
 & Programação Procedimental & Obrigatória & 30 & 30 & 60 & Programação Script & Livre & FEELT\\
 & Sistemas Digitais & Obrigatória & 30 & 0 & 30 & Lógica e Matemática Discreta para Computação & Livre & FEELT\\
 & Experimental de Sistemas Digitais & Obrigatória & 0 & 30 & 30 & Livre & Sistemas Digitais & FEELT\\
\hline[AzulEscuro]
                \SetCell[r=8]{h,bg=white} 3\textordmasculine & Cálculo Diferencial e Integral III & Obrigatória & 90 & 0 & 90 & Cálculo Diferencial e Integral II & Livre & FAMAT\\
 & Elementos de Sistemas Computacionais (EaD) & Obrigatória & 30 & 30 & 60 & Sistemas Digitais & Livre & FEELT\\
 & Estrutura de Dados & Obrigatória & 60 & 0 & 60 & Programação Procedimental & Livre & FEELT\\
 & Física Básica: Eletricidade e Magnetismo & Obrigatória & 60 & 0 & 60 & Física Básica: Mecânica & Livre & INFIS\\
 & Experimental de Física Básica: Eletricidade e Magnetismo & Obrigatória & 0 & 30 & 30 & Livre & Física Básica: Eletricidade e Magnetismo & INFIS\\
 & Inteligência Artificial, Ética e Sociedade & Obrigatória & 30 & 0 & 30 & Livre & Livre & FEELT\\
 & Programação Funcional & Obrigatória & 30 & 15 & 45 & Livre & Livre & FEELT\\
 & Programação Orientada a Objetos & Obrigatória & 30 & 30 & 60 & Programação Procedimental & Livre & FEELT\\
\hline[AzulEscuro]
                \SetCell[r=8]{h,bg=white} 4\textordmasculine & Análise de Algoritmos & Obrigatória & 60 & 0 & 60 & Estrutura de Dados & Livre & FEELT\\
 & Arquitetura e Organização de Computadores & Obrigatória & 60 & 0 & 60 & Elementos de Sistemas Computacionais & Livre & FEELT\\
 & Banco de Dados & Obrigatória & 15 & 45 & 60 & Estrutura de Dados & Livre & FEELT\\
 & Circuitos Elétricos I & Obrigatória & 75 & 0 & 75 & Física Básica: Eletricidade e Magnetismo & Livre & FEELT\\
 & Experimental de Circuitos Elétricos I & Obrigatória & 0 & 15 & 15 & Metrologia & Circuitos Elétricos I & FEELT\\
 & Estatística & Obrigatória & 60 & 0 & 60 & Livre & Livre & FAMAT\\
 & Métodos Matemáticos & Obrigatória & 75 & 0 & 75 & Cálculo Diferencial e Integral III & Livre & FAMAT\\
 & Atividades Curriculares de Extensão I\TblrNote{$1$} & Obrigatória & 0 & 90 & 90 & Livre & Livre & FEELT\\
\hline[AzulEscuro]
                \SetCell[r=8]{h,bg=white} 5\textordmasculine & Cálculo Numérico & Obrigatória & 60 & 0 & 60 & Cálculo Diferencial e Integral III; Álgebra Linear & Livre & FAMAT\\
 & Eletrônica Analógica I & Obrigatória & 60 & 0 & 60 & Circuitos Elétricos I & Livre & FEELT\\
 & Experimental de Eletrônica Analógica I & Obrigatória & 0 & 30 & 30 & Livre & Eletrônica Analógica I & FEELT\\
 & Empreendedorismo e Inovação & Obrigatória & 30 & 0 & 30 & Livre & Livre & FAGEN\\
 & Engenharia de Software & Obrigatória & 45 & 0 & 45 & Banco de Dados; Programação Orientada a Objetos & Livre & FEELT\\
 & Fundamentos da Inteligência Artificial & Obrigatória & 30 & 15 & 45 & Estatística & Livre & FEELT\\
 & Sistemas Operacionais & Obrigatória & 45 & 0 & 45 & Programação Procedimental; Arquitetura e Organização de Computadores & Livre & FEELT\\
 & Atividades Curriculares de Extensão II\TblrNote{$1$} & Obrigatória & 0 & 90 & 90 & Atividades Curriculares de Extensão I & Livre & FEELT\\
\hline[AzulEscuro]
                \SetCell[r=6]{h,bg=white} 6\textordmasculine & Aprendizagem de Máquina & Obrigatória & 30 & 15 & 45 & Cálculo Numérico; Fundamentos da Inteligência Artificial & Livre & FEELT\\
 & Processamento Digital de Sinais & Obrigatória & 45 & 15 & 60 & Métodos Matemáticos & Livre & FEELT\\
 & Redes de Comunicações I & Obrigatória & 45 & 15 & 60 & Sistemas Operacionais & Livre & FEELT\\
 & Sistemas Embarcados I & Obrigatória & 45 & 30 & 75 & Eletrônica Analógica I; Sistemas Operacionais & Livre & FEELT\\
 & Teoria da Computação & Obrigatória & 45 & 0 & 45 & Análise de Algoritmos & Livre & FEELT\\
 & Atividades Curriculares de Extensão III\TblrNote{$1$} & Obrigatória & 0 & 75 & 75 & Atividades Curriculares de Extensão II & Livre & FEELT\\
\hline[AzulEscuro]
                \SetCell[r=6]{h,bg=white} 7\textordmasculine & Aprendizagem Profunda e Modelos Generativos & Obrigatória & 30 & 15 & 45 & Aprendizagem de Máquina & Livre & FEELT\\
 & Redes de Comunicações II & Obrigatória & 45 & 15 & 60 & Redes de Comunicações I & Livre & FEELT\\
 & Segurança de Sistemas Computacionais (EaD) & Obrigatória & 15 & 30 & 45 & Redes de Comunicações I & Livre & FEELT\\
 & Sistemas Computacionais em Tempo Real & Obrigatória & 15 & 30 & 45 & Sistemas Embarcados I & Livre & FEELT\\
 & Sistemas e Controle & Obrigatória & 30 & 15 & 45 & Processamento Digital de Sinais & Livre & FEELT\\
 & Atividades Curriculares de Extensão IV\TblrNote{$1$} & Obrigatória & 0 & 60 & 60 & Atividades Curriculares de Extensão III & Livre & FEELT\\
\hline[AzulEscuro]
                \SetCell[r=7]{h,bg=white} 8\textordmasculine & Arquitetura de Software Aplicada & Obrigatória & 15 & 45 & 60 & Engenharia de Software; Redes de Comunicações I & Livre & FEELT\\
 & Engenharia Econômica & Obrigatória & 30 & 0 & 30 & Livre & Livre & IERI\\
 & Engenharia de Contexto em IA Generativa (EaD) & Obrigatória & 15 & 30 & 45 & Aprendizagem Profunda e Modelos Generativos & Livre & FEELT\\
 & IA Aplicada à Cibersegurança (EaD) & Obrigatória & 15 & 30 & 45 & Segurança de Sistemas Computacionais; Aprendizagem Profunda e Modelos Generativos & Livre & FEELT\\
 & Sistemas Distribuídos & Obrigatória & 45 & 0 & 45 & Redes de Comunicações I & Livre & FEELT\\
 & Memorial de Atividades Curriculares de Extensão\TblrNote{$1$} & Obrigatória & 0 & 30 & 30 & Atividades Curriculares de Extensão IV & Livre & FEELT\\
 & ENADE (Concluinte)\TblrNote{$2$} & Obrigatória & -- & -- & -- & Livre & Livre & --\\
\hline[AzulEscuro]
                \SetCell[c=2]{l}{Atividades Acadêmicas Complementares\TblrNote{$3$}} & & Obrigatória & -- & -- & 90 & Livre & Livre & --\\
\hline[AzulEscuro]
                \SetCell[c=2]{l}{Atividade de Conclusão de Curso\TblrNote{$4$}} & & Obrigatória & -- & -- & 300 & \textit{1875 horas} & Livre & --\\
\hline[AzulEscuro]
                \SetCell[c=2]{l}{Disciplinas Optativas\TblrNote{$5$}} & & Optativa & -- & -- & 90 & Livre & Livre & --\\
\hline[AzulEscuro]
                \SetCell[r=8]{h,bg=white} \rotatebox[origin=r]{90}{Seleção de Optativas (ver Seção~\ref{sec:optativas} para mais)} & Computação Evolutiva & Optativa & 30 & 15 & 45 & Fundamentos da Inteligência Artificial & Livre & FEELT\\
 & Computação Gráfica RV-RA & Optativa & 30 & 15 & 45 & Álgebra Linear & Livre & FEELT\\
 & Eletrônica Analógica II & Optativa & 60 & 0 & 60 & Eletrônica Analógica I & Livre & FEELT\\
 & Experimental de Eletrônica Analógica II & Optativa & 0 & 30 & 30 & Livre & Eletrônica Analógica II & FEELT\\
 & Língua Brasileira de Sinais - Libras I & Optativa & 30 & 30 & 60 & Livre & Livre & FACED\\
 & Projetos de Hardware e Firmware para IOT & Optativa & 0 & 45 & 45 & Sistemas Embarcados I & Livre & FEELT\\
 & Robótica & Optativa & 45 & 15 & 60 & Sistemas e Controle & Livre & FEELT\\
 & Sistemas Embarcados II & Optativa & 30 & 30 & 60 & Sistemas Embarcados I & Livre & FEELT\\
%\SetRow{bg=AzulEscuro} & & & & & & & & \\
    \end{longtblr}

\newpage

\section{Regime e Tempo de Integralização}
\label{sec:regime}

% O Curso de \nomecursonovo é oferecido em \textbf{regime semestral} e \textbf{período vespertino}~\cite{MEC:Portaria:21:2017}. De acordo com a definição legal, o curso vespertino é aquele em que a maior parte da carga horária é oferecida entre 12h e 18h todos os dias da semana. No curso de \nomecursonovo, a maior parte da carga horária segue esse turno, mas aulas podem ser alocadas nos períodos da manhã e da noite, de acordo com a necessidade, incluindo manhãs de sábado. É requerido que o(a) estudante tenha ciência que as atividades extra-classes precisam de dedicação do estudante e, mesmo que não contabilizadas na carga horária total do curso, serão significativas no seu percurso formativo.

% O(a) estudante tem um prazo mínimo de \textbf{oito semestres} e um prazo máximo de \textbf{doze semestres} para a integralização regulamentar do curso. O Quadro~\ref{tab:ch_semanal_por_periodo} a seguir mostra a carga horária semanal de disciplinas obrigatórias e optativas, além das cargas horárias dos componentes de Atividades Curriculares de Extensão (ACE), Atividade de Conclusão do Curso (ACC) e Atividades Acadêmicas Complementares (AAC).

O Curso de Graduação em \nomecursonovo é oferecido em \textbf{regime semestral}, com \textit{a maior parte de sua carga horária} concentrada no período vespertino, conforme definido pela legislação~\cite{MEC:Portaria:21:2017}. No entanto, o aluno deve estar ciente de que disciplinas regulares do currículo poderão ser ofertadas em seu contraturno (no período da manhã ou da noite), inclusive aos sábados pela manhã, para viabilizar sua integralização curricular. Dessa forma, o curso opera em um \textbf{regime de turno vespertino com oferta flexível}, exigindo disponibilidade do estudante em outros horários para a conclusão de componentes curriculares obrigatórios, atendimento docente ou outras atividades acadêmicas complementares. Além disso, espera-se do(a) estudante um comprometimento significativo com estudos autônomos, trabalhos em grupo e projetos, atividades essenciais para a consolidação da aprendizagem, ainda que essas horas não componham a carga horária formal do curso.

O prazo para integralização curricular é de mínimo de \textbf{oito semestres} e máximo de \textbf{doze}. A distribuição da carga horária ao longo dos períodos, incluindo disciplinas obrigatórias, optativas, Atividades Curriculares de Extensão (ACE), Atividade de Conclusão do Curso (ACC) e Atividades Acadêmicas Complementares (AAC), é detalhada no Quadro~\ref{tab:ch_semanal_por_periodo}.



% \addtocounter{footnote}{1}
% \renewcommand{\familydefault}{\sfdefault} \normalfont
% \begin{longtblr}[
%     theme = ecp,
%     caption = {Carga Horária Semanal por Período},%\protect\footnotemark[\value{footnote}]},
%     label = {tab:ch_semanal},
% ]{
%   colspec = {Q[c,m,wd=30mm]Q[c,m,wd=30mm]Q[c,m,wd=30mm]},
%   rowhead = 1,
%   row{odd} = {bg=CinzaClaro},
%   row{1} = {bg=AzulEscuro, fg=white},
%   row{13} = {bg=AzulClaro, fg=white},
%   cells  = {font = \fontsize{10pt}{12pt}\selectfont},
% } 
%     \textbf{Período} & \textbf{CH Semanal (horas)} & \textbf{Total (horas)} \\
%     1º & 20 & 300 \\
%     2º & 24 & 360 \\
%     3º & 24 & 360 \\
%     4º & 25 & 375 \\
%     5º & 25 & 375 \\
%     6º & 23 & 345 \\
%     7º & 22 & 330 \\
%     8º & 19\protect\footnotemark[\value{footnote}] & 285\protect\footnotemark[\value{footnote}] \\
%     ACE & {\footnotesize distribuida} & 345\\
%     ACC & {\footnotesize distribuida} & 300\\
%     AAC & {\footnotesize distribuida} & 75\\
%     \textbf{TOTAL} &   & \textbf{3450} \\
% \end{longtblr}
% \renewcommand{\familydefault}{\rmdefault} \normalfont
% % \footnotetext{Exceto Atividades Curriculares de Extensão, Atividades Acadêmicas Complementares e Atividade de Conclusão de Curso.}
% \footnotetext{Carga Horária Semanal de 21 horas com Total de 315 horas, se considerarmos as 30 horas teóricas da disciplina ``Memorial de Atividades Curriculares de Extensão'' incluída em ACE.}
\begin{longtblr}[
        theme = ecp,
        caption = {Carga Horária Semanal por Período},%\protect\footnotemark[\value{footnote}]},
        label = {tab:ch_semanal_por_periodo}
    ]{
    colspec = {Q[c,m,wd=30mm]Q[c,m,wd=30mm]Q[c,m,wd=30mm]Q[c,m,wd=30mm]},
    rowhead = 1,
    row{odd} = {bg=CinzaClaro},
    row{1} = {bg=AzulEscuro, fg=white},
    row{14} = {bg=AzulClaro, fg=white},
    cells  = {font = \fontsize{10pt}{12pt}\selectfont},
    } 
        \textbf{Período} & \textbf{CH Semanal (horas)} & \textbf{Total (horas)} & \textbf{Percentual} \\
        1\textordmasculine{} & 21 & 315 & 9,1\% \\
    2\textordmasculine{} & 27 & 405 & 11,7\% \\
    3\textordmasculine{} & 29 & 435 & 12,6\% \\
    4\textordmasculine{} & 27 & 405 & 11,7\% \\
    5\textordmasculine{} & 21 & 315 & 9,1\% \\
    6\textordmasculine{} & 19 & 285 & 8,3\% \\
    7\textordmasculine{} & 16 & 240 & 7,0\% \\
    8\textordmasculine{} & 15 & 225 & 6,5\% \\
    Optativas & variada & 90 & 2,6\% \\
    ACE & variada & 345 & 10,0\% \\
    ACC & variada & 300 & 8,7\% \\
    AAC & variada & 90 & 2,6\% \\
   \textbf{TOTAL} &  & \textbf{3450} & \textbf{100,0\%} \\
    \end{longtblr}


% \clearpage

% \KOMAoptions{paper=a3,paper=landscape}
% \recalctypearea{}

% \newpage
% \section{Representação Gráfica do Currículo} %https://tex.stackexchange.com/questions/168748/optimal-arrangement-of-pictures-boxes-in-a-page
% \todo{A redigir}

\newgeometry{margin=1cm}
\clearpage

\thispagestyle{empty} % Removes headers/footers from this page
% %%%%%%%%%%%%%%%%%%%%%%%%%%%%%%%%%%%%%%%%%%%%%%%%%%%%%%%%%%%%%%%%%%%%%%
%     % \hspace{-1.3cm}
%     % \begin{table}
%     %     \centering
%     %     \caption*{{\scriptsize Representação Gráfica:} \textbf{\nomecursonovo}}
%     \begin{table}
%     \centering
%     \rotatebox{90}{
    %         % \begin{center}
    
    %            % \begin{landscape}
    %             \resizebox{0.85\pdfpageheight}{!}{
        %                 \nohyph{}
        %                 % \begin{minipage}          
        %                 % \paragraph{Representação Gráfica:} \textbf{\nomecursonovo}
        \newcommand{\ccheader}[6]{
    \begin{tikzpicture}
        \draw (-0.2,1.5) node[font=\itshape\sffamily\scriptsize, yshift=-7pt, anchor=center, align=right, text width=2ex]{};
        \draw (0,0.5) rectangle node[font=\bfseries\sffamily\scriptsize, text width=3cm, align=center]{#1Período} (3,1);
        
        \draw (0,0) rectangle node[font=\bfseries\sffamily\tiny, text width=0.75cm, align=center]{CHT} (0.6,0.5);
        \draw (0.6,0) rectangle node[font=\bfseries\sffamily\tiny, text width=0.6cm, align=center]{CHP} (1.2,0.5);
        \draw (1.2,0) rectangle node[font=\bfseries\sffamily\tiny, text width=0.6cm, align=center]{CHD} (1.8,0.5);
        \draw (1.8,0) rectangle node[font=\bfseries\sffamily\tiny, text width=0.6cm, align=center]{CHE} (2.4,0.5);
        \draw (2.4,0) rectangle node[font=\bfseries\sffamily\tiny, text width=0.6cm, align=center]{TOT} (3,0.5);
        
        \draw (0,-0.5) rectangle node[font=\bfseries\sffamily\scriptsize, text width=0.6cm, align=center]{#2} (0.6,0);
        \draw (0.6,-0.5) rectangle node[font=\bfseries\sffamily\scriptsize, text width=0.6cm, align=center]{#3} (1.2,0);
        \draw (1.2,-0.5) rectangle node[font=\bfseries\sffamily\scriptsize, text width=0.6cm, align=center]{#4} (1.8,0);
        \draw (1.8,-0.5) rectangle node[font=\bfseries\sffamily\scriptsize, text width=0.6cm, align=center]{#5} (2.4,0);
        \draw (2.4,-0.5) rectangle node[font=\bfseries\sffamily\scriptsize, text width=0.6cm, align=center]{#6} (3,0);
        
    \end{tikzpicture}
    }
    
    %\newcounter{contbloco}
    %\setcounter{contbloco}{0}

    \newcommand{\ccbloco}[9]{
    %\addtocounter{contbloco}{1}
    \begin{tikzpicture}
        \draw (0,0.5) rectangle node[font=\sffamily\scriptsize, yshift=-2pt, text width=2.8cm, align=center]{#2} (3,2.5);
        \draw (1.5,2.5) node[font=\bfseries\sffamily\scriptsize, yshift=-7pt, anchor=mid]{#1};
        \draw (0,0) rectangle node[font=\sffamily\scriptsize, text width=1cm, align=center]{#3} (0.6,0.5);
        \draw (0.6,0) rectangle node[font=\sffamily\scriptsize, text width=1cm, align=center]{#4} (1.2,0.5);
        \draw (1.2,0) rectangle node[font=\sffamily\scriptsize, text width=1cm, align=center]{#5} (1.8,0.5);
        \draw (1.8,0) rectangle node[font=\sffamily\scriptsize, text width=1cm, align=center]{#6} (2.4,0.5);
        \draw (2.4,0) rectangle node[font=\sffamily\scriptsize, text width=1cm, align=center]{#7} (3,0.5);
        \ifthenelse{\isempty{#8}}
        {
        \draw (-0.475,1.75) node[font=\sffamily\scriptsize, yshift=-7pt, anchor=center, align=right, text width=2ex]{};
        }
        {
        \draw (-0.475,1.75) node[font=\sffamily\scriptsize, yshift=-7pt, anchor=center, align=right, text width=2ex]{\textbf{#8}};
        \draw (-0.10,1.75) node[font=\sffamily\footnotesize, yshift=-7pt, anchor=mid, rotate=90]{$\boldsymbol{\blacktriangledown}$};
        }
        \ifthenelse{\isempty{#9}}
        {
        \draw (-0.475,0.5) node[font=\sffamily\scriptsize, yshift=-7pt, anchor=center, align=right, text width=2ex]{};
        }
        {
        \draw (-0.475,0.5) node[font=\sffamily\scriptsize, yshift=-7pt, anchor=center, align=right, text width=2ex]{\textbf{#9}};
        \draw (-0.30,0.5) node[font=\sffamily\footnotesize, yshift=-7pt, anchor=mid, rotate=90]{\rotatebox{180}{$\boldsymbol{\bigstar}$}};
        }
    \end{tikzpicture}
    }

    
    
    \begin{table}
    \refstepcounter{section}
    \addcontentsline{toc}{section}{\protect\numberline{\thesection}Representação Gráfica do Fluxo Curricular}
    \centering
    \rotatebox{90}{\resizebox{0.85\pdfpageheight}{!}{
    \nohyph{}
    \begin{tabular}{||cccccccc|c||}
    \hline\hline
    \multicolumn{9}{l}{{\fontfamily{phv}\selectfont{\scriptsize \framebox[1.1\width]{Seção~\thesection} PPC \ppcversao - Representação Gráfica:} \textbf{\nomecursonovo} / FEELT / UFU}}\\\hline
    \hspace*{0pt}\ccheader{1º }{255}{60}{0}{0}{315} & \hspace*{0pt}\ccheader{2º }{285}{120}{0}{0}{405} & \hspace*{0pt}\ccheader{3º }{330}{75}{30}{0}{435} & \hspace*{0pt}\ccheader{4º }{345}{60}{0}{90}{495} & \hspace*{0pt}\ccheader{5º }{270}{45}{0}{90}{405} & \hspace*{0pt}\ccheader{6º }{210}{75}{0}{75}{360} & \hspace*{0pt}\ccheader{7º }{135}{75}{30}{60}{300} & \hspace*{0pt}\ccheader{8º }{120}{45}{60}{30}{255} & \hspace*{-4pt}\ccheader{Multi-}{0}{0}{0}{0}{480}\\
    \ccbloco{(01)}{Cálculo Diferencial e Integral I}{90}{0}{0}{0}{90}{}{} & \ccbloco{(07)}{Álgebra Linear}{45}{0}{0}{0}{45}{02}{} & \ccbloco{(15)}{Cálculo Diferencial e Integral III}{90}{0}{0}{0}{90}{08}{} & \ccbloco{(23)}{Estatística}{60}{0}{0}{0}{60}{}{} & \ccbloco{(31)}{Cálculo Numérico}{60}{0}{0}{0}{60}{07 15}{} & \ccbloco{(39)}{Aprendizagem de Máquina}{30}{15}{0}{0}{45}{31 36}{} & \ccbloco{(45)}{Aprendizagem Profunda e Modelos Generativos}{30}{15}{0}{0}{45}{39}{} & \ccbloco{(51)}{Engenharia Econômica}{30}{0}{0}{0}{30}{}{} & \ccbloco{(59)}{Disciplinas Optativas}{--}{--}{--}{--}{90}{}{} \\
        \ccbloco{(02)}{Geometria Analítica}{60}{0}{0}{0}{60}{}{} & \ccbloco{(08)}{Cálculo Diferencial e Integral II}{90}{0}{0}{0}{90}{01}{} & \ccbloco{(16)}{Física Básica: Eletricidade e Magnetismo}{60}{0}{0}{0}{60}{09}{} & \ccbloco{(24)}{Métodos Matemáticos}{75}{0}{0}{0}{75}{15}{} & \ccbloco{(32)}{Empreendedorismo e Inovação}{30}{0}{0}{0}{30}{}{} & \ccbloco{(40)}{Processamento Digital de Sinais}{45}{15}{0}{0}{60}{24}{} & \ccbloco{(46)}{Redes de Comunicações II}{45}{15}{0}{0}{60}{41}{} & \ccbloco{(52)}{Arquitetura de Software Aplicada}{15}{45}{0}{0}{60}{35 41}{} & \ccbloco{(57)}{Atividade de Conclusão de Curso}{--}{--}{--}{--}{300}{\rotatebox[origin=c]{90}{\tiny 1875 horas}}{} \\
        \ccbloco{(03)}{Introdução à Engenharia de Computação}{30}{0}{0}{0}{30}{}{} & \ccbloco{(09)}{Física Básica: Mecânica}{60}{0}{0}{0}{60}{01}{} & \ccbloco{(17)}{Física Básica: Eletricidade e Magnetismo, Experimental de}{0}{30}{0}{0}{30}{}{16} & \ccbloco{(25)}{Análise de Algoritmos}{60}{0}{0}{0}{60}{19}{} & \ccbloco{(33)}{Eletrônica Analógica I}{60}{0}{0}{0}{60}{28}{} & \ccbloco{(41)}{Redes de Comunicações I}{45}{15}{0}{0}{60}{37}{} & \ccbloco{(47)}{Segurança de Sistemas Computacionais}{15}{0}{30}{0}{45}{41}{} & \ccbloco{(53)}{Engenharia de Contexto em IA Generativa}{15}{0}{30}{0}{45}{45}{} & \ccbloco{(58)}{Atividades Acadêmicas Complementares}{--}{--}{--}{--}{90}{}{} \\
        \ccbloco{(04)}{Introdução à Prototipagem Eletrônica}{0}{30}{0}{0}{30}{}{} & \ccbloco{(10)}{Física Básica: Mecânica, Experimental de}{0}{30}{0}{0}{30}{}{09} & \ccbloco{(18)}{Elementos de Sistemas Computacionais}{30}{0}{30}{0}{60}{13}{} & \ccbloco{(26)}{Arquitetura e Organização de Computadores}{60}{0}{0}{0}{60}{18}{} & \ccbloco{(34)}{Eletrônica Analógica I, Experimental de}{0}{30}{0}{0}{30}{}{33} & \ccbloco{(42)}{Sistemas Embarcados I}{45}{30}{0}{0}{75}{33 37}{} & \ccbloco{(48)}{Sistemas Computacionais em Tempo Real}{15}{30}{0}{0}{45}{42}{} & \ccbloco{(54)}{IA Aplicada à Cibersegurança}{15}{0}{30}{0}{45}{45 47}{} &   \\
        \ccbloco{(05)}{Lógica e Matemática Discreta para Computação}{45}{0}{0}{0}{45}{}{} & \ccbloco{(11)}{Metrologia}{30}{30}{0}{0}{60}{}{} & \ccbloco{(19)}{Estrutura de Dados}{60}{0}{0}{0}{60}{12}{} & \ccbloco{(27)}{Banco de Dados}{15}{45}{0}{0}{60}{19}{} & \ccbloco{(35)}{Engenharia de Software}{45}{0}{0}{0}{45}{22 27}{} & \ccbloco{(43)}{Teoria da Computação}{45}{0}{0}{0}{45}{25}{} & \ccbloco{(49)}{Sistemas e Controle}{30}{15}{0}{0}{45}{40}{} & \ccbloco{(55)}{Sistemas Distribuídos}{45}{0}{0}{0}{45}{41}{} &   \\
        \ccbloco{(06)}{Programação Script}{30}{30}{0}{0}{60}{}{} & \ccbloco{(12)}{Programação Procedimental}{30}{30}{0}{0}{60}{06}{} & \ccbloco{(20)}{Inteligência Artificial, Ética e Sociedade}{30}{0}{0}{0}{30}{}{} & \ccbloco{(28)}{Circuitos Elétricos I}{75}{0}{0}{0}{75}{16}{} & \ccbloco{(36)}{Fundamentos da Inteligência Artificial}{30}{15}{0}{0}{45}{23}{} &   &   &   &   \\
          & \ccbloco{(13)}{Sistemas Digitais}{30}{0}{0}{0}{30}{05}{} & \ccbloco{(21)}{Programação Funcional}{30}{15}{0}{0}{45}{}{} & \ccbloco{(29)}{Circuitos Elétricos I, Experimental de}{0}{15}{0}{0}{15}{11}{28} & \ccbloco{(37)}{Sistemas Operacionais}{45}{0}{0}{0}{45}{12 26}{} &   &   &   &   \\
          & \ccbloco{(14)}{Sistemas Digitais, Experimental de}{0}{30}{0}{0}{30}{}{13} & \ccbloco{(22)}{Programação Orientada a Objetos}{30}{30}{0}{0}{60}{12}{} &   &   &   &   &   &   {\fontfamily{phv}\selectfont\footnotesize\textbf{CH Extensão:} 345 horas} \\
    \hline
      &   &   & \ccbloco{(30)}{Atividades Curriculares de Extensão I}{0}{0}{0}{90}{90}{}{} & \ccbloco{(38)}{Atividades Curriculares de Extensão II}{0}{0}{0}{90}{90}{30}{} & \ccbloco{(44)}{Atividades Curriculares de Extensão III}{0}{0}{0}{75}{75}{38}{} & \ccbloco{(50)}{Atividades Curriculares de Extensão IV}{0}{0}{0}{60}{60}{44}{} & \ccbloco{(56)}{Memorial de Atividades Curriculares de Extensão}{0}{0}{0}{30}{30}{50}{} & \framebox[1.1\width]{\textbf{CH Total:} 3450 horas}\\
        
        \hline
        \multicolumn{9}{l}{{\fontfamily{phv}\selectfont\footnotesize\textbf{Legenda:}\quad \rotatebox[origin=c]{90}{$\boldsymbol{\blacktriangledown}$} Pré-requisito \quad \rotatebox[origin=c]{-90}{$\boldsymbol{\bigstar}$} Correquisito \quad -- CHT/CHP/CHD/CHE indefinidas ou definidas conforme o componente curricular escolhido}}\\
            \hline\hline
        \end{tabular}
           
        }
     }
     \end{table}
     
    
        %             % \end{minipage}
        %             }
        
        %         % \end{center}
        %         }
        %      \end{table}
        %     % \end{landscape}       
        %%%%%%%%%%%%%%%%%%%%%%%%%%%%%%%%%%%%%%%%%%%%%%%%%%%%%%%%%%%%%%%%%%%%%%
\clearpage

\restoregeometry

\fancyhead[LE,RO]{\sffamily\nouppercase{\leftmark}}

% \hspace{-2cm}
%     \begin{table}
%         \centering
%         \caption*{{\scriptsize Representação Gráfica:} \textbf{\nomecursonovo}}
%         \noindent % Evita indentação
%         \hyphenpenalty=10000
%         \exhyphenpenalty=10000
%         \rotatebox{90}{ % Rotaciona apenas a tabela
%             \resizebox{0.8\pdfpageheight}{!}{ % Redimensiona a tabela para caber na página
%                 \begin{minipage}{\textheight} % Usa minipage para conter \input{}
%                     \newcommand{\ccheader}[6]{
    \begin{tikzpicture}
        \draw (-0.2,1.5) node[font=\itshape\sffamily\scriptsize, yshift=-7pt, anchor=center, align=right, text width=2ex]{};
        \draw (0,0.5) rectangle node[font=\bfseries\sffamily\scriptsize, text width=3cm, align=center]{#1Período} (3,1);
        
        \draw (0,0) rectangle node[font=\bfseries\sffamily\tiny, text width=0.75cm, align=center]{CHT} (0.6,0.5);
        \draw (0.6,0) rectangle node[font=\bfseries\sffamily\tiny, text width=0.6cm, align=center]{CHP} (1.2,0.5);
        \draw (1.2,0) rectangle node[font=\bfseries\sffamily\tiny, text width=0.6cm, align=center]{CHD} (1.8,0.5);
        \draw (1.8,0) rectangle node[font=\bfseries\sffamily\tiny, text width=0.6cm, align=center]{CHE} (2.4,0.5);
        \draw (2.4,0) rectangle node[font=\bfseries\sffamily\tiny, text width=0.6cm, align=center]{TOT} (3,0.5);
        
        \draw (0,-0.5) rectangle node[font=\bfseries\sffamily\scriptsize, text width=0.6cm, align=center]{#2} (0.6,0);
        \draw (0.6,-0.5) rectangle node[font=\bfseries\sffamily\scriptsize, text width=0.6cm, align=center]{#3} (1.2,0);
        \draw (1.2,-0.5) rectangle node[font=\bfseries\sffamily\scriptsize, text width=0.6cm, align=center]{#4} (1.8,0);
        \draw (1.8,-0.5) rectangle node[font=\bfseries\sffamily\scriptsize, text width=0.6cm, align=center]{#5} (2.4,0);
        \draw (2.4,-0.5) rectangle node[font=\bfseries\sffamily\scriptsize, text width=0.6cm, align=center]{#6} (3,0);
        
    \end{tikzpicture}
    }
    
    %\newcounter{contbloco}
    %\setcounter{contbloco}{0}

    \newcommand{\ccbloco}[9]{
    %\addtocounter{contbloco}{1}
    \begin{tikzpicture}
        \draw (0,0.5) rectangle node[font=\sffamily\scriptsize, yshift=-2pt, text width=2.8cm, align=center]{#2} (3,2.5);
        \draw (1.5,2.5) node[font=\bfseries\sffamily\scriptsize, yshift=-7pt, anchor=mid]{#1};
        \draw (0,0) rectangle node[font=\sffamily\scriptsize, text width=1cm, align=center]{#3} (0.6,0.5);
        \draw (0.6,0) rectangle node[font=\sffamily\scriptsize, text width=1cm, align=center]{#4} (1.2,0.5);
        \draw (1.2,0) rectangle node[font=\sffamily\scriptsize, text width=1cm, align=center]{#5} (1.8,0.5);
        \draw (1.8,0) rectangle node[font=\sffamily\scriptsize, text width=1cm, align=center]{#6} (2.4,0.5);
        \draw (2.4,0) rectangle node[font=\sffamily\scriptsize, text width=1cm, align=center]{#7} (3,0.5);
        \ifthenelse{\isempty{#8}}
        {
        \draw (-0.475,1.75) node[font=\sffamily\scriptsize, yshift=-7pt, anchor=center, align=right, text width=2ex]{};
        }
        {
        \draw (-0.475,1.75) node[font=\sffamily\scriptsize, yshift=-7pt, anchor=center, align=right, text width=2ex]{\textbf{#8}};
        \draw (-0.10,1.75) node[font=\sffamily\footnotesize, yshift=-7pt, anchor=mid, rotate=90]{$\boldsymbol{\blacktriangledown}$};
        }
        \ifthenelse{\isempty{#9}}
        {
        \draw (-0.475,0.5) node[font=\sffamily\scriptsize, yshift=-7pt, anchor=center, align=right, text width=2ex]{};
        }
        {
        \draw (-0.475,0.5) node[font=\sffamily\scriptsize, yshift=-7pt, anchor=center, align=right, text width=2ex]{\textbf{#9}};
        \draw (-0.30,0.5) node[font=\sffamily\footnotesize, yshift=-7pt, anchor=mid, rotate=90]{\rotatebox{180}{$\boldsymbol{\bigstar}$}};
        }
    \end{tikzpicture}
    }

    
    
    \begin{table}
    \refstepcounter{section}
    \addcontentsline{toc}{section}{\protect\numberline{\thesection}Representação Gráfica do Fluxo Curricular}
    \centering
    \rotatebox{90}{\resizebox{0.85\pdfpageheight}{!}{
    \nohyph{}
    \begin{tabular}{||cccccccc|c||}
    \hline\hline
    \multicolumn{9}{l}{{\fontfamily{phv}\selectfont{\scriptsize \framebox[1.1\width]{Seção~\thesection} PPC \ppcversao - Representação Gráfica:} \textbf{\nomecursonovo} / FEELT / UFU}}\\\hline
    \hspace*{0pt}\ccheader{1º }{255}{60}{0}{0}{315} & \hspace*{0pt}\ccheader{2º }{285}{120}{0}{0}{405} & \hspace*{0pt}\ccheader{3º }{330}{75}{30}{0}{435} & \hspace*{0pt}\ccheader{4º }{345}{60}{0}{90}{495} & \hspace*{0pt}\ccheader{5º }{270}{45}{0}{90}{405} & \hspace*{0pt}\ccheader{6º }{210}{75}{0}{75}{360} & \hspace*{0pt}\ccheader{7º }{135}{75}{30}{60}{300} & \hspace*{0pt}\ccheader{8º }{120}{45}{60}{30}{255} & \hspace*{-4pt}\ccheader{Multi-}{0}{0}{0}{0}{480}\\
    \ccbloco{(01)}{Cálculo Diferencial e Integral I}{90}{0}{0}{0}{90}{}{} & \ccbloco{(07)}{Álgebra Linear}{45}{0}{0}{0}{45}{02}{} & \ccbloco{(15)}{Cálculo Diferencial e Integral III}{90}{0}{0}{0}{90}{08}{} & \ccbloco{(23)}{Estatística}{60}{0}{0}{0}{60}{}{} & \ccbloco{(31)}{Cálculo Numérico}{60}{0}{0}{0}{60}{07 15}{} & \ccbloco{(39)}{Aprendizagem de Máquina}{30}{15}{0}{0}{45}{31 36}{} & \ccbloco{(45)}{Aprendizagem Profunda e Modelos Generativos}{30}{15}{0}{0}{45}{39}{} & \ccbloco{(51)}{Engenharia Econômica}{30}{0}{0}{0}{30}{}{} & \ccbloco{(59)}{Disciplinas Optativas}{--}{--}{--}{--}{90}{}{} \\
        \ccbloco{(02)}{Geometria Analítica}{60}{0}{0}{0}{60}{}{} & \ccbloco{(08)}{Cálculo Diferencial e Integral II}{90}{0}{0}{0}{90}{01}{} & \ccbloco{(16)}{Física Básica: Eletricidade e Magnetismo}{60}{0}{0}{0}{60}{09}{} & \ccbloco{(24)}{Métodos Matemáticos}{75}{0}{0}{0}{75}{15}{} & \ccbloco{(32)}{Empreendedorismo e Inovação}{30}{0}{0}{0}{30}{}{} & \ccbloco{(40)}{Processamento Digital de Sinais}{45}{15}{0}{0}{60}{24}{} & \ccbloco{(46)}{Redes de Comunicações II}{45}{15}{0}{0}{60}{41}{} & \ccbloco{(52)}{Arquitetura de Software Aplicada}{15}{45}{0}{0}{60}{35 41}{} & \ccbloco{(57)}{Atividade de Conclusão de Curso}{--}{--}{--}{--}{300}{\rotatebox[origin=c]{90}{\tiny 1875 horas}}{} \\
        \ccbloco{(03)}{Introdução à Engenharia de Computação}{30}{0}{0}{0}{30}{}{} & \ccbloco{(09)}{Física Básica: Mecânica}{60}{0}{0}{0}{60}{01}{} & \ccbloco{(17)}{Física Básica: Eletricidade e Magnetismo, Experimental de}{0}{30}{0}{0}{30}{}{16} & \ccbloco{(25)}{Análise de Algoritmos}{60}{0}{0}{0}{60}{19}{} & \ccbloco{(33)}{Eletrônica Analógica I}{60}{0}{0}{0}{60}{28}{} & \ccbloco{(41)}{Redes de Comunicações I}{45}{15}{0}{0}{60}{37}{} & \ccbloco{(47)}{Segurança de Sistemas Computacionais}{15}{0}{30}{0}{45}{41}{} & \ccbloco{(53)}{Engenharia de Contexto em IA Generativa}{15}{0}{30}{0}{45}{45}{} & \ccbloco{(58)}{Atividades Acadêmicas Complementares}{--}{--}{--}{--}{90}{}{} \\
        \ccbloco{(04)}{Introdução à Prototipagem Eletrônica}{0}{30}{0}{0}{30}{}{} & \ccbloco{(10)}{Física Básica: Mecânica, Experimental de}{0}{30}{0}{0}{30}{}{09} & \ccbloco{(18)}{Elementos de Sistemas Computacionais}{30}{0}{30}{0}{60}{13}{} & \ccbloco{(26)}{Arquitetura e Organização de Computadores}{60}{0}{0}{0}{60}{18}{} & \ccbloco{(34)}{Eletrônica Analógica I, Experimental de}{0}{30}{0}{0}{30}{}{33} & \ccbloco{(42)}{Sistemas Embarcados I}{45}{30}{0}{0}{75}{33 37}{} & \ccbloco{(48)}{Sistemas Computacionais em Tempo Real}{15}{30}{0}{0}{45}{42}{} & \ccbloco{(54)}{IA Aplicada à Cibersegurança}{15}{0}{30}{0}{45}{45 47}{} &   \\
        \ccbloco{(05)}{Lógica e Matemática Discreta para Computação}{45}{0}{0}{0}{45}{}{} & \ccbloco{(11)}{Metrologia}{30}{30}{0}{0}{60}{}{} & \ccbloco{(19)}{Estrutura de Dados}{60}{0}{0}{0}{60}{12}{} & \ccbloco{(27)}{Banco de Dados}{15}{45}{0}{0}{60}{19}{} & \ccbloco{(35)}{Engenharia de Software}{45}{0}{0}{0}{45}{22 27}{} & \ccbloco{(43)}{Teoria da Computação}{45}{0}{0}{0}{45}{25}{} & \ccbloco{(49)}{Sistemas e Controle}{30}{15}{0}{0}{45}{40}{} & \ccbloco{(55)}{Sistemas Distribuídos}{45}{0}{0}{0}{45}{41}{} &   \\
        \ccbloco{(06)}{Programação Script}{30}{30}{0}{0}{60}{}{} & \ccbloco{(12)}{Programação Procedimental}{30}{30}{0}{0}{60}{06}{} & \ccbloco{(20)}{Inteligência Artificial, Ética e Sociedade}{30}{0}{0}{0}{30}{}{} & \ccbloco{(28)}{Circuitos Elétricos I}{75}{0}{0}{0}{75}{16}{} & \ccbloco{(36)}{Fundamentos da Inteligência Artificial}{30}{15}{0}{0}{45}{23}{} &   &   &   &   \\
          & \ccbloco{(13)}{Sistemas Digitais}{30}{0}{0}{0}{30}{05}{} & \ccbloco{(21)}{Programação Funcional}{30}{15}{0}{0}{45}{}{} & \ccbloco{(29)}{Circuitos Elétricos I, Experimental de}{0}{15}{0}{0}{15}{11}{28} & \ccbloco{(37)}{Sistemas Operacionais}{45}{0}{0}{0}{45}{12 26}{} &   &   &   &   \\
          & \ccbloco{(14)}{Sistemas Digitais, Experimental de}{0}{30}{0}{0}{30}{}{13} & \ccbloco{(22)}{Programação Orientada a Objetos}{30}{30}{0}{0}{60}{12}{} &   &   &   &   &   &   {\fontfamily{phv}\selectfont\footnotesize\textbf{CH Extensão:} 345 horas} \\
    \hline
      &   &   & \ccbloco{(30)}{Atividades Curriculares de Extensão I}{0}{0}{0}{90}{90}{}{} & \ccbloco{(38)}{Atividades Curriculares de Extensão II}{0}{0}{0}{90}{90}{30}{} & \ccbloco{(44)}{Atividades Curriculares de Extensão III}{0}{0}{0}{75}{75}{38}{} & \ccbloco{(50)}{Atividades Curriculares de Extensão IV}{0}{0}{0}{60}{60}{44}{} & \ccbloco{(56)}{Memorial de Atividades Curriculares de Extensão}{0}{0}{0}{30}{30}{50}{} & \framebox[1.1\width]{\textbf{CH Total:} 3450 horas}\\
        
        \hline
        \multicolumn{9}{l}{{\fontfamily{phv}\selectfont\footnotesize\textbf{Legenda:}\quad \rotatebox[origin=c]{90}{$\boldsymbol{\blacktriangledown}$} Pré-requisito \quad \rotatebox[origin=c]{-90}{$\boldsymbol{\bigstar}$} Correquisito \quad -- CHT/CHP/CHD/CHE indefinidas ou definidas conforme o componente curricular escolhido}}\\
            \hline\hline
        \end{tabular}
           
        }
     }
     \end{table}
     
    
%                 \end{minipage}
%             }
%         }
%     \end{table}

%%%%%%%%%%%%%%%%%%%%%%%%%%%%%%%%%%%%%%%%%%%%%%%%%%%%%%%%%%%%%%%%%%%%%%

% \newpage
% % \tikzset{
%   mynode/.style={
%     draw,
%     fit={(0,0.5) (3,2)},
%     %xshift=0cm,
%     inner sep=0pt,
%     align=center,below,
%     execute at begin node={\baselineskip=3pt}
%   }
% }

% \tikzset{
%   mynode2/.style={
%     draw,
%     fit={(0,0) (1,0.5)},
%     %xshift=0cm,
%     %yshift=-1cm,
%     inner sep=0pt,
%     text width=1cm, 
%     anchor=center,below
%     %align=middle,
%   }
% }

% \begin{tikzpicture}
%     % \draw[black, thick] (0,0.5) rectangle (3,2) node[pos=0.5]{Intersection point};
%     \node[mynode] {\footnotesize Computação Gráfica e Realidade Virtual e Aumentada};
%     \node[mynode2] {\footnotesize 30};
%     % \draw[black, thick] (0,0) rectangle (1,0.5);
%     % \draw[black, thick] (1,0) rectangle (2,0.5);
%     % \draw[black, thick] (2,0) rectangle (3,0.5);
% \end{tikzpicture}

% \begin{tikzpicture}%text width=3cm,
%     \node [font=\footnotesize, fit={(0,0.5) (3,2)}, anchor=base, align=center, draw] {Computação Gráfica e Realidade Virtual e Aumentada};
%     \node [font=\footnotesize, fit={(0,0) (1,0.5)}, anchor=base, align=center, draw] {30};
%     % \node [font=\footnotesize, fit={(0,0) (3,1.5)},  align=center, draw] {Enriquecimento Instrumental};
% \end{tikzpicture}


\newcommand{\ccheader}[1]{
  \begin{tikzpicture}
    \draw (0,0.5) rectangle node[font=\bfseries\sffamily\scriptsize, text width=3cm, align=center]{#1º Período} (3,1);
    \draw (0,0) rectangle node[font=\bfseries\sffamily\scriptsize, text width=0.75cm, align=center]{CHT} (0.75,0.5);
    \draw (0.75,0) rectangle node[font=\bfseries\sffamily\scriptsize, text width=0.75cm, align=center]{CHP} (1.5,0.5);
    \draw (1.5,0) rectangle node[font=\bfseries\sffamily\scriptsize, text width=0.75cm, align=center]{CHR} (2.25,0.5);
    \draw (2.25,0) rectangle node[font=\bfseries\sffamily\scriptsize, text width=0.75cm, align=center]{TOT} (3,0.5);
  \end{tikzpicture}
}

% \begin{tikzpicture}
%   \draw (0,0.5) rectangle node[font=\sffamily\scriptsize, text width=3cm, align=center]{Computação Gráfica e Realidade Virtual e Aumentada\\(1)} (3,2.5);
%   \draw (0,0) rectangle node[font=\sffamily\scriptsize, text width=1cm, align=center]{30} (0.75,0.5);
%   \draw (0.75,0) rectangle node[font=\sffamily\scriptsize, text width=1cm, align=center]{30} (1.5,0.5);
%   \draw (1.5,0) rectangle node[font=\sffamily\scriptsize, text width=1cm, align=center]{0} (2.25,0.5);
%   \draw (2.25,0) rectangle node[font=\sffamily\scriptsize, text width=1cm, align=center]{60} (3,0.5);
% \end{tikzpicture}

% \newcounter{contbloco}
% \setcounter{contbloco}{0}

\newcommand{\ccbloco}[8]{
% \addtocounter{contbloco}{1}

% \ccbloco{(" + str(dscpl.enum) + r")}{" + dscpl.Nome + r"}{" + str(dscpl.CHT) + r"}{" + str(dscpl.CHP) + r"}{" + str(dscpl.CHR) + r"}{" + str(dscpl.TOT) + r"}
\begin{tikzpicture}
    \draw (0,0.5) rectangle node[font=\sffamily\scriptsize, text width=2.8cm, align=center]{#2} (3,2.5);
    \draw (1.5,2.5) node[font=\bfseries\sffamily\scriptsize, yshift=-7pt, anchor=mid]{#1};
    \draw (0,0) rectangle node[font=\sffamily\scriptsize, text width=1cm, align=center]{#3} (0.75,0.5);
    \draw (0.75,0) rectangle node[font=\sffamily\scriptsize, text width=1cm, align=center]{#4} (1.5,0.5);
    \draw (1.5,0) rectangle node[font=\sffamily\scriptsize, text width=1cm, align=center]{#5} (2.25,0.5);
    \draw (2.25,0) rectangle node[font=\sffamily\scriptsize, text width=1cm, align=center]{#6} (3,0.5);

    \ifthenelse{\isempty{#7}}
    {
    \draw (-0.475,1.5) node[font=\itshape\sffamily\scriptsize, yshift=-7pt, anchor=center, align=right, text width=2ex]{};
    }
    {
    \draw (-0.475,1.5) node[font=\itshape\sffamily\scriptsize, yshift=-7pt, anchor=center, align=right, text width=2ex]{#7};
    \draw (-0.125,1.5) node[font=\sffamily\scriptsize, yshift=-7pt, anchor=mid, rotate=90]{$\boldsymbol{\bigtriangledown}$};
    }
    % \draw (3.2,1.5) node[font=\sffamily\scriptsize, yshift=-7pt, anchor=mid, rotate=270]{#8};
    \draw (1.5,1) node[font=\sffamily\scriptsize, yshift=-7pt, anchor=mid]{+ ( #8 )};

\end{tikzpicture}
}

\begin{tabular}{c}
  \hspace*{1pt}\ccheader{2}\\
  \ccbloco{1}{Computação Gráfica e Realidade Virtual e Aumentada}{30}{30}{0}{60}{1 2 73 5 6 7 8 9}{5,6}\\
  \ccbloco{2}{Segurança de Sistemas Computacionais}{30}{0}{15}{45}{}{}\\
  \ccbloco{3}{Enriquecimento Instrumental}{0}{30}{0}{30}{}{}\\
\end{tabular}

% \newpage

% \clearpage

% \KOMAoptions{paper=a4,paper=portrait}
% \recalctypearea{}

\section{Relação de Disciplinas Optativas}
\label{sec:optativas}

No Quadro~\ref{tab:disciplinas_optativas}, apresenta-se a relação de disciplinas optativas pré-aprovadas para o curso de \nomecursonovo. Essas disciplinas são ofertadas de forma periódica, conforme a demanda, a disponibilidade de docentes ou a oferta por outros cursos e Unidades Acadêmicas. Como cada disciplina pode ou não possuir pré-requisitos, o(a) estudante deve verificar se já concluiu os componentes necessários ou se tem conhecimento suficiente para poder cursá-la. A carga horária mínima que o estudante deve integralizar por meio de disciplinas optativas é de \input{include/auto/par_carga_horaria_opt.tex}\unskip.

\begin{longtblr}[
        theme = ecp,
        caption = {Disciplinas Optativas Pré-Aprovadas},
        label = {tab:disciplinas_optativas},]{
    colspec = {Q[l,m,wd=63mm]Q[l,m,wd=23mm]Q[c,m,wd=7mm]Q[c,m,wd=7mm]Q[c,m,wd=7mm]Q[c,m,wd=7mm]Q[c,m,wd=9mm]},
    rowhead = 1,
    row{odd} = {bg=CinzaClaro},
    row{1} = {bg=AzulEscuro, fg=white},
    row{78} = {bg=AzulClaro, fg=white},
    cells  = {font=\fontsize{10pt}{12pt}\selectfont},
    }
        \textbf{Componente} & \textbf{Código}  & \textbf{CH Teór.} & \textbf{CH Prát.}  & \textbf{CH EaD} & \textbf{CH Ext.} & \textbf{CH Total} \\
        Bioinformática & FACOM39042 & 60 & 0 & 0 & 0 & 60 \\
    Ciência de Dados I & FACOM32602 & 30 & 30 & 0 & 0 & 60 \\
    Ciência de Dados II & FACOM32701 & 30 & 30 & 0 & 0 & 60 \\
    Ciências Sociais e Jurídicas & FADIR39901 & 60 & 0 & 0 & 0 & 60 \\
    Circuitos Elétricos II & FEELT31403 & 60 & 0 & 0 & 0 & 60 \\
    Experimental de Circuitos Elétricos II & FEELT31404 & 0 & 30 & 0 & 0 & 30 \\
    Cloud Computing Architecture AWS & FEELT!CCAWS & 0 & 0 & 45 & 0 & 45 \\
    Compiladores & FACOM39043 & 60 & 0 & 0 & 0 & 60 \\
    Computação Evolutiva & GBC210 & 60 & 0 & 0 & 0 & 60 \\
    Computação Evolutiva & FEELT!CEVO & 30 & 15 & 0 & 0 & 45 \\
    Computação Gráfica & GBC204 & 60 & 0 & 0 & 0 & 60 \\
    Computação Gráfica RV-RA & FEELT31721 & 30 & 15 & 0 & 0 & 45 \\
    Computação Quântica & FEELT!CQUA & 30 & 15 & 0 & 0 & 45 \\
    Construção de Compiladores & GBC071 & 60 & 0 & 0 & 0 & 60 \\
    Contratos Inteligentes e Segurança Blockchain & FEELT39040O & 30 & 15 & 0 & 0 & 45 \\
    Criptografia & FACOM39044 & 60 & 0 & 0 & 0 & 60 \\
    Data Warehouse & FACOM39049 & 60 & 0 & 0 & 0 & 60 \\
    Desenvolvimento de MOOCs & FEELT39041 & 30 & 15 & 0 & 0 & 45 \\
    Design Colaborativo & FEELT31623 & 30 & 15 & 0 & 0 & 45 \\
    Direito e Legislação & FADIR39801 & 30 & 0 & 0 & 0 & 30 \\
    Eletrônica Analógica II & FEELT32502 & 60 & 0 & 0 & 0 & 60 \\
    Experimental de Eletrônica Analógica II & FEELT31502 & 0 & 30 & 0 & 0 & 30 \\
    Explorações em Arquiteturas Computacionais [por tipo] & FEELT!EAC & 0 & 45 & 0 & 0 & 45 \\
    Explorações em Linguagens Computacionais [por tipo] & FEELT!ELC & 30 & 0 & 15 & 0 & 45 \\
    Física Básica: Oscilações, Ondas e Óptica & INFIS31402 & 60 & 0 & 0 & 0 & 60 \\
    Experimental de Física Básica: Oscilações, Ondas e Óptica & INFIS39056 & 0 & 30 & 0 & 0 & 30 \\
    Fontes Renováveis não Convencionais - Técnicas e Aplicações & FEELT39047 & 45 & 15 & 0 & 0 & 60 \\
    Gerenciamento de Banco de Dados & GBC053 & 60 & 0 & 0 & 0 & 60 \\
    Instalações Elétricas & FEELT31603 & 30 & 0 & 0 & 0 & 30 \\
    Inteligência Artificial Aplicada aos Negócios & GBC209 & 60 & 0 & 0 & 0 & 60 \\
    Inteligência Artificial Aplicada aos Negócios & GSI051 & 60 & 0 & 0 & 0 & 60 \\
    Inteligência Computacional & GBC073 & 60 & 0 & 0 & 0 & 60 \\
    Interação Humano-Computador & GSI037 & 60 & 0 & 0 & 0 & 60 \\
    Língua Brasileira de Sinais - Libras I & LIBRAS01 & 30 & 30 & 0 & 0 & 60 \\
    Língua Brasileira de Sinais - Libras II & LIBRAS02 & 30 & 30 & 0 & 0 & 60 \\
    Micro e Nanoeletrônica & FEELT39043 & 30 & 15 & 0 & 0 & 45 \\
    Mineração de Dados & GSI055 & 60 & 0 & 0 & 0 & 60 \\
    Mineração de Dados & GBC212 & 60 & 0 & 0 & 0 & 60 \\
    Mineração de Dados & GBC212 & 60 & 0 & 0 & 0 & 60 \\
    Modelagem e Simulação & GBC065 & 60 & 0 & 0 & 0 & 60 \\
    Modelagem e Simulação de Sistemas a Eventos Discretos & FEELT32901 & 45 & 15 & 0 & 0 & 60 \\
    Modelos de Linguagem e Agentes de IA & FEELT!MLIA & 30 & 15 & 0 & 0 & 45 \\
    Multimídia & GBC213 & 60 & 0 & 0 & 0 & 60 \\
    Organização e Recuperação da Informação & FACOM32605 & 30 & 30 & 0 & 0 & 60 \\
    Otimização & FACOM31601 & 60 & 0 & 0 & 0 & 60 \\
    Otimização & GSI027 & 60 & 0 & 0 & 0 & 60 \\
    Otimização e Simulação & FEELT31723 & 30 & 15 & 0 & 0 & 45 \\
    Pesquisa Operacional & FAGEN31703 & 30 & 30 & 0 & 0 & 60 \\
    Princípios e Padrões de Projeto & FACOM31401 & 30 & 30 & 0 & 0 & 60 \\
    Processamento Digital de Imagens & FACOM39053 & 60 & 0 & 0 & 0 & 60 \\
    Processamento Digital de Imagens & GSI058 & 60 & 0 & 0 & 0 & 60 \\
    Programação Lógica & GBC025 & 30 & 30 & 0 & 0 & 60 \\
    Programação Lógica e Inteligência Artificial & FEELT31407 & 30 & 15 & 0 & 0 & 45 \\
    Programação Paralela e Distribuída & GSI059 & 60 & 0 & 0 & 0 & 60 \\
    Programação para Dispositivos Móveis & FACOM32503 & 30 & 30 & 0 & 0 & 60 \\
    Programação para Internet & GBC084 & 30 & 30 & 0 & 0 & 60 \\
    Projeto de Redes de Computadores & GBC219 & 45 & 15 & 0 & 0 & 60 \\
    Projetos Digitais com FPGA & FEELT!FPGA & 0 & 45 & 0 & 0 & 45 \\
    Projetos de Hardware e Firmware para IOT & FEELT!PHFI & 0 & 45 & 0 & 0 & 45 \\
    Redes de Sensores e Internet das Coisas & FEELT39057 & 45 & 15 & 0 & 0 & 60 \\
    Resolução de Problemas & FEELT39042 & 30 & 15 & 0 & 0 & 45 \\
    Resolução de Problemas I & FEELT!RP1 & 0 & 45 & 0 & 0 & 45 \\
    Resolução de Problemas II & FEELT!RP2 & 0 & 45 & 0 & 0 & 45 \\
    Robótica & FEELT31722 & 45 & 15 & 0 & 0 & 60 \\
    Sinais e Multimídia & FEELT31526 & 30 & 15 & 0 & 0 & 45 \\
    Sistemas Embarcados II & FEELT39015 & 30 & 30 & 0 & 0 & 60 \\
    Tecnologias Web e Mobile & FEELT31521 & 30 & 15 & 0 & 0 & 45 \\
    Tópicos Especiais em Engenharia de Computação: Bancos de Dados Não Relacionais & FEELT39040E & 30 & 15 & 0 & 0 & 45 \\
    Tópicos Especiais em Engenharia de Computação: Cloud Computing Architecture AWS & FEELT39040C & 30 & 15 & 0 & 0 & 45 \\
    Tópicos Especiais em Engenharia de Computação: Computação Evolutiva & FEELT39040G & 30 & 15 & 0 & 0 & 45 \\
    Tópicos Especiais em Engenharia de Computação: Desenvolvimento de Jogos com uso de Realidade Virtual e Aumentada & FEELT39040F & 30 & 15 & 0 & 0 & 45 \\
    Tópicos Especiais em Engenharia de Computação: Desenvolvimento de hardware microcontrolado e firmware para IoT & FEELT39040P & 30 & 15 & 0 & 0 & 45 \\
    Tópicos em Engenharia de Computação TEC30 [por tipo] & FEELT!TEC30 & -- & -- & -- & -- & 30 \\
    Tópicos em Engenharia de Computação TEC45 [por tipo] & FEELT!TEC45 & -- & -- & -- & -- & 45 \\
    Tópicos em Engenharia de Computação TEC60 [por tipo] & FEELT!TEC60 & -- & -- & -- & -- & 60 \\
    Tópicos em Engenharia de Computação TEC75 [por tipo] & FEELT!TEC75 & -- & -- & -- & -- & 75 \\*

    \end{longtblr}
    

Qualquer outra disciplina cursada em qualquer unidade acadêmica da UFU poderá ser incorporada a essa relação\footnote{\url{https://www.feelt.ufu.br/graduacao/engenharia-de-computacao/saiba-mais/optativas-aprovadas}}, mediante requerimento do interessado. Após a anuência do Núcleo Docente Estruturante (NDE) e a aprovação do Colegiado do Curso, a solicitação seguirá os trâmites institucionais necessários para efetivação da incorporação.


\section{Relação de Disciplinas Equivalentes}
\label{sec:equivalentes}

No Quadro~\ref{tab:disciplinas_equivalentes}, apresenta-se a relação de disciplinas equivalentes pré-aprovadas para o curso de \nomecursonovo. Essas disciplinas também são ofertadas de forma periódica, conforme a demanda, a disponibilidade de docentes ou a oferta por outros cursos e Unidades Acadêmicas.

Entende-se por disciplinas equivalentes aquelas que compartilham conteúdos programáticos similares ou que, embora adotem abordagens distintas, promovem impactos formativos equivalentes na trajetória do estudante. A adoção de equivalências visa respeitar a dinamicidade do cenário tecnológico, reconhecendo que determinados conhecimentos podem ser transmitidos por diferentes caminhos pedagógicos, desde que se preservem os objetivos educacionais essenciais do curso. Essa abordagem confere flexibilidade curricular responsável, sem comprometer a coerência e a qualidade da formação em \nomecursonovo.

\begin{longtblr}[
        theme = ecp,
        caption = {Disciplinas Equivalentes Pré-Aprovadas},
        label = {tab:disciplinas_equivalentes},]{
    colspec = {Q[l,m,wd=31mm]Q[l,m,wd=19mm]Q[c,m,wd=7mm]Q[c,m,wd=7mm]Q[c,m,wd=7mm]Q[c,m,wd=7mm]Q[c,m,wd=9mm]Q[l,m,wd=31mm]},
    rowhead = 1,
    row{odd} = {bg=CinzaClaro},
    row{1} = {bg=AzulEscuro, fg=white},
    row{24} = {bg=AzulClaro, fg=white},
    cells  = {font=\fontsize{9pt}{9pt}\selectfont},
    }
        \textbf{Componente} & \textbf{Código}  & \textbf{CH Teór.} & \textbf{CH Prát.}  & \textbf{CH EaD} & \textbf{CH Ext.} & \textbf{CH Total}  & \textbf{Equivale a} \\
        Administração & FAGEN39901 & 60 & 0 & 0 & 0 & 60 & FAGEN39401 Empreendorismo e Inovação \\
    Algoritmos e Estrutura de Dados I & FACOM31201 & 45 & 15 & 0 & 0 & 60 & FEELT!ED Estrutura de Dados \\
    Algoritmos e Estruturas de Dados II & FACOM31303 & 60 & 0 & 0 & 0 & 60 & FEELT!AA Análise de Algoritmos \\
    Arquitetura de Redes TCP/IP & GBC066 & 30 & 30 & 0 & 0 & 60 & FEELT31831 Redes de Comunicações II \\
    Arquitetura de Redes de Computadores & GBC056 & 60 & 0 & 0 & 0 & 60 & FEELT31724 Redes de Comunicações I \\
    Ciências Econômicas & IEUFU39901 & 60 & 0 & 0 & 0 & 60 & IERI39601 Engenharia Econômica \\
    Ciências Econômicas & IERI39001 & 60 & 0 & 0 & 0 & 60 & IERI39601 Engenharia Econômica \\
    Gerenciamento de Projetos & FAGEN32502 & 30 & 30 & 0 & 0 & 60 & FAGEN39401 Empreendorismo e Inovação \\
    Inteligência Artificial & GBC063 & 60 & 0 & 0 & 0 & 60 & FEELT!FIA Fundamentos da Inteligência Artificial \\
    Inteligência Artificial & FACOM39050 & 60 & 0 & 0 & 0 & 60 & FEELT!FIA Fundamentos da Inteligência Artificial \\
    Linguagens Formais e Autômatos & GBC044 & 60 & 0 & 0 & 0 & 60 & FEELT!TCOMP Teoria da Computação \\
    Perspectivas em Engenharia de Computação PEC30 [por tipo] & FEELT!PEC30 & -- & -- & -- & -- & 30 & [por tipo] \\
    Perspectivas em Engenharia de Computação PEC45 [por tipo] & FEELT!PEC45 & -- & -- & -- & -- & 45 & [por tipo] \\
    Perspectivas em Engenharia de Computação PEC60 [por tipo] & FEELT!PEC60 & -- & -- & -- & -- & 60 & [por tipo] \\
    Perspectivas em Engenharia de Computação PEC75 [por tipo] & FEELT!PEC75 & -- & -- & -- & -- & 75 & [por tipo] \\
    Programação Funcional & GBC033 & 30 & 30 & 0 & 0 & 60 & FEELT!PF Programação Funcional \\
    Redes Industriais para Controle e Automação I & FEELT31713 & 60 & 15 & 0 & 0 & 75 & FEELT31724 Redes de Comunicações I \\
    Redes Industriais para Controle e Automação II & FEELT32801 & 30 & 60 & 0 & 0 & 90 & FEELT31831 Redes de Comunicações II \\
    Sistemas Distribuídos & FACOM32606 & 60 & 0 & 0 & 0 & 60 & FEELT!SD Sistemas Distribuídos \\
    Sistemas Distribuídos & GBC074 & 60 & 0 & 0 & 0 & 60 & FEELT!SD Sistemas Distribuídos \\
    Sistemas Operacionais & GBC045 & 60 & 30 & 0 & 0 & 90 & FEELT!SO Sistemas Operacionais \\
    Sistemas de Banco de Dados & GBC043 & 60 & 30 & 0 & 0 & 90 & FEELT!BD Banco de Dados \\*

    \end{longtblr}
    

Outras disciplinas poderão ser incorporadas a essa relação\footnote{\url{https://www.feelt.ufu.br/graduacao/engenharia-de-computacao/saiba-mais/equivalencias-aprovadas}}, mediante requerimento do interessado. Após a anuência do Núcleo Docente Estruturante (NDE) e a aprovação do Colegiado do Curso, a solicitação seguirá os trâmites institucionais necessários para efetivação da incorporação.

\section{Conteúdos de Destaque e/ou Transversais}

O curso de \nomecursonovo da Universidade Federal de Uberlândia contempla, ao longo de sua estrutura curricular, diversos conteúdos de destaque e natureza transversal, alinhados às Diretrizes Curriculares Nacionais e às demandas contemporâneas da formação profissional. Esses conteúdos são abordados de maneira integrada, transversal ou interdisciplinar, permeando múltiplas unidades curriculares e projetos formativos. Dentre os principais, destacam-se:

\begin{itemize}
\item \textbf{Ética e Responsabilidade Profissional}: Discussões sobre responsabilidade técnica, segurança, privacidade, uso ético de dados e impactos sociais da computação e automação estão presentes em componentes como Introdução à Engenharia de Computação, Projeto de Sistemas Computacionais e nas Atividades Curriculares de Extensão.

\item \textbf{Sustentabilidade e Impacto Socioambiental}: O curso promove a reflexão crítica sobre os impactos ambientais e sociais das tecnologias desenvolvidas, com ênfase em soluções energeticamente eficientes, reaproveitamento de materiais e uso consciente de recursos computacionais.

\item \textbf{Inovação, Empreendedorismo e Propriedade Intelectual}: Tais temas são tratados em disciplinas como Empreendedorismo e Inovação Tecnológica, Projeto Integrador e Trabalho de Conclusão de Curso, além de serem estimulados em atividades extensionistas e nos projetos com empresas juniores e incubadoras.

\item \textbf{Diversidade, Inclusão e Direitos Humanos}: A formação propicia espaços para o debate e valorização da diversidade cultural, de gênero, étnico-racial e de acessibilidade, tanto em conteúdos da formação geral quanto em contextos de aplicação tecnológica.

\item \textbf{Extensão Universitária}: A curricularização da extensão assegura que os(as) estudantes interajam com demandas reais da sociedade, contribuindo para a formação cidadã e o fortalecimento da função social da universidade.

\item \textbf{Trabalho Colaborativo, Comunicação e Liderança}: Estimulados por metodologias ativas (PBL, projetos integradores) e práticas em equipe, esses conteúdos desenvolvem competências interpessoais fundamentais para o exercício da engenharia em ambientes multidisciplinares.

\item \textbf{Interdisciplinaridade e Integração Ciência-Tecnologia-Sociedade (CTS)}: Presente em atividades práticas, integradoras e de pesquisa, esse eixo fortalece a capacidade crítica e a atuação ética e técnica diante de problemas complexos e contextos sociotécnicos.

\end{itemize}

% \newpage
% \section{Certificados}
% \todo{Completar}

% Nesta Seção do projeto, apresentam-se as certificações oferecidas no curso, um caminho significativo para os estudantes aprofundarem suas habilidades em áreas específicas da Engenharia de Computação. Cada certificação representa um conjunto estruturado de disciplinas optativas pré-selecionadas, destinadas a fornecer uma base sólida e especializada na área correspondente. Para alcançar uma dessas certificações reconhecidas, os alunos deverão completar um conjunto específico de disciplinas optativas, pré-determinadas e alinhadas com os requisitos e competências essenciais daquela certificação em particular, para cada certificação que desejar.

% Cada certificação é projetada para oferecer uma experiência acadêmica enriquecedora e prática, capacitando os estudantes a adquirirem conhecimentos especializados e habilidades direcionadas. Ao selecionar disciplinas entre as opções disponíveis no curso, os alunos têm a oportunidade de personalizar seu percurso educacional, adaptando-o aos seus interesses e aspirações profissionais. Essas certificações não apenas enriquecem o conhecimento técnico dos alunos, mas também os preparam para desafios específicos e demandas crescentes dentro do campo da Engenharia de Computação. A seguir, apresentamos as listas de disciplinas optativas associadas a cada certificação, fornecendo um roteiro claro e detalhado para os alunos que desejam buscar uma especialização particular durante seus estudos.

% O Colegiado do Curso poderá incluir, remover ou modificar as disciplinas que compõem os requisitos para qualquer uma das certificações existentes, bem como estabelecer ou encerrar certificações, sempre que forem justificadas a necessidade de tais ações. Quando devidas, no entanto, essas mudanças serão feitas de modo a não prejudicar os estudantes que já estejam avançando em busca de uma determinada certificação. O objetivo é assegurar que qualquer ajuste no conjunto de disciplinas necessárias para as certificações seja feito de maneira a beneficiar os alunos e aprimorar a qualidade do programa, sem causar impactos negativos à trajetória acadêmica daqueles que estejam perseguindo esses objetivos educacionais específicos.

% \todoreview{sugestão para o curso sem certificado}

% \subsection{Certificado em Engenharia de Algoritmos}

% Engenharia de Algoritmos se refere à especialização na compreensão, projeto e implementação eficiente de algoritmos em diversos campos da computação. Esse certificado oferece um conjunto selecionado de disciplinas optativas que aprofundam o conhecimento em estruturas algorítmicas, otimização de algoritmos e análise de complexidade. A Engenharia de Algoritmos capacita os estudantes a enfrentar demandas crescentes na criação, adaptação e implementação de algoritmos eficientes, preparando-os para uma ampla gama de oportunidades em áreas técnicas e de pesquisa.

% \renewcommand{\familydefault}{\sfdefault} \normalfont
% \begin{longtblr}[
%     theme = ecp,
%     caption = {Engenharia de Algoritmos},
%     label = {tab:engenharia_algoritmos},
% ]{
%   colspec = {Q[l,m,wd=65mm]Q[c,m,wd=24mm]Q[c,m,wd=8mm]Q[c,m,wd=8mm]Q[c,m,wd=8mm]Q[c,m,wd=10mm]},
%   rowhead = 1,
%   row{odd} = {bg=CinzaClaro},
%   row{1} = {bg=AzulEscuro, fg=white},
%   row{2} = {bg=AzulClaro, fg=white},
%   row{6} = {bg=AzulClaro, fg=white},
%   row{11} = {bg=AzulClaro, fg=white},
%   cells  = {font=\fontsize{10pt}{12pt}\selectfont},
% } 
%     % Componente & Código & CH Teór. & CH Prát.  & CH Rem. & CH Total \\
%     \textbf{Componente} & \textbf{Código}  & \textbf{CH Teór.} & \textbf{CH Prát.}  & \textbf{CH Rem.} & \textbf{CH Total} \\
%     \ \ \ \textit{Optativas A} &&&&&\\
%     Resolução de Problemas I & FEELT? &  &  & 0 & 45 \\
%     Teoria dos Grafos & GBC042 & & & 0 & 60 \\
%     Otimização & FACOM31601 & & & 0 & 60 \\
%     \ \ \ \textit{Optativas B} &&&&&\\
%     Resolução de Problemas II & FEELT? &  &  & 0 & 45 \\
%     Algoritmos Avançados & FEELT? &  &  & 0 & 45 \\
%     Construção de Compiladores & GBC071 & & & 0 & 60 \\
%     Programação Paralela e Distribuída & GBC218 & & & 0 & 60 \\
%     % Outras optativas à escolha & -- & & & & 120 \\
%     \textbf{TOTAL} &  & \textbf{X} & \textbf{X} & \textbf{0} & \textbf{375} \\
% \end{longtblr}
% \renewcommand{\familydefault}{\rmdefault} \normalfont

% \subsection{Certificado em Engenharia de Algoritmos}

Engenharia de Algoritmos se refere à especialização na compreensão, projeto e implementação eficiente de algoritmos em diversos campos da computação. Esse certificado oferece um conjunto selecionado de disciplinas optativas que aprofundam o conhecimento em estruturas algorítmicas, otimização de algoritmos e análise de complexidade. A Engenharia de Algoritmos capacita os estudantes a enfrentar demandas crescentes na criação, adaptação e implementação de algoritmos eficientes, preparando-os para uma ampla gama de oportunidades em áreas técnicas e de pesquisa.

% \renewcommand{\familydefault}{\sfdefault} \normalfont
% \begin{longtblr}[
%     theme = ecp,
%     caption = {Engenharia de Algoritmos},
%     label = {tab:engenharia_algoritmos},
% ]{
%   colspec = {Q[l,m,wd=65mm]Q[c,m,wd=24mm]Q[c,m,wd=8mm]Q[c,m,wd=8mm]Q[c,m,wd=8mm]Q[c,m,wd=10mm]},
%   rowhead = 1,
%   row{odd} = {bg=CinzaClaro},
%   row{1} = {bg=AzulEscuro, fg=white},
%   row{2} = {bg=AzulClaro, fg=white},
%   row{6} = {bg=AzulClaro, fg=white},
%   row{11} = {bg=AzulClaro, fg=white},
%   cells  = {font=\fontsize{10pt}{12pt}\selectfont},
% } 
%     % Componente & Código & CH Teór. & CH Prát.  & CH Rem. & CH Total \\
%     \textbf{Componente} & \textbf{Código}  & \textbf{CH Teór.} & \textbf{CH Prát.}  & \textbf{CH Rem.} & \textbf{CH Total} \\
%     \ \ \ \textit{Optativas A} &&&&&\\
%     Resolução de Problemas I & FEELT? &  &  & 0 & 45 \\
%     Teoria dos Grafos & GBC042 & & & 0 & 60 \\
%     Otimização & FACOM31601 & & & 0 & 60 \\
%     \ \ \ \textit{Optativas B} &&&&&\\
%     Resolução de Problemas II & FEELT? &  &  & 0 & 45 \\
%     Algoritmos Avançados & FEELT? &  &  & 0 & 45 \\
%     Construção de Compiladores & GBC071 & & & 0 & 60 \\
%     Programação Paralela e Distribuída & GBC218 & & & 0 & 60 \\
%     % Outras optativas à escolha & -- & & & & 120 \\
%     \textbf{TOTAL} &  & \textbf{X} & \textbf{X} & \textbf{0} & \textbf{375} \\
% \end{longtblr}
% \renewcommand{\familydefault}{\rmdefault} \normalfont

\renewcommand{\familydefault}{\sfdefault} \normalfont{}
    \begin{longtblr}[
        theme = ecp,
        caption = {Engenharia de Algoritmos},
        label = {tab:engenharia_algoritmos},
    ]{
    colspec = {Q[l,m,wd=65mm]Q[c,m,wd=24mm]Q[c,m,wd=8mm]Q[c,m,wd=8mm]Q[c,m,wd=8mm]Q[c,m,wd=10mm]},
    rowhead = 1,
    row{odd} = {bg=CinzaClaro},
    row{1} = {bg=AzulEscuro, fg=white},
    row{9} = {bg=AzulClaro, fg=white},
    cells  = {font=\fontsize{10pt}{12pt}\selectfont},
    }
        \textbf{Componente} & \textbf{Código}  & \textbf{CH Teór.} & \textbf{CH Prát.}  & \textbf{CH Rem.} & \textbf{CH Total} \\
        Algoritmos Avançados & FEELT:AAV & 0 & 0 & 0 & 45 \\
    Construção de Compiladores & GBC071 & 60 & 0 & 0 & 60 \\
    Otimização & FACOM31601 & 60 & 0 & 0 & 60 \\
    Programação Paralela e Distribuída & GBC218 & 60 & 0 & 0 & 60 \\
    Resolução de Problemas I & FEELT:RP1 & 0 & 0 & 0 & 45 \\
    Resolução de Problemas II & FEELT:RP2 & 0 & 0 & 0 & 45 \\
    Teoria dos Grafos & FEELT:TGRAF & 0 & 0 & 0 & 45 \\
    \textbf{TOTAL} &  & \textbf{180} & \textbf{0} & \textbf{0} & \textbf{360} \\
    \end{longtblr}
    \renewcommand{\familydefault}{\rmdefault} \normalfont{}
    
% \subsection{Certificado em Engenharia de Hardware}

A Engenharia de Hardware refere-se a uma especialização focada no projeto, desenvolvimento e implementação de sistemas eletrônicos e dispositivos físicos na área da Engenharia de Computação. Esse certificado oferece um conjunto específico de disciplinas optativas que aprofundam o conhecimento em eletrônica, design de circuitos, arquitetura de computadores, sistemas embarcados e tecnologias de hardware, capacitando os estudantes projetar e otimizar componentes de hardware, compreendendo desde a eletrônica básica até a integração avançada de sistemas. A Engenharia de Hardware prepara os estudantes para enfrentar desafios complexos no desenvolvimento de tecnologias eletrônicas e contribuir para a inovação em setores como computação embarcada, Internet das Coisas (IoT), sistemas integrados e novas fronteiras da computação.

% \renewcommand{\familydefault}{\sfdefault} \normalfont
% \begin{longtblr}[
%     theme = ecp,
%     caption = {Engenharia de Hardware},
%     label = {tab:engenharia_hardware},
% ]{
%   colspec = {Q[l,m,wd=65mm]Q[c,m,wd=24mm]Q[c,m,wd=8mm]Q[c,m,wd=8mm]Q[c,m,wd=8mm]Q[c,m,wd=10mm]},
%   rowhead = 1,
%   row{odd} = {bg=CinzaClaro},
%   row{1} = {bg=AzulEscuro, fg=white},
%   row{2} = {bg=AzulClaro, fg=white},
%   row{5} = {bg=AzulClaro, fg=white},
%   row{10} = {bg=AzulClaro, fg=white},
%   cells  = {font=\fontsize{10pt}{12pt}\selectfont},
% } 
%     % Componente & Código & CH Teór. & CH Prát.  & CH Rem. & CH Total \\
%     \textbf{Componente} & \textbf{Código}  & \textbf{CH Teór.} & \textbf{CH Prát.}  & \textbf{CH Rem.} & \textbf{CH Total} \\
%     \ \ \ \textit{Optativas A} &&&&&\\
%     Física Básica: Oscilações, Ondas e Óptica & INFIS31402 & 60 & 0 & 0 & 60 \\
%     Instrumentação Industrial I & FEELT32503 & 45 & 15 & 0 & 60 \\
%     \ \ \ \textit{Optativas B} &&&&&\\
%     Projetos Eletrônicos e Microprocessados & FEELT? & 0 & 45 & 0 & 45\\%https://blog.raisa.com.br/construcao-e-fabricacao-de-eletronicos-guia-basico/
%     % Robótica & FEELT31722 & 45 & 15 & 0 & 60 \\
%     Sistemas de Controle Realimentado & FEELT32603 & 60 & 0 & 0 & 60 \\
%     Sistemas Computacionais em Tempo Real & FEELT31725 & 30 & 15 & 0 & 45 \\
%     Sistemas Embarcados II & FEELT39015 & 30 & 30 & 0 & 60 \\
%     % Outras optativas à escolha & -- & & & & 120 \\
%     \textbf{TOTAL} &  & \textbf{225} & \textbf{105} & \textbf{0} & \textbf{330} \\
% \end{longtblr}
% \renewcommand{\familydefault}{\rmdefault} \normalfont

\renewcommand{\familydefault}{\sfdefault} \normalfont{}
    \begin{longtblr}[
        theme = ecp,
        caption = {Engenharia de Hardware},
        label = {tab:engenharia_hardware},
    ]{
    colspec = {Q[l,m,wd=65mm]Q[c,m,wd=24mm]Q[c,m,wd=8mm]Q[c,m,wd=8mm]Q[c,m,wd=8mm]Q[c,m,wd=10mm]},
    rowhead = 1,
    row{odd} = {bg=CinzaClaro},
    row{1} = {bg=AzulEscuro, fg=white},
    row{8} = {bg=AzulClaro, fg=white},
    cells  = {font=\fontsize{10pt}{12pt}\selectfont},
    }
        \textbf{Componente} & \textbf{Código}  & \textbf{CH Teór.} & \textbf{CH Prát.}  & \textbf{CH Rem.} & \textbf{CH Total} \\
        Física Básica: Oscilações Ondas e Óptica & INFIS31402 & 60 & 0 & 0 & 60 \\
    Instrumentação Industrial I & FEELT32503 & 45 & 15 & 0 & 60 \\
    Projetos Eletrônicos e Microprocessados & FEELT:PEM & 0 & 45 & 0 & 45 \\
    Sistemas Computacionais em Tempo Real & FEELT31725 & 30 & 15 & 0 & 45 \\
    Sistemas Embarcados II & FEELT39015 & 30 & 30 & 0 & 60 \\
    Sistemas de Controle Realimentado & FEELT32603 & 60 & 0 & 0 & 60 \\
    \textbf{TOTAL} &  & \textbf{225} & \textbf{105} & \textbf{0} & \textbf{330} \\
    \end{longtblr}
    \renewcommand{\familydefault}{\rmdefault} \normalfont{}
    
% \subsection{Certificado em Computação Gráfica e Realidade\ \mbox{Virtual} e Aumentada}

Um certificado em Computação Gráfica e Realidade Virtual e Aumentada oferece uma especialização focada na exploração e aplicação de tecnologias avançadas para criar ambientes virtuais e experiências imersivas. Esse certificado abrange disciplinas que aprofundam o entendimento dos alunos em áreas como gráficos computacionais, modelagem 3D, renderização, interação humano-computador e conceitos fundamentais de realidade virtual e aumentada. Esta linha de optativas prepara os alunos para trabalhar na vanguarda da tecnologia, capacitando-os a contribuir para a criação de conteúdo imersivo, jogos, simulações, treinamentos virtuais, aplicações de design e muito mais, em indústrias como entretenimento, educação, saúde, arquitetura e engenharia. 

% \renewcommand{\familydefault}{\sfdefault} \normalfont
% \begin{longtblr}[
%     theme = ecp,
%     caption = {Computação Gráfica e Realidade Virtual e Aumentada},
%     label = {tab:cg_rva},
% ]{
%   colspec = {Q[l,m,wd=65mm]Q[c,m,wd=24mm]Q[c,m,wd=8mm]Q[c,m,wd=8mm]Q[c,m,wd=8mm]Q[c,m,wd=10mm]},
%   rowhead = 1,
%   row{odd} = {bg=CinzaClaro},
%   row{1} = {bg=AzulEscuro, fg=white},
%   row{2} = {bg=AzulClaro, fg=white},
%   row{5} = {bg=AzulClaro, fg=white},
%   row{12} = {bg=AzulClaro, fg=white},
%   cells  = {font=\fontsize{10pt}{12pt}\selectfont},
% } 
%     % Componente & Código & CH Teór. & CH Prát.  & CH Rem. & CH Total \\
%     \textbf{Componente} & \textbf{Código}  & \textbf{CH Teór.} & \textbf{CH Prát.}  & \textbf{CH Rem.} & \textbf{CH Total} \\
%     \ \ \ \textit{Optativas A} &&&&&\\
%     Computação Gráfica RV-RA & FEELT31721 & 30 & 15 & 0 & 45 \\
%     Tecnologias Web & FEELT? &  &  & 0 & 45\\
%     \ \ \ \textit{Optativas B} &&&&&\\
%     Arquitetura de Computação em Nuvem AWS & FEELT? & 0 & 0 & 45 & 45 \\
%     RV-RA Jogos Sérios & FEELT? &  &  & 0 & 45\\
%     Tecnologias Mobile & FEELT?  &  &  & 0 & 45 \\
%     Processamento Digital de Imagens & FEELT39022 ou GBC216 & 60 & 0 & 0 & 60\\
%     Visão Computacional & FEELT?  &  &  & 0 & 45 \\
%     Interação Humano-Computador & FACOM32704 & 60 & 0 & 0 & 60 \\
%     \textbf{TOTAL} &  & \textbf{225} & \textbf{105} & \textbf{0} & \textbf{390} \\
% \end{longtblr}
% \renewcommand{\familydefault}{\rmdefault} \normalfont

\renewcommand{\familydefault}{\sfdefault} \normalfont{}
    \begin{longtblr}[
        theme = ecp,
        caption = {Computação Gráfica e Realidade Virtual e Aumentada},
        label = {tab:computacao_grafica_realidade_virtual_aumentada},
    ]{
    colspec = {Q[l,m,wd=65mm]Q[c,m,wd=24mm]Q[c,m,wd=8mm]Q[c,m,wd=8mm]Q[c,m,wd=8mm]Q[c,m,wd=10mm]},
    rowhead = 1,
    row{odd} = {bg=CinzaClaro},
    row{1} = {bg=AzulEscuro, fg=white},
    row{10} = {bg=AzulClaro, fg=white},
    cells  = {font=\fontsize{10pt}{12pt}\selectfont},
    }
        \textbf{Componente} & \textbf{Código}  & \textbf{CH Teór.} & \textbf{CH Prát.}  & \textbf{CH Rem.} & \textbf{CH Total} \\
        Arquitetura de Computação em Nuvem AWS & FEELT:CCAWS & 0 & 0 & 45 & 45 \\
    Computação Gráfica RV-RA & FEELT31721 & 30 & 15 & 0 & 45 \\
    Interação Humano-Computador & FACOM32704 & 60 & 0 & 0 & 60 \\
    Processamento Digital de Imagens & GBC216 & 60 & 0 & 0 & 60 \\
    Programação Web & FEELT:PWEB & 30 & 15 & 0 & 45 \\
    RV-RA Jogos Sérios & FEELT:JOGS & 0 & 0 & 0 & 45 \\
    Tecnologias Mobile & FEELT:TMOB & 0 & 0 & 0 & 45 \\
    Visão Computacional & FEELT:VCOM & 0 & 0 & 0 & 45 \\
    \textbf{TOTAL} &  & \textbf{180} & \textbf{30} & \textbf{45} & \textbf{390} \\
    \end{longtblr}
    \renewcommand{\familydefault}{\rmdefault} \normalfont{}
    
% \subsection{Certificado em Engenharia de Redes de \mbox{Computadores}}

Um certificado em Engenharia de Redes de Computadores oferece uma especialização que abrange os fundamentos, princípios avançados e práticas atualizadas relacionadas ao projeto, implementação e gerenciamento de redes de computadores. Os alunos que buscam esse certificado exploram disciplinas específicas que incluem, entre outras, protocolos de rede, arquitetura de redes, segurança de redes, administração de sistemas, tecnologias de comunicação e serviços de rede. Os estudantes são capacitados para enfrentar os desafios complexos do mundo das redes de computadores, capacitando-os para carreiras que exigem conhecimentos especializados em infraestrutura de rede, segurança cibernética, comunicação de dados e administração de sistemas.

% \renewcommand{\familydefault}{\sfdefault} \normalfont
% \begin{longtblr}[
%     theme = ecp,
%     caption = {Redes de Computadores},
%     label = {tab:redes_computadores},
% ]{
%   colspec = {Q[l,m,wd=65mm]Q[c,m,wd=24mm]Q[c,m,wd=8mm]Q[c,m,wd=8mm]Q[c,m,wd=8mm]Q[c,m,wd=10mm]},
%   rowhead = 1,
%   row{odd} = {bg=CinzaClaro},
%   row{1} = {bg=AzulEscuro, fg=white},
%   row{2} = {bg=AzulClaro, fg=white},
%   row{4} = {bg=AzulClaro, fg=white},
%   row{11} = {bg=AzulClaro, fg=white},
%   cells  = {font=\fontsize{10pt}{12pt}\selectfont},
% } 
%     % Componente & Código & CH Teór. & CH Prát.  & CH Rem. & CH Total \\
%     \textbf{Componente} & \textbf{Código}  & \textbf{CH Teór.} & \textbf{CH Prát.}  & \textbf{CH Rem.} & \textbf{CH Total} \\
%     \ \ \ \textit{Optativas A} &&&&&\\
%     Tecnologias Web & FEELT? & 0 & 45 & 0 & 45\\
%     \ \ \ \textit{Optativas B} &&&&&\\
%     Segurança de Sistemas Computacionais & FEELT? &  &  & 0 & 45 \\
%     Arquitetura de Computação em Nuvem AWS & FEELT? & 0 & 0 & 45 & 45 \\
%     Redes de Comunicações II & FEELT31831 & 45 & 15 & 0 & 60 \\
%     Projeto de Protocolos de Comunicação & GSI061 & 60 & 0 & 0 & 60 \\
%     Projeto de Redes de Computadores & GSI062 & 60 & 0 & 0 & 60 \\
%     Infraestrutura e Planejamento para Telecomunicações & FEELT39055 & 45 & 15 & 0 & 60 \\
%     \textbf{TOTAL} &  & \textbf{225} & \textbf{105} & \textbf{0} & \textbf{375} \\
% \end{longtblr}
% \renewcommand{\familydefault}{\rmdefault} \normalfont

\renewcommand{\familydefault}{\sfdefault} \normalfont{}
    \begin{longtblr}[
        theme = ecp,
        caption = {Engenharia de Redes de Computadores},
        label = {tab:engenharia_redes_computadores},
    ]{
    colspec = {Q[l,m,wd=65mm]Q[c,m,wd=24mm]Q[c,m,wd=8mm]Q[c,m,wd=8mm]Q[c,m,wd=8mm]Q[c,m,wd=10mm]},
    rowhead = 1,
    row{odd} = {bg=CinzaClaro},
    row{1} = {bg=AzulEscuro, fg=white},
    row{9} = {bg=AzulClaro, fg=white},
    cells  = {font=\fontsize{10pt}{12pt}\selectfont},
    }
        \textbf{Componente} & \textbf{Código}  & \textbf{CH Teór.} & \textbf{CH Prát.}  & \textbf{CH Rem.} & \textbf{CH Total} \\
        Arquitetura de Computação em Nuvem AWS & FEELT:CCAWS & 0 & 0 & 45 & 45 \\
    Infraestrutura e Planejamento para Telecomunicações & FEELT39055 & 45 & 15 & 0 & 60 \\
    Programação Web & FEELT:PWEB & 30 & 15 & 0 & 45 \\
    Projeto de Protocolos de Comunicação & GSI061 & 60 & 0 & 0 & 60 \\
    Projeto de Redes de Computadores & GSI062 & 60 & 0 & 0 & 60 \\
    Redes de Comunicações II & FEELT31831 & 45 & 15 & 0 & 60 \\
    Segurança de Sistemas Computacionais & FEELT:SEG & 0 & 0 & 0 & 45 \\
    \textbf{TOTAL} &  & \textbf{240} & \textbf{45} & \textbf{45} & \textbf{375} \\
    \end{longtblr}
    \renewcommand{\familydefault}{\rmdefault} \normalfont{}
    








\section{Atividades Curriculares de Extensão / ACE}

% A extensão universitária consiste em ações da universidade junto à comunidade na qual está inserida, disponibilizando, ao público externo, o conhecimento adquirido com o ensino e a pesquisa desenvolvidos dentro da própria universidade. O conceito de extensão está associado ao entendimento de que o conhecimento gerado pelas instituições de ensino e pesquisa deve ser indissociável da intenção de transformar uma realidade da sociedade através da tentativa de intervir para suprir as  deficiências identificadas, sem simplesmente se limitar à formação dos alunos regulares daquela instituição.

% Em atendimento a:

% \begin{itemize}
%     \item Lei nº 13.005 de 25 de junho de 2014, que aprova o Plano Nacional de Educação~-- PNE e dá outras providências;
%     \item Resolução CNE/CES nº 07/2018 de 18 de dezembro de 2018, que estabelece as Diretrizes para a Extensão na Educação Superior Brasileira e regimenta o disposto na Meta 12.7 da Lei nº 13.005/2014, que aprova o Plano Nacional de Educação (PNE 2014-2024) e dá outras providências;
%     \item Resolução nº 25/2019, do Conselho Universitário, que estabelece a Política de Extensão da Universidade Federal de Uberlândia, e dá outras providências;
%     \item Resolução nº 05/2020, do Conselho de Extensão, Cultura e Assuntos Estudantis, que dispõe sobre a elaboração do Plano de Extensão da Unidade (PEX) nas Unidades Acadêmicas e Unidades Especiais de Ensino, e dá outras providências;
%     \item Resolução nº 13/2019, do Conselho de Graduação, que regulamenta a inserção das atividades de extensão nos Currículos dos Cursos de Graduação da Universidade Federal de Uberlândia, na forma de componente curricular específico,
% \end{itemize}

% \noindent este projeto opta por destacar o controle e registro das atividades de extensão universitária de cada discente junto à comunidade. Para se adequar à legislação vigente, são criadas as seguintes disciplinas:

% \begin{itemize}
%     \item Atividades Curriculares de Extensão I, 75 horas práticas, sem pré-requisito;
%     \item Atividades Curriculares de Extensão II, 75 horas práticas, sem pré-requisito;
%     \item Atividades Curriculares de Extensão III, 90 horas práticas, sem pré-requisito;
%     \item Atividades Curriculares de Extensão IV, 90 horas práticas, sem pré-requisito;
%     \item Memorial de Atividades Curriculares de Extensão, 30 horas teóricas, com pré-requisito das disciplinas Atividades Curriculares de Entensão I a IV.
% \end{itemize}

% Na disciplina Introdução à Engenharia de Computação, existe um momento onde as práticas e normativas específicas de extensão serão apresentadas, preparando os estudantes para suas jornadas acadêmicas extensionistas, com apresentação dos docentes do curso e de possíveis projetos extensionistas.

% As disciplinas Atividades Curriculares de Extensão I a IV, sem pré-requisitos, são flexíveis e possuem atividades programadas pelo próprio discente em conjunto com o coordenador do projeto de extensão ao qual esteja vinculado. Portanto, não dependem de encontros semanais nem do calendário acadêmico oficial. A quitação de tais disciplinas fica condicionada à entrega dos certificados devidos, com indicativo de horas, obtidos pelos discentes através do Sistema de Informação de Extensão e Cultura conhecido como a plataforma SIEX da PROEXC/UFU. O registro de tais projetos de extensão na plataforma SIEX cabem ao docente coordenador dos respectivos projeto.

% Já a disciplina Memorial de Atividades Curriculares de Extensão, com pré-requisito das disciplinas Atividades Curriculares de Extensão I a IV, é uma disciplina teórica e pretende reunir os estudantes para que os mesmos construam um memorial, discutam e apresentem suas experiências na extensão universitária.

% O detalhamento do processo de convalidação das Atividades Curriculares de Extensão consta em normas específicas aprovadas nos âmbitos do Colegiado do Curso com anuência do NDE. Os casos omissos deverão ser deliberados pelo Colegiado do Curso.

A extensão universitária se destina a conectar a universidade à comunidade em que está inserida, compartilhando com o público externo o conhecimento obtido por meio do ensino e da pesquisa realizados internamente. Essa abordagem busca atuar na transformação da realidade social, intervindo para suprir deficiências identificadas, indo além da formação dos alunos regulares da instituição. Sobre extensão universitária, este projeto pedagógico atende às legislações e diretrizes em vigor~\cite{Lei:13005:2014,MEC:CNE:CES:7:2018,UFU:CONSUN:25:2019,UFU:CONGRAD:13:2019,UFU:CONSEX:5:2020}.

% \begin{itemize}
%     \item Lei nº 13.005 de 25 de junho de 2014, que aprova o Plano Nacional de Educação -- PNE, em especial sua Meta 12.7;
%     \item Resolução CNE/CES nº 07/2018 de 18 de dezembro de 2018, estabelecendo as Diretrizes para a Extensão na Educação Superior Brasileira;
%     \item Resolução nº 25/2019 do Conselho Universitário, definindo a Política de Extensão da Universidade Federal de Uberlândia;
%     \item Resolução nº 05/2020 do Conselho de Extensão, Cultura e Assuntos Estudantis, relacionada ao Plano de Extensão da Unidade (PEX);
%     \item Resolução nº 13/2019 do Conselho de Graduação, regulamentando a inserção de atividades de extensão nos Currículos dos Cursos de Graduação da Universidade Federal de Uberlândia.
% \end{itemize}

No presente projeto, as atividades de extensão universitária junto à comunidade foram destacadas das atividades complementares, para possibilitar o controle e o registro independente das ações extensionistas de cada discente visando ao atendimento da meta 12.7 da Lei nº 13.005/2014\footnote{Meta de 10\% dos créditos curriculares destinados à extensão universitária}~\cite{Lei:13005:2014}.

Este projeto propõe o controle e registro das atividades de extensão universitária de cada discente junto à comunidade, em conformidade com a legislação vigente, por meio da criação das disciplinas:

\begin{itemize}
    \item Atividades Curriculares de Extensão I, II, III e IV, totalizando horas práticas (90h, 90h, 75h, 60h respectivamente)\footnote{Reservas de horário sugeridos para estudantes participarem em projetos de extensão.} sem pré-requisitos;
    \item Memorial de Atividades Curriculares de Extensão, exigindo pré-requisito das disciplinas anteriores para 30 horas teóricas.
\end{itemize}

Durante a disciplina ``Introdução à Engenharia de Computação'', serão apresentadas as práticas e normativas específicas relacionadas à extensão universitária, com o objetivo de preparar os(as) estudantes para jornadas acadêmicas extensionistas e explorar possibilidades de envolvimento em projetos.

Os componentes curriculares de ``Atividades Curriculares de Extensão I a IV'' (ACEs), flexíveis e sem pré-requisitos, envolvem atividades programadas pelo discente e coordenadação do projeto de extensão ao qual seja vinculado, devidamente registrados na plataforma SIEX\footnote{Sistema de Informação de Extensão da UFU - \url{http://www.siex.proex.ufu.br/}} da PROEXC/UFU, e que não dependem do calendário letivo em vigor. São exemplos de projetos extensionistas reconhecidos nos termos destas atividades:

\begin{itemize}
\item Projetos desenvolvidos em parceria ativa com a comunidade externa por docentes e/ou técnico-administrativos do curso ou da FEELT, devidamente registrados na plataforma SIEX;
\item Projetos desenvolvidos em parceria ativa com a comunidade externa por docentes e/ou técnico-administrativos de outras unidades acadêmicas da UFU, também registrados no SIEX;
\item Projetos desenvolvidos em parceria ativa pelo(a) próprio(a) estudante com a comunidade externa, representada por empresas (inclusive o local de estágio externo), instituições, organizações públicas ou da sociedade civil, ou outras iniciativas com relevância social, desde que orientados por docentes ou técnico-administrativos da UFU e registrados no SIEX;
\item Participação em projetos ou na administração da CONSELT\footnote{\url{https://conselt.com.br/}}, empresa júnior vinculada à FEELT, ou em outras empresas juniores da UFU, em razão de sua natureza extensionista e do vínculo institucional estabelecido por meio de docentes ou técnico-administrativos da UFU que atuam como tutores;
\item Contribuições efetivas — compreendidas como participações registradas, validadas por orientador(a) da UFU e preferencialmente aceitas nos repositórios principais — em projetos FOSS\footnote{\textit{Free and Open Source Software}, software de código aberto e livre, conforme os princípios da \textit{Open Source Initiative} disponíveis em \url{https://opensource.org/osd}, com tutoriais de colaboração e sites para encontrar projetos em \url{https://www.firsttimersonly.com/}}, OSHW\footnote{\textit{Open Source Hardware}, hardware livre e aberto, conforme os princípios da \textit{Open Source Hardware Association}, disponíveis em \url{https://oshwa.org/resources/open-source-hardware-definition/}}, OSAI\footnote{\textit{Open Source Artificial Intelligence}, inteligência artificial aberta, conforme os princípios disponíveis em \url{https://opensource.org/ai/open-source-ai-definition}} ou outras iniciativas assim reconhecidas, com documentação técnica e códigos-fonte publicados em repositórios públicos amplamente reconhecidos, desde que a participação seja orientada por docentes ou técnico-administrativos da UFU e registrada no SIEX;
\item Ações de comunicação pública do conhecimento, tais como podcasts, blogs, vlogs, vídeos ou outras mídias digitais, desde que fundamentadas na interação dialógica entre a universidade e a comunidade externa e resultem em produção conjunta ou retroalimentada de saberes, sob orientação de docentes ou técnico-administrativos da UFU e com registro no SIEX;
\item Casos omissos deverão ser avaliados e aprovados pelo Colegiado do Curso com anuência do Núcleo Docente Estruturante (NDE) e da Coordenação de Extensão da Unidade Acadêmica.
\end{itemize}

A disciplina ``Memorial de Atividades Curriculares de Extensão'', cujo pré-requisito é o cumprimento prévio de atividades extensionistas, tem como objetivo proporcionar um espaço coletivo de reflexão e elaboração individual de um memorial. Durante a disciplina, os(as) estudantes participarão de discussões sobre suas experiências extensionistas e produzirão um registro crítico e pessoal dessas vivências.

A aprovação na disciplina está condicionada à entrega do memorial e à comprovação, por meio de certificados emitidos via plataforma SIEX, da carga horária mínima exigida de atividades de extensão (\textbf{345 horas}
\unskip). O registro dos projetos de extensão no SIEX, necessário para a emissão dos certificados, é de responsabilidade dos respectivos coordenadores vinculados à UFU. Os(as) estudantes têm liberdade para participar de quantos projetos desejarem, desde que cumpram integralmente a carga horária mínima prevista pelas diretrizes curriculares.


\section{Atividades de Conclusão de Curso / ACC}
\label{sec:acc_conclusao}

Baseado na legislação vigente~\cite{MEC:CNE:CES:5:2016}, o presente projeto pedagógico permite ao estudante validar o componente curricular ``Atividade de Conclusão de Curso'' com a opção entre ``Estágio Supervisionado'' e ``Trabalho de Conclusão de Curso''.

Ambas as opções possuem carga horária mínima a ser integralizada igual a \textbf{300 horas}
\unskip, com pré-requisito de \textbf{1875 horas}
de integralização do curso, o equivalente a estar começando o \input{include/auto/par_equivale_min_acc.tex}\unskipº Período.

Ainda, o estudante tem a liberdade de defender um ``Trabalho de Conclusão de Curso'' nas modalidades ``projeto de pesquisa acadêmica'' ou ``empresa do tipo startup''.

\subsection{Estágio Supervisionado}

Regulamentado como ``Estágio Obrigatório'' pelas Normas Gerais de Estágio de Graduação da Universidade Federal de Uberlândia~\cite{UFU:CONGRAD:93:2023}, o Estágio Supervisionado é uma das atividades possíveis necessárias para a conclusão do curso de \nomecursonovo quando da sua validação como Atividade de Conclusão de Curso. O estudante necessita, obrigatoriamente, cumprir uma carga horária mínima estipulada de \textbf{300 horas}
de estágio na sua área de formação. São necessários o acompanhamento de um supervisor -- um profissional da mesma área de formação (ou área afim) que faça parte do quadro de empregados da parte cedente do estágio -- e a realização de horas supervisionadas por um professor do curso. Ao final do estágio, o estudante deve apresentar um relatório para o registro final das atividades realizadas.

Para realizar essa atividade, o discente deve optar pela modalidade de ``estágio supervisionado'' como ``Atividade de Conclusão do Curso'' e ter, como pré-requisito, um mínimo de \textbf{1875 horas}
integralizadas do currículo. Quando da conclusão do Estágio Supervisionado, a carga horária integralizada com essa atividade é de no máximo \textbf{300 horas}
\unskip.

Um certificado de conclusão de estágio deverá ser emitido pela Coordenação de Estágio do Curso. O detalhamento dos procedimentos relativos ao Estágio Supervisionado consta nas Normas Complementares de Estágio do Curso de Graduação em \nomecursonovo\footnote{\url{http://www.feelt.ufu.br/graduacao/engenharia-de-computacao/saiba-mais/estagio-supervisionado}}, aprovada nos âmbitos do Colegiado do Curso e da Unidade Acadêmica com anuência do NDE.

\subsubsection{Estágio Não Obrigatório}

Embora não seja uma modalidade aceita para a quitação da Atividade de Conclusão de Curso, de acordo com as Normas Gerais de Estágio de Graduação da Universidade Federal de Uberlândia~\cite{UFU:CONGRAD:93:2023}, o estudante é autorizado a desenvolver um Estágio Não Obrigatório. Essa modalidade de estágio é desenvolvida como atividade opcional e complementar, acrescida à carga horária regular e obrigatória de acordo com as normas complementares de estágio e as normas de atividades complementares do curso. São necessários o acompanhamento de um supervisor -- um profissional da mesma área de formação (ou área afim) que faça parte do quadro de empregados da parte cedente do estágio -- e a realização de horas supervisionadas por um professor do curso.

Para realizar essa atividade, o discente tem o pré-requisito mínimo de \textbf{720 horas}
integralizadas do currículo com a necessidade de aprovação em todos os componentes curriculares dos 1º e 2º
períodos do curso. Quando da conclusão do Estágio Não Obrigatório, a carga horária integralizada com essa atividade a partir das devidas comprovações está definida em Atividades Acadêmicas Complementares (ver \autoref{sec:AAC}).

O detalhamento do Estágio Não Obrigatório consta nas Normas Complementares de Estágio do Curso de Graduação em \nomecursonovo\footnote{\url{http://www.feelt.ufu.br/graduacao/engenharia-de-computacao/saiba-mais/estagio-supervisionado}}, aprovada nos âmbitos do Colegiado do Curso e da Unidade Acadêmica com anuência do NDE.

\subsection{Trabalho de Conclusão de Curso}

O estudante do curso de \nomecursonovo poderá desenvolver um Trabalho de Conclusão de Curso (TCC), caso opte por essa modalidade como forma de cumprimento da Atividade de Conclusão de Curso. O TCC poderá assumir diferentes formatos, tais como um projeto de pesquisa acadêmica ou a especificação e desenvolvimento de uma proposta de empresa startup, desde que orientado por docente da UFU e avaliado por banca designada.

Para realizar essa atividade, o discente deve optar pelo TCC como Atividade de Conclusão do Curso e ter, como pré-requisito mínimo, \textbf{1875 horas}
 integralizadas do currículo. Quando da conclusão do TCC em qualquer das modalidades, a carga horária integralizada com essa atividade é de no máximo \textbf{300 horas}
em Atividade de Conclusão de Curso.

O detalhamento dos procedimentos relativos ao TCC consta consta nas Normas Gerais para os Trabalhos de Conclusão de Curso (TCC) do Curso de Graduação em \nomecursonovo\footnote{\url{https://www.feelt.ufu.br/graduacao/engenharia-de-computacao/saiba-mais/tcc}}, aprovada nos âmbitos do Colegiado do Curso e da Unidade Acadêmica com anuência do NDE.

\subsubsection{TCC: Projeto de Pesquisa Acadêmica}

No caso da modalidade de projeto de pesquisa, o desenvolvimento de um Trabalho de Conclusão de Curso (TCC) consolida o processo formativo do estudante, integrando conhecimentos adquiridos ao longo do curso. Essa experiência estimula a prática investigativa, a capacidade de estabelecer conexões entre diferentes correntes teóricas, o aprimoramento da pesquisa bibliográfica e a análise crítica de dados obtidos por múltiplas fontes. Busca-se, com isso, fomentar a criatividade, a autonomia intelectual e a produção científica fundamentada. Nesta modalidade, a defesa do TCC ocorre por meio da apresentação formal de uma monografia científica que, embora possa dialogar com outras áreas do conhecimento, deve ter como eixo central a aplicação dos saberes desenvolvidos ao longo da formação em \nomecursonovo.

\subsubsection{TCC: Empresa Tipo Startup}

Uma empresa do tipo startup é, em geral, caracterizada por um grupo de pessoas — preferencialmente interdisciplinar — que atua em condições de incerteza e é impulsionado por ideias inovadoras, buscando modelos de negócio sustentáveis que gerem valor de forma escalável. Na modalidade de TCC voltada à criação de uma startup, objetiva-se estimular a criatividade do estudante por meio do enfrentamento de desafios reais e do desenvolvimento de soluções inovadoras que dialoguem com demandas concretas da sociedade. Busca-se, ainda, fomentar o pensamento empreendedor e a capacidade de estruturar produtos e serviços com potencial de impacto. Nesta modalidade, a defesa do TCC ocorre por meio da apresentação formal de um \textit{pitch}, acompanhada da construção de um Produto Mínimo Viável (PMV) e da elaboração de um estudo de viabilidade utilizando ferramentas como o \textit{Business Model Canvas}.

\section{Atividades Acadêmicas Complementares / AAC}
\label{sec:AAC}

De acordo com o Parecer do CNE/CES nº 136/2012~\cite{MEC:CNE:CES:136:2012}, as atividades complementares são componentes curriculares que têm como objetivo principal enriquecer e expandir o perfil do egresso com atividades que privilegiem aspectos diversos da sua formação, incluindo atividades desenvolvidas fora do ambiente acadêmico. Tais atividades constituem instrumental importante para o desenvolvimento pleno do aluno, servindo de estímulo a uma formação prática independente e interdisciplinar, sobretudo nas relações com o mundo do trabalho.

Ainda, as atividades podem ser cumpridas em diversos ambientes, como a instituição a que o estudante está vinculado, outras instituições e variados ambientes sociais, técnico-científicos ou profissionais, em modalidades tais como: formação profissional (cursos de formação profissional, experiências de trabalho ou estágios não obrigatórios), de pesquisa e publicação (iniciação científica e participação em eventos técnico-científicos, publicações científicas), de ensino (programas de monitoria e tutoria), de gestão e política (representação discente em comissões e comitês), de empreendedorismo e inovação (participação em incubadoras, \textit{startups} ou outros mecanismos), de qualificação e experiência internacional, entre outras. Essas atividades devem ser permanentemente incentivadas no cotidiano acadêmico, permitindo a diversificação das atividades complementares desenvolvidas pelos estudantes.

No Quadro~\ref{tab:aac}, temos a relação de Atividades Acadêmicas Complementares (AAC) propostas neste projeto, com carga horária máxima a ser integralizada por tipo de atividade. A carga horária mínima a ser integralizada com essas atividades é de \textbf{90 horas}, sem necessidade de pré-requisito.
São consideradas atividades válidas para a quitação desse componente as que \textbf{foram realizadas durante o período de vínculo do estudante ao curso}.

% \renewcommand{\familydefault}{\sfdefault} \normalfont
\begin{longtblr}[
    theme = ecp,
    caption = {Carga Horária Máxima das AAC},
    label = {tab:aac},
]{
  colspec = {Q[c,m,wd=17mm]Q[l,m,wd=105mm]Q[c,m,wd=14mm]},
  rowhead = 1,
  rowsep = 1pt,
  row{odd} = {bg=CinzaClaro},
  row{1} = {bg=AzulEscuro, fg=white},
  row{2} = {bg=AzulClaro, fg=white},
  row{8} = {bg=AzulClaro, fg=white},
  row{18} = {bg=AzulClaro, fg=white},
  row{26} = {bg=AzulClaro, fg=white},
  row{31} = {bg=AzulClaro, fg=white},
  row{35} = {bg=AzulClaro, fg=white},
  row{39} = {bg=AzulClaro, fg=white},
  cells  = {font = \fontsize{10pt}{12pt}\selectfont},
} 
    \textbf{Código} & \textbf{Descrição} & \textbf{CH Máx.} \\
    & \textit{Formação Profissional} \\
   ATCO1181 & Realização e conclusão de Curso Online Aberto e Massivo (MOOC) aprovado pelo Colegiado do Curso & 90 \\
   ATCO1135 & Participação em oficinas, cursos ou mini-cursos relacionados ao aprendizado de técnicas úteis à profissão & 45 \\
   ATCO1182 & Obtenção de certificações técnicas na área de Computação & 60 \\
   ATCO0725 & Participação em visitas técnicas orientadas & 10 \\
   ATCO0254 & Estágio não obrigatório & 36 \\
    & \textit{Pesquisa} \\
   ATCO1104 & Participação em Iniciação Científica com bolsa (PIBIC ou PIBITI) & 90 \\
   ATCO1105 & Participação em Iniciação Científica sem bolsa (PIVIC) & 90\\
   ATCO0044 & Apresentação de trabalhos em eventos científicos na forma oral ou pôster & 45 \\
   ATCO0964 & Publicação de trabalhos científicos - resumo e/ou pôster  & 20 \\
   ATCO0965 & Publicação de Trabalhos completos em anais de eventos & 45 \\
   ATCO0993 & Publicações em periódicos especializados (revistas indexadas da área) & 90\\
   ATCO0994 & Publicações em periódicos não especializados (revistas de outras áreas, jornais e revistas não indexadas) & 15\\
   ATCO1184 & Publicação de livro ou capítulo de livro especializado com código ISBN e corpo editorial técnico-científico & 90\\
   ATCO1185 & Publicação de livro ou capítulo de livro, especializado ou não, sem código ISBN ou corpo editorial técnico-científico & 15\\
    & \textit{Ensino} \\
    ATCO0354 & Monitoria em disciplina ministrada na UFU, fora do escopo do curso & 15\\
    ATCO1919 & Monitoria em disciplina de graduação relativas aos 1º e 2º semestres do curso & 30\\
    ATCO2020 & Monitoria em disciplina de graduação relativas ao 3º semestre em diante do curso & 45\\
    ATCO0753 & Participação no Programa de Educação Tutorial -- PET & 60\\
    ATCO0599 & Participação em grupo de estudos de temas específicos registrado e certificado pela Instituição & 30\\
    ATCO1186 & Participação orientada por docente no desenvolvimento de material informacional ou didático para uso interno à UFU & 30 \\
    ATCO1187 & Ministrante de palestras, minicursos, seminários e oficinas para comunidade interna da UFU & 30\\
    & \textit{Gestão e Representação Estudantil} \\
    ATCO0319 & Membro de Diretório Acadêmico & 30\\
    ATCO0327 & Membro do Diretório Central dos Estudantes & 30\\
    ATCO1019 & Representante Discente no Conselho de Unidade ou Colegiado de Curso & 45\\
    ATCO0315 & Membro de Conselho Superior da UFU & 45\\
    & \textit{Empreendorismo e Inovação} \\
    ATCO1188 & Participação ou desenvolvimento de projetos junto a incubadoras de empresas & 45\\
    ATCO1189 & Fundador ou membro de empresa do tipo \textit{startup} de tecnologia & 90\\
    ATCO2121 & Patente ou Registro de Software & 90\\
    & \textit{Qualificação e Experiência Internacional}  \\
    ATCO1099 & Curso de língua estrangeira ou aprovação em exame de proficiência em língua estrangeira & 30\\
    ATCO0344 & Mobilidade Internacional oficializada pela DRII/UFU & 60\\
    ATCO1190 & Realização de intercâmbio internacional para estágio ou pesquisa na área de formação & 60\\
    & \textit{Outras}  \\
    ATCO0750 & Participação no Exame Nacional do Desempenho de Estudante (ENADE) & 15\\
    ATCO0345 & Mobilidade Nacional & 45\\
    ATCO0285 & Frequência e aprovação em disciplinas facultativas (outras Unidades Acadêmicas da UFU ou outra IES) & 45\\
    ATCO0706 & Participação em projetos institucionais (PIBEG, PGB, PIBID, PROGRAD), sem registro no SIEX & 30\\
    ATCO0492 & Participação em Competições Técnicas & 90\\
    ATCO0491 & Participação em Competições Culturais, Artísticas, ou Esportivas & 35\\
    ATCO1191 & Participação como ouvinte em eventos técnicos e/ou científicos (congressos, simpósios, seminários, mesa-redonda, workshops) & 30\\
    ATCO1192 & Organização ou participação na organização de eventos institucionais, técnicos ou científicos para comunidade interna da UFU & 30\\
\end{longtblr}
% \renewcommand{\familydefault}{\rmdefault} \normalfont

A requisição para a quitação desse componente curricular é de responsabilidade do estudante que deverá apresentar requerimento com esse objetivo. O detalhamento das Atividades Acadêmicas Complementares  consta em normas específicas\footnote{\url{https://www.feelt.ufu.br/graduacao/engenharia-de-computacao/saiba-mais/atividades-complementares}} aprovadas nos âmbitos do Colegiado do Curso com anuência do NDE e da Unidade Acadêmica. Os casos omissos deverão ser tratados pelo Colegiado do Curso.



